% HYDROGEOLOGY — Geology Upper Sixth — Cours signé OnBuch+
% Official MINESEC syllabus. Compiler avec : tectonic GEO07-hydrogeology.tex
\newcommand{\DOCMATIERE}{Geology}
\newcommand{\DOCNIVEAU}{Upper Sixth}
\newcommand{\DOCMODULE}{Upper Sixth — Geology}
\newcommand{\DOCLECON}{7}
\newcommand{\DOCTITRE}{HYDROGEOLOGY}
% OnBuch+ — préambule « cours signés » ENRICHI (compilé avec tectonic / XeLaTeX)
% Système visuel « Pop » d'OnBuch+ : fond crème (jamais blanc), bordures encre
% épaisses, ombres « dures » décalées, titres Archivo Black en capitales.
% Version MATHÉMATIQUES : graphes (pgfplots/tikz), boîtes pédagogiques
% (définition, propriété, méthode, attention, le savais-tu, exercice résolu,
% exercice, TP, bilan), page de garde et signature OnBuch+ ancrée en bas.
%
% Macros attendues dans main.tex AVANT \input{preamble} :
%   \DOCMATIERE  \DOCNIVEAU  \DOCMODULE  \DOCLECON (numéro)  \DOCTITRE
\documentclass[11pt]{article}

\usepackage{fontspec}
\usepackage[english]{babel}
\usepackage[a4paper,margin=2cm,top=3.4cm,bottom=2.8cm,headheight=1.4cm,footskip=1.3cm]{geometry}
\usepackage{amsmath,amssymb,amsfonts}
\usepackage{mathtools} % \xmapsto, \coloneqq... souvent utilisés par les modèles
\usepackage{cancel} % simplifications barrées (bilans de forces)
% \abs{z} (module, valeur absolue) et \norm{u} (norme), utilisés par les modèles
\providecommand{\abs}{}\let\abs\relax\DeclarePairedDelimiter{\abs}{\lvert}{\rvert}
\providecommand{\norm}{}\let\norm\relax\DeclarePairedDelimiter{\norm}{\lVert}{\rVert}
\usepackage[table]{xcolor} % [table] : fournit aussi \rowcolors (alternance de lignes)
\usepackage{pagecolor}
\usepackage{tikz}
\usetikzlibrary{shapes.symbols,circuits.logic.US,circuits.logic.IEC,arrows.meta,positioning,calc,shapes.geometric,shapes.misc,shapes.arrows,decorations.pathmorphing,decorations.pathreplacing,decorations.markings,patterns,shadows,mindmap,trees,fit,backgrounds,matrix,intersections,angles,quotes}
\usepackage{pgfplots}
\usepackage[version=4]{mhchem} % équations nucléaires \ce{^{235}_{92}U}
\usepackage[european,siunitx]{circuitikz} % schémas électriques
\pgfplotsset{compat=1.18}
\usepgfplotslibrary{fillbetween,groupplots} % aires entre courbes (calcul intégral)
\usepackage{enumitem}
\usepackage{fancyhdr}
% [most] charge toutes les bibliothèques tcolorbox (listings, minted, raster,
% documentation...) et gonfle inutilement la mémoire TeX sur un document long
% (~300 boîtes) : on ne charge que ce qui est réellement utilisé.
\usepackage[skins,breakable]{tcolorbox}
\usepackage{graphicx}
\usepackage{colortbl}
\usepackage{booktabs}
\usepackage{tabularx}
\usepackage{diagbox} % case d'en-tête diagonale des tableaux à double entrée
% Colonnes extensibles centrée / gauche / droite (C, L, R), souvent utilisées par les modèles
\newcolumntype{C}{>{\centering\arraybackslash}X}
\newcolumntype{L}{>{\raggedright\arraybackslash}X}
\newcolumntype{R}{>{\raggedleft\arraybackslash}X}
\usepackage{array}
\usepackage{multicol}
\usepackage{siunitx}
% parse-numbers=false : accepte aussi \SI{2,5 \times 10^{-3}}{...}
\sisetup{locale=UK,output-decimal-marker={.},group-minimum-digits=4,parse-numbers=false}
\DeclareSIUnit\ohm{\ensuremath{\symup{\Omega}}} % U+2126 absent de la police
\DeclareSIUnit\division{div} % réglages d'oscilloscope (ms/div, V/div)
\DeclareSIUnit\year{yr}\DeclareSIUnit\annum{yr}\DeclareSIUnit\Ma{Ma}\DeclareSIUnit\tr{tr} % tours (vitesse de rotation, ex. \SI{1500}{\tr\per\minute})
\usepackage{qrcode}
% \degree : commande fréquemment utilisée par les modèles pour un angle ou une
% température en dehors de \SI{}{} (non fournie par les paquets chargés).
\providecommand{\degree}{^{\circ}}
% \qed (fin de démonstration), utilisé par les modèles sans amsthm
\providecommand{\qed}{\hfill$\square$}
% \barre : double barre des valeurs interdites dans un tableau de variations (tkz-tab)
\providecommand{\barre}{\ensuremath{\|}}
% environnement proof (amsthm non chargé) utilisé par les modèles
\ifdefined\proof\else\newenvironment{proof}[1][Proof]{\par\noindent\textit{#1.}\ }{\qed\par}\fi
\usepackage{lastpage}
\usepackage{hyperref}
\hypersetup{hidelinks}

% ============================================================
% Palette « Pop » OnBuch+ — mêmes hex que la classe P (plus_ui.dart)
% ============================================================
\definecolor{poporange}{HTML}{F59321}
\definecolor{popdark}{HTML}{E07A0C}
\definecolor{popcream}{HTML}{FFF6E8}
\definecolor{popink}{HTML}{1C1714}
\definecolor{poppurple}{HTML}{7A5AE0}
\definecolor{popgreen}{HTML}{2E9E6A}
\definecolor{poppink}{HTML}{E0457B}
\definecolor{popblue}{HTML}{2F7DE1}
\definecolor{popgold}{HTML}{C98A12}
\definecolor{popmuted}{HTML}{8A7F76}
\definecolor{popline}{HTML}{EADFD2}
\definecolor{popgray}{HTML}{8A7F76}
\colorlet{popgrey}{popgray}
% Alias tolérants (le modèle nomme parfois les couleurs différemment)
\colorlet{popwhite}{white}
\colorlet{popblack}{popink}
\colorlet{popred}{poppink}
\colorlet{popyellow}{popgold}
% Teintes claires pour fonds de tableaux / zones
\colorlet{popblueL}{popblue!10!white}
\colorlet{poporangeL}{poporange!14!white}
\colorlet{popgreenL}{popgreen!12!white}
\colorlet{poppurpleL}{poppurple!12!white}
\colorlet{poppinkL}{poppink!12!white}
\colorlet{popgoldL}{popgold!14!white}

\pagecolor{popcream}

% ============================================================
% Typographie
% ============================================================
\defaultfontfeatures{Ligatures=TeX}
\setmainfont{PlusJakartaSans-Medium.ttf}[Path=../../../fonts/,BoldFont=PlusJakartaSans-Bold.ttf]
% Maths en sans-serif (Fira Math) avec chiffres et lettres droites de Plus
% Jakarta, pour que les formules se fondent dans le texte.
\usepackage[math-style=ISO,bold-style=ISO]{unicode-math}
\setmathfont{FiraMath-Regular.otf}[Path=../../../fonts/]
\setmathfont{PlusJakartaSans-Medium.ttf}[Path=../../../fonts/,range={up/num,up/{Latin,latin},\mathup}]
% (icomma retiré : décimales avec point en anglais)
% Caractères mathématiques Unicode absents des polices de texte (Plus Jakarta,
% Archivo Black) : rendus par leur équivalent mathématique, en texte comme en
% formule — sinon ils s'affichent en carré vide (ex. « Arithmétique dans ℤ »).
\usepackage{newunicodechar}
\newunicodechar{ℝ}{\ensuremath{\mathbb{R}}}
\newunicodechar{ℤ}{\ensuremath{\mathbb{Z}}}
\newunicodechar{ℂ}{\ensuremath{\mathbb{C}}}
\newunicodechar{ℕ}{\ensuremath{\mathbb{N}}}
\newunicodechar{ℚ}{\ensuremath{\mathbb{Q}}}
\newunicodechar{𝔻}{\ensuremath{\mathbb{D}}}
\newunicodechar{α}{\ensuremath{\alpha}}
\newunicodechar{β}{\ensuremath{\beta}}
\newunicodechar{γ}{\ensuremath{\gamma}}
\newunicodechar{δ}{\ensuremath{\delta}}
\newunicodechar{ε}{\ensuremath{\varepsilon}}
\newunicodechar{θ}{\ensuremath{\theta}}
\newunicodechar{λ}{\ensuremath{\lambda}}
\newunicodechar{μ}{\ensuremath{\mu}}
\newunicodechar{σ}{\ensuremath{\sigma}}
\newunicodechar{τ}{\ensuremath{\tau}}
\newunicodechar{φ}{\ensuremath{\varphi}}
\newunicodechar{ψ}{\ensuremath{\psi}}
\newunicodechar{ω}{\ensuremath{\omega}}
\newunicodechar{Σ}{\ensuremath{\Sigma}}
% majuscules produites par \MakeUppercase dans les titres (α-aminés -> Α)
\newunicodechar{Α}{\ensuremath{\alpha}}
\newunicodechar{Β}{\ensuremath{\beta}}
\newunicodechar{Γ}{\ensuremath{\Gamma}}
\newunicodechar{Δ}{\ensuremath{\Delta}}
\newunicodechar{Ε}{\ensuremath{\varepsilon}}
\newunicodechar{Θ}{\ensuremath{\Theta}}
\newunicodechar{Λ}{\ensuremath{\Lambda}}
\newunicodechar{Μ}{\ensuremath{\mu}}
\newunicodechar{Τ}{\ensuremath{\tau}}
\newunicodechar{Φ}{\ensuremath{\Phi}}
\newunicodechar{Ψ}{\ensuremath{\Psi}}
\newunicodechar{Ω}{\ensuremath{\Omega}}
\newunicodechar{↦}{\ensuremath{\mapsto}}
\newunicodechar{∈}{\ensuremath{\in}}
\newunicodechar{∉}{\ensuremath{\notin}}
\newunicodechar{∧}{\ensuremath{\wedge}}
\newunicodechar{⇔}{\ensuremath{\Leftrightarrow}}
\newunicodechar{⟺}{\ensuremath{\Longleftrightarrow}}
\newunicodechar{⇒}{\ensuremath{\Rightarrow}}
\newunicodechar{⟹}{\ensuremath{\Longrightarrow}}
\newunicodechar{∘}{\ensuremath{\circ}}
\newunicodechar{∪}{\ensuremath{\cup}}
\newunicodechar{∩}{\ensuremath{\cap}}
\newunicodechar{≡}{\ensuremath{\equiv}}
\newunicodechar{⊂}{\ensuremath{\subset}}
\newunicodechar{⊥}{\ensuremath{\perp}}
\newunicodechar{∃}{\ensuremath{\exists}}
\newunicodechar{∀}{\ensuremath{\forall}}
\newunicodechar{∛}{\ensuremath{\sqrt[3]{\,}}}
\newunicodechar{∜}{\ensuremath{\sqrt[4]{\,}}}
\newunicodechar{⃗}{\ensuremath{{}^{\to}}}
\newunicodechar{ⁿ}{\ifmmode^{n}\else\textsuperscript{\ensuremath{n}}\fi}
\newunicodechar{ˣ}{\ifmmode^{x}\else\textsuperscript{\ensuremath{x}}\fi}
\newunicodechar{ᵇ}{\ifmmode^{b}\else\textsuperscript{\ensuremath{b}}\fi}
\newunicodechar{ʳ}{\ifmmode^{r}\else\textsuperscript{\ensuremath{r}}\fi}
\newunicodechar{ᵘ}{\ifmmode^{u}\else\textsuperscript{\ensuremath{u}}\fi}
\newunicodechar{ʸ}{\ifmmode^{y}\else\textsuperscript{\ensuremath{y}}\fi}
\newunicodechar{ᵃ}{\ifmmode^{a}\else\textsuperscript{\ensuremath{a}}\fi}
\newunicodechar{ᵉ}{\ifmmode^{e}\else\textsuperscript{\ensuremath{e}}\fi}
\newunicodechar{ᵏ}{\ifmmode^{k}\else\textsuperscript{\ensuremath{k}}\fi}
\newunicodechar{ᵖ}{\ifmmode^{p}\else\textsuperscript{\ensuremath{p}}\fi}
\newunicodechar{ᴺ}{\ifmmode^{N}\else\textsuperscript{\ensuremath{N}}\fi}
\newunicodechar{ᶜ}{\ifmmode^{c}\else\textsuperscript{\ensuremath{c}}\fi}
\newunicodechar{ʰ}{\ifmmode^{h}\else\textsuperscript{\ensuremath{h}}\fi}
\newunicodechar{ᵠ}{\ifmmode^{q}\else\textsuperscript{\ensuremath{q}}\fi}
\newunicodechar{ₙ}{\ifmmode_{n}\else\textsubscript{\ensuremath{n}}\fi}
\newunicodechar{ₐ}{\ifmmode_{a}\else\textsubscript{\ensuremath{a}}\fi}
\newunicodechar{ᵢ}{\ifmmode_{i}\else\textsubscript{\ensuremath{i}}\fi}
\newunicodechar{ᵣ}{\ifmmode_{r}\else\textsubscript{\ensuremath{r}}\fi}
\newunicodechar{ₖ}{\ifmmode_{k}\else\textsubscript{\ensuremath{k}}\fi}
\newunicodechar{ᵦ}{\ifmmode_{\beta}\else\textsubscript{\ensuremath{\beta}}\fi}
\newunicodechar{ₓ}{\ifmmode_{x}\else\textsubscript{\ensuremath{x}}\fi}
\newfontfamily\popdisplay{ArchivoBlack-Regular.ttf}[Path=../../../fonts/]
\newfontfamily\poplogo{SpaceGrotesk-Bold.ttf}[Path=../../../fonts/]
\newfontfamily\popsemi{PlusJakartaSans-SemiBold.ttf}[Path=../../../fonts/]
\newfontfamily\popxbold{PlusJakartaSans-ExtraBold.ttf}[Path=../../../fonts/]
\newfontfamily\popxboldls{PlusJakartaSans-ExtraBold.ttf}[Path=../../../fonts/,LetterSpace=3.0]

\setlist[enumerate]{leftmargin=1.7em,itemsep=5pt,topsep=4pt}
\setlist[itemize]{leftmargin=1.7em,itemsep=4pt,topsep=4pt,label={\color{poporange}\textbullet}}
\setlength{\parindent}{0pt}
\setlength{\parskip}{5pt}
\renewcommand{\arraystretch}{1.25}

% pgfplots : style Pop par défaut
\pgfplotsset{
  every axis/.append style={axis line style={popink,line width=1pt},tick style={popink},
    grid style={popline,line width=0.5pt},label style={font=\small},tick label style={font=\footnotesize}},
  popcurve/.style={poporange,line width=1.8pt,no markers,smooth},
}
% Même style disponible dans un \draw TikZ simple (hors pgfplots)
\tikzset{popcurve/.style={poporange,line width=1.8pt}}
\tikzset{popgrid/.style={popline,very thin}}
\colorlet{popcurve}{poporange} % le nom du style est parfois utilisé comme couleur

% ============================================================
% Pill / squiggle
% ============================================================
\newcommand{\poppill}[3]{%
  \tikz[baseline=-0.55ex]\node[draw=popink,line width=1.2pt,rounded corners=2.6pt,%
    fill=#2,inner xsep=6.5pt,inner ysep=3pt]{%
    \popxboldls\fontsize{7}{7}\selectfont\color{#3}\MakeUppercase{#1}};%
}
\newcommand{\popsquiggle}[1]{%
  \tikz[baseline=-0.4ex]{
    \draw[#1,line width=2.6pt,line cap=round]
      plot [smooth] coordinates {(0,0.09) (0.34,0.01) (0.68,0.21) (1.02,0.01) (1.36,0.10)};
  }%
}
% Mot-clé surligné (marqueur orange sous le texte)
\newcommand{\cle}[1]{{\popxbold\color{popdark}#1}}

% ============================================================
% En-tête / pied
% ============================================================
\pagestyle{fancy}
\fancyhf{}
\renewcommand{\headrulewidth}{0pt}
\renewcommand{\footrulewidth}{0pt}
\lhead{\raisebox{-0.15ex}{\poplogo\fontsize{13}{13}\selectfont\color{popink}OnBuch\color{poporange}+}}
\rhead{\poppill{\DOCMATIERE\ \textbullet\ \DOCNIVEAU\ \textbullet\ Chapter \DOCLECON}{white}{popink}}
\lfoot{\popxbold\fontsize{7.6}{7.6}\selectfont\color{popmuted}\MakeUppercase{Signed course OnBuch+}\ \textbullet\ {\popsemi\DOCTITRE}}
\rfoot{\tikz[baseline=-0.6ex]\node[rounded corners=8.5pt,fill=popink,minimum height=17pt,minimum width=17pt,inner xsep=5pt,inner sep=0]{\popxbold\fontsize{8}{8}\selectfont\color{popcream}\thepage\,/\,\pageref*{LastPage}};}

% ============================================================
% Titres
% ============================================================
\newcommand{\chaptitre}[2]{%
  \begin{tcolorbox}[enhanced,colback=poporange,colframe=popink,boxrule=2.6pt,arc=11pt,
      drop shadow=popink,left=15pt,right=15pt,top=12pt,bottom=11pt,before skip=3pt,after skip=18pt]
    \poppill{#2}{popink}{popcream}\par\vspace{9pt}
    {\raggedright\hyphenpenalty=10000\exhyphenpenalty=10000\popdisplay\fontsize{22}{24}\selectfont\color{popcream}\MakeUppercase{#1}\par}
  \end{tcolorbox}
}
% Section numérotée : pastille orange avec le numéro + titre Archivo
\newcounter{popsec}
\newcommand{\coursec}[1]{%
  \stepcounter{popsec}\setcounter{popsub}{0}%
  \par\vspace{16pt}\noindent
  \begin{minipage}{\linewidth}
  \tikz[baseline=-0.7ex]\node[draw=popink,line width=1.6pt,fill=poporange,rounded corners=4pt,minimum size=22pt,inner sep=0]{\popdisplay\fontsize{12}{12}\selectfont\color{popcream}\thepopsec};%
  \hspace{8pt}{\popdisplay\fontsize{15}{17}\selectfont\color{popink}\MakeUppercase{#1}}\par
  \vspace{4pt}\noindent{\color{poporange}\rule{36pt}{3.6pt}}
  \end{minipage}\par\nopagebreak\vspace{8pt}
}
\newcounter{popsub}
\newcommand{\courssub}[1]{%
  \stepcounter{popsub}%
  \par\vspace{9pt}\noindent{\popxbold\fontsize{11.5}{13}\selectfont\color{popink}{\color{poporange}\thepopsec.\thepopsub}\hspace{6pt}#1}\par\nopagebreak\vspace{4pt}
}

% ============================================================
% Boîtes pédagogiques — même recette : fond blanc, bordure épaisse colorée,
% ombre dure encre, badge plein. Titre optionnel [#1].
% ============================================================
\tcbset{popbase/.style={breakable,enhanced,colback=white,boxrule=2.3pt,arc=9pt,
  shadow={3pt}{-3pt}{0pt}{popink},left=14pt,right=14pt,top=11pt,bottom=11pt,before skip=11pt,after skip=15pt,
  fontupper=\popsemi\fontsize{10.6}{14.5}\selectfont\color{popink},
  fontlower=\popsemi\fontsize{10.6}{14.5}\selectfont\color{popink}}}
% Coche dessinée (le glyphe ✓ n'existe pas dans les polices du document)
\newcommand{\popcheck}[1]{\tikz[baseline=-0.5ex]\draw[#1,line width=1.6pt,line cap=round,line join=round](-0.13,0.0)--(-0.04,-0.09)--(0.13,0.1);}
\providecommand{\coche}{\popcheck{popgreen}}
\providecommand{\resultat}[1]{\par\smallskip\noindent\fcolorbox{popgreen}{popgreenL}{\parbox{\dimexpr\linewidth-2\fboxsep-2\fboxrule\relax}{\textbf{Result.} #1}}\par\smallskip}
\newcommand{\popboxhead}[3]{\poppill{#1}{#2}{white}\ifx\relax#3\relax\else\hspace{6pt}{\popxbold\fontsize{10.6}{12}\selectfont\color{popink}#3}\fi\par\vspace{7pt}}

\newtcolorbox{aretenir}[1][]{popbase,colframe=popgreen,before upper={\popboxhead{Key points}{popgreen}{#1}}}
\newtcolorbox{exemplebox}[1][]{popbase,colframe=poporange,before upper={\popboxhead{Exemple}{poporange}{#1}}}
\newtcolorbox{definition}[1][]{popbase,colframe=popblue,before upper={\popboxhead{Definition}{popblue}{#1}}}
\newtcolorbox{propriete}[1][]{popbase,colframe=poppurple,before upper={\popboxhead{Property}{poppurple}{#1}}}
\newtcolorbox{methode}[1][]{popbase,colframe=popgold,before upper={\popboxhead{Method}{popgold}{#1}}}
\newtcolorbox{attention}[1][]{popbase,colframe=poppink,before upper={\popboxhead{Warning}{poppink}{#1}}}
\newtcolorbox{experience}[1][]{popbase,colframe=popblue,colback=popblueL,before upper={\popboxhead{Experiment}{popblue}{#1}}}
% « Le savais-tu ? » : carte violette pleine (contexte, culture, Cameroun)
\newtcolorbox{savaistu}[1][]{popbase,colback=poppurple,colframe=popink,
  fontupper=\popsemi\fontsize{10.4}{14.2}\selectfont\color{white},
  before upper={\poppill{Did you know?}{white}{poppurple}\ifx\relax#1\relax\else\hspace{6pt}{\popxbold\color{white}#1}\fi\par\vspace{7pt}}}
% Exercice résolu : énoncé en haut, \tcblower puis la solution sur fond crème
\newcounter{popexo}\newcounter{popexr}
\newtcolorbox{exoresolu}[1][]{popbase,colframe=poporange,colbacklower=popcream,
  segmentation style={popink,line width=1.2pt,dashed},
  before upper={\stepcounter{popexr}\popboxhead{Worked example \thepopexr}{poporange}{#1}},
  before lower={\poppill{Solution}{popink}{popcream}\par\vspace{6pt}}}
% Exercice (non résolu en place) — la correction est dans la partie Corrigés
\newtcolorbox{exercice}[1][]{popbase,colframe=popink,
  before upper={\stepcounter{popexo}\popboxhead{Exercise \thepopexo}{popink}{#1}}}
% Corrigé
\newtcolorbox{corrige}[1][]{popbase,colframe=popgreen,colback=popgreenL,
  before upper={\popboxhead{Solution}{popgreen}{#1}}}
% Objectifs / prérequis / bilan
\newtcolorbox{objectifs}{popbase,colframe=popink,colback=poporangeL,
  before upper={\popboxhead{Chapter objectives}{popink}{}}}
\newtcolorbox{prerequis}{popbase,colframe=popink,colback=popblueL,
  before upper={\popboxhead{Prerequisites}{popblue}{}}}
\newtcolorbox{bilan}{popbase,colframe=popink,colback=white,boxrule=2.8pt,
  before upper={\popboxhead{Fiche bilan}{poporange}{L'essentiel à connaître pour l'examen}}}
% Figure encadrée avec légende
\newtcolorbox{popfigure}[1][]{enhanced,breakable=false,colback=white,colframe=popline,boxrule=1.4pt,arc=8pt,
  left=10pt,right=10pt,top=10pt,bottom=8pt,before skip=10pt,after skip=12pt,halign=center,
  lower separated=false,halign lower=center,
  fontlower=\popsemi\fontsize{9}{11.5}\selectfont\color{popmuted},
  #1}
\newcommand{\legende}[1]{\par\vspace{6pt}{\popsemi\fontsize{9}{11.5}\selectfont\color{popmuted}#1}}

% ============================================================
% Signature OnBuch+
% ============================================================
\newcommand{\poplogoinline}[1]{{\poplogo\fontsize{#1}{#1}\selectfont\color{popink}OnBuch\color{poporange}+}}
\newcommand{\signatureonbuch}{%
  \begin{tcolorbox}[enhanced,colback=popink,colframe=popink,boxrule=2.2pt,arc=10pt,
      drop shadow=poporange,left=14pt,right=14pt,top=10pt,bottom=10pt,before skip=0pt,after skip=0pt]
    \tikz[baseline=-0.7ex]\node[circle,fill=popgreen,draw=popcream,line width=1.3pt,minimum size=22pt,inner sep=0]{\popcheck{white}};%
    \hspace{9pt}\begin{minipage}[c]{0.62\linewidth}
      {\popxbold\fontsize{12}{13}\selectfont\color{popcream}Signed\ \textbullet\ {\poplogo OnBuch\color{poporange}+}}\par\vspace{2pt}
      {\raggedright\popsemi\fontsize{8}{10.5}\selectfont\color{popline}\MakeUppercase{OnBuch+ teaching team}\\\MakeUppercase{Follows the official MINESEC syllabus}\par}
    \end{minipage}\hfill
    {\poplogo\fontsize{22}{22}\selectfont\color{popcream}OnBuch\color{poporange}+}
  \end{tcolorbox}%
}

% ============================================================
% Page de garde
% #1 titre · #2 matière · #3 niveau · #4 module · #5 n° leçon · #6 durée ·
% #7 description / objectifs (texte ou itemize)
% ============================================================
\newcommand{\pagegarde}[7]{%
  \thispagestyle{empty}
  \newgeometry{a4paper,margin=1.7cm,top=1.6cm,bottom=1.4cm}%
  \begin{tikzpicture}[remember picture,overlay]
    \fill[poporange!18] ([xshift=-3.2cm,yshift=-3.6cm]current page.north east) circle (4.6cm);
    \fill[poppurple!14] ([xshift=2.4cm,yshift=7.6cm]current page.south west) circle (3.8cm);
  \end{tikzpicture}%
  {\poplogo\fontsize{30}{30}\selectfont\color{popink}OnBuch\color{poporange}+}\hfill
  \raisebox{0.6ex}{\poppill{Signed course \textbullet\ Premium}{popink}{popcream}}\par
  \vspace{1pt}\popsquiggle{poppurple}\par
  \vspace{8pt}
  \begin{tcolorbox}[enhanced,colback=poporange,colframe=popink,boxrule=2.8pt,arc=13pt,
      drop shadow=popink,left=18pt,right=18pt,top=12pt,bottom=13pt,after skip=10pt]
    \poppill{#2 \textbullet\ #3}{popink}{popcream}\hspace{5pt}\poppill{#4}{white}{popink}\par\vspace{8pt}
    {\popdisplay\fontsize{14}{14}\selectfont\color{popink}CHAPTER #5}\par\vspace{4pt}
    {\raggedright\hyphenpenalty=10000\exhyphenpenalty=10000\popdisplay\fontsize{25}{27}\selectfont\color{popcream}\MakeUppercase{#1}\par}
    \vspace{8pt}
    {\popxbold\fontsize{9.5}{9.5}\selectfont\color{popink}\MakeUppercase{Suggested time: #6}}
  \end{tcolorbox}
  \begin{tcolorbox}[enhanced,colback=white,colframe=popink,boxrule=2.2pt,arc=10pt,
      drop shadow=popink,left=16pt,right=16pt,top=10pt,bottom=10pt,after skip=8pt]
    \poppill{In this chapter}{poppurple}{white}\par\vspace{5pt}
    \popsemi\fontsize{9.3}{12.6}\selectfont\color{popink}#7
  \end{tcolorbox}
  \noindent\begin{minipage}[t]{0.487\linewidth}
    \begin{tcolorbox}[enhanced,colback=popgreen,colframe=popink,boxrule=2.2pt,arc=10pt,
        drop shadow=popink,left=11pt,right=11pt,top=10pt,bottom=10pt,height=3.1cm,valign=center]
      \tcbox[colback=white,colframe=popink,boxrule=1.3pt,arc=5pt,boxsep=0pt,nobeforeafter,
        left=3pt,right=3pt,top=3pt,bottom=3pt,valign=center]{\qrcode[height=1.9cm]{https://play.google.com/store/apps/details?id=com.luvvix.onbuch&hl=en}}%
      \hspace{8pt}\begin{minipage}[c]{0.5\linewidth}
        {\popdisplay\fontsize{11}{12.5}\selectfont\color{white}\MakeUppercase{Download OnBuch}}\par\vspace{4pt}
        {\raggedright\popsemi\fontsize{8.4}{11}\selectfont\color{white}Free \textbullet\ Google Play\\AI tutor, past papers, results\par}
      \end{minipage}
    \end{tcolorbox}
  \end{minipage}\hfill
  \begin{minipage}[t]{0.487\linewidth}
    \begin{tcolorbox}[enhanced,colback=poppurple,colframe=popink,boxrule=2.2pt,arc=10pt,
        drop shadow=popink,left=11pt,right=11pt,top=10pt,bottom=10pt,height=3.1cm,valign=center]
      \tikz[baseline=-0.7ex]\node[circle,fill=white,draw=popink,line width=1.3pt,minimum size=1.6cm,inner sep=0]{\popdisplay\fontsize{24}{24}\selectfont\color{poppurple}+};%
      \hspace{8pt}\begin{minipage}[c]{0.6\linewidth}
        {\popdisplay\fontsize{11}{12.5}\selectfont\color{white}\MakeUppercase{Go OnBuch+}}\par\vspace{4pt}
        {\raggedright\popsemi\fontsize{8.4}{11}\selectfont\color{white}All signed courses, unlimited AI tutor, PDF sheets — OnBuch+ tab in the app\par}
      \end{minipage}
    \end{tcolorbox}
  \end{minipage}
  \vspace{8pt}
  \popcontenu
  \vfill
  \signatureonbuch
  \restoregeometry
  \newpage
}

% Rangée « Ce cours contient » (page de garde)
\newcommand{\poptile}[3]{% #1 couleur #2 gros texte #3 légende
  \begin{tcolorbox}[enhanced,colback=#1,colframe=popink,boxrule=2pt,arc=9pt,shadow={2.5pt}{-2.5pt}{0pt}{popink},
      left=6pt,right=6pt,top=9pt,bottom=9pt,halign=center,valign=center,height=1.75cm,width=\linewidth,nobeforeafter]
    {\popdisplay\fontsize{12.5}{13}\selectfont\color{white}\MakeUppercase{#2}}\par\vspace{3pt}
    {\centering\popsemi\fontsize{7.8}{9.5}\selectfont\color{white}#3\par}
  \end{tcolorbox}}
\newcommand{\popcontenu}{%
  \par\noindent{\popxboldls\fontsize{7.5}{7.5}\selectfont\color{popmuted}THIS SIGNED COURSE CONTAINS}\par\vspace{7pt}
  \noindent\begin{minipage}[t]{0.235\linewidth}\poptile{popblue}{Course}{complete, illustrated, step by step}\end{minipage}\hfill
  \begin{minipage}[t]{0.235\linewidth}\poptile{poporange}{Methods}{and worked examples}\end{minipage}\hfill
  \begin{minipage}[t]{0.235\linewidth}\poptile{poppink}{Exercises}{exam style + solutions}\end{minipage}\hfill
  \begin{minipage}[t]{0.235\linewidth}\poptile{popgreen}{Summary sheet}{the essentials to remember}\end{minipage}\par}

% Sommaire
\newtcolorbox{sommaire}{enhanced,colback=white,colframe=popink,boxrule=2.2pt,arc=10pt,
  drop shadow=popink,left=18pt,right=18pt,top=13pt,bottom=13pt,before skip=6pt,after skip=20pt,
  before upper={\poppill{Contents}{poppurple}{white}\par\vspace{9pt}\popsemi\fontsize{10.6}{14.5}\selectfont\color{popink}}}

% Signature de fin de cours — poussée tout en bas de la dernière page
\newcommand{\findecours}{%
  \par\vfill\nopagebreak
  \begin{center}\popsquiggle{poporange}\end{center}\vspace{4pt}
  \signatureonbuch
}
\AtBeginDocument{\def\llbracket{\mathopen{[\mkern-3.5mu[}}\def\rrbracket{\mathclose{]\mkern-3.5mu]}}} % intervalles d'entiers (glyphe ⟦ absent de la police)
% Glyphes absents de Fira Math : redéfinis à partir de glyphes disponibles
\AtBeginDocument{%
  \def\setminus{\mathbin{\backslash}}\let\smallsetminus\setminus
  \def\vdots{\mathord{\vbox{\baselineskip3.2pt\lineskiplimit0pt\kern4pt\hbox{.}\hbox{.}\hbox{.}}}}%
  \def\ddots{\mathinner{\mkern1mu\raise6.5pt\hbox{.}\mkern2mu\raise3.6pt\hbox{.}\mkern2mu\raise0.7pt\hbox{.}\mkern1mu}}%
  \def\triangle{\mathord{\scalebox{0.9}{$\symup{\Delta}$}}}%
  \def\nabla{\mathord{\raisebox{\depth}{\scalebox{1}[-1]{$\symup{\Delta}$}}}}%
  \def\checkmark{\popcheck{popdark}}%
  \def\star{\mathbin{\ast}}\def\bigstar{\mathord{\ast}}%
  \def\diamond{\mathbin{\scalebox{0.6}{$\Diamond$}}}%
  \def\bigcup{\mathop{\mathchoice{\raisebox{-0.3em}{\scalebox{1.7}{$\cup$}}}{\scalebox{1.25}{$\cup$}}{\cup}{\cup}}\displaylimits}%
  \def\bigcap{\mathop{\mathchoice{\raisebox{-0.3em}{\scalebox{1.7}{$\cap$}}}{\scalebox{1.25}{$\cap$}}{\cap}{\cap}}\displaylimits}%
  \def\lesseqqgtr{\mathrel{\lessgtr}}\def\bigoplus{\mathop{\oplus}\displaylimits}%
}
\providecommand{\ssi}{if and only if} % abréviation fréquente
\AtBeginDocument{\providecommand{\wideparen}{\overparen}} % arc de cercle

% Fonctions usuelles que les modèles utilisent sans que amsmath les fournisse
\providecommand{\arsinh}{\operatorname{arsinh}}\providecommand{\arcosh}{\operatorname{arcosh}}
\providecommand{\artanh}{\operatorname{artanh}}\providecommand{\arsech}{\operatorname{arsech}}
\providecommand{\arcsch}{\operatorname{arcsch}}\providecommand{\arcoth}{\operatorname{arcoth}}
\providecommand{\arcsinh}{\operatorname{arsinh}}\providecommand{\arccosh}{\operatorname{arcosh}}
\providecommand{\arctanh}{\operatorname{artanh}}\providecommand{\sech}{\operatorname{sech}}
\providecommand{\csch}{\operatorname{csch}}\providecommand{\arcsec}{\operatorname{arcsec}}
\providecommand{\arccsc}{\operatorname{arccsc}}\providecommand{\arccot}{\operatorname{arccot}}
\providecommand{\sgn}{\operatorname{sgn}}\providecommand{\lcm}{\operatorname{lcm}}
\providecommand{\Var}{\operatorname{Var}}\providecommand{\Cov}{\operatorname{Cov}}
\providecommand{\permil}{\ensuremath{{}^{0}\!/_{00}}}\providecommand{\textperthousand}{\ensuremath{{}^{0}\!/_{00}}}

%% FIGFIX-BEGIN v3 — correctifs globaux des figures (docs/onbuch-generation/tools/figfix)
\usepackage{adjustbox}\usepackage{etoolbox}
% toute figure TikZ/pgfplots plus large que la ligne est réduite à la largeur disponible
\BeforeBeginEnvironment{tikzpicture}{\begin{adjustbox}{max width=\linewidth}}
\AfterEndEnvironment{tikzpicture}{\end{adjustbox}}
% légendes pgfplots lisibles par-dessus les courbes
\pgfplotsset{every axis/.append style={legend style={fill=white,fill opacity=0.92,text opacity=1,draw=black!15,inner sep=3pt}}}
% étiquettes d'arêtes (syntaxe edge["..."]) avec fond blanc
\tikzset{every edge quotes/.append style={fill=white,inner sep=1.2pt}}
% bulles de cartes mentales : texte à la ligne dans la bulle, bulle un peu plus grande
\tikzset{every concept/.append style={text width=2.0cm,align=center,minimum size=2.3cm,inner sep=2pt}}
%% FIGFIX-END
\begin{document}

\pagegarde{HYDROGEOLOGY}{Geology}{Upper Sixth}{Upper Sixth — Geology}{7}{8 periods}{This chapter studies groundwater: its origin in the hydrologic cycle, its circulation and storage in rocks, and its natural or artificial outlets (springs, rivers, wells). It also examines the rock properties that control groundwater flow, the geometry of aquifers, and the threats of contamination, with direct applications to water supply in Cameroon.
\vspace{4pt}
\begin{multicols}{2}\small
\begin{itemize}
\item By the end of this chapter I can describe and draw the main stages of the hydrologic cycle with their water fluxes.
\item By the end of this chapter I can identify and explain the different sources of groundwater (meteoric, juvenile, connate, vadose).
\item By the end of this chapter I can describe the zones of underground water (aeration and saturation) and the factors controlling infiltration and circulation.
\item By the end of this chapter I can explain the formation and types of springs and the relationship between groundwater, rivers and oceans.
\item By the end of this chapter I can describe the structure, equipment and functioning of a water well and distinguish wells from boreholes.
\item By the end of this chapter I can define and compare porosity and permeability and relate them to rock types.
\end{itemize}\end{multicols}}

\begin{sommaire}
\begin{multicols}{2}
\begin{enumerate}
\item The Hydrologic Cycle and Sources of Groundwater
\item Underground Water Circulation and Accumulation
\item Springs, Rivers, Oceans and Water Wells
\item Rock Properties Controlling Groundwater Flow; Aquifers and Aquicludes
\item Water Quality and Groundwater Contamination
\item Integration activity
\item Methods and worked examples
\item Exercises
\item Solutions to the exercises
\item Summary sheet
\end{enumerate}
\end{multicols}
\end{sommaire}

\begin{objectifs}
\begin{itemize}
\item By the end of this chapter I can describe and draw the main stages of the hydrologic cycle with their water fluxes.
\item By the end of this chapter I can identify and explain the different sources of groundwater (meteoric, juvenile, connate, vadose).
\item By the end of this chapter I can describe the zones of underground water (aeration and saturation) and the factors controlling infiltration and circulation.
\item By the end of this chapter I can explain the formation and types of springs and the relationship between groundwater, rivers and oceans.
\item By the end of this chapter I can describe the structure, equipment and functioning of a water well and distinguish wells from boreholes.
\item By the end of this chapter I can define and compare porosity and permeability and relate them to rock types.
\item By the end of this chapter I can distinguish aquifers, aquicludes, aquitards and aquifuges and classify aquifers (unconfined, confined, perched, artesian).
\item By the end of this chapter I can state the criteria of water quality and identify the main sources and pathways of groundwater contamination.
\item By the end of this chapter I can interpret a hydrogeological cross-section to locate a favourable drilling site and propose protection measures.
\end{itemize}
\end{objectifs}

\begin{prerequis}
\begin{itemize}
\item Rock cycle and the three rock families (igneous, sedimentary, metamorphic) — Lower Sixth
\item Weathering, erosion and sedimentation processes — Lower Sixth
\item Basic properties of rocks: texture, structure, joints and bedding planes — Lower Sixth
\item The water states and simple notions of precipitation and evaporation from Geography — Form 5 / Lower Sixth
\item Reading and interpreting geological maps and cross-sections — Lower Sixth
\end{itemize}
\end{prerequis}

\newpage

\coursec{The Hydrologic Cycle and Sources of Groundwater}

\textbf{Introduction: the hidden water beneath our feet}\par

During the long dry season in Bamenda and Bafoussam, taps run dry while a borehole drilled a few streets away yields clear water all year round. Why does water hide underground in some places and not in others? Why does a spring gush out on the flank of Mount Cameroon while the plain below stays parched? And why are some wells in Douala's quarters contaminated while others remain safe? This chapter reveals the hidden journey of water beneath our feet and explains how geology controls where groundwater is stored, how it moves, and how we can exploit and protect it.

The answers to all these questions begin with one simple observation: \emph{water never stops moving}. The rain that fell on the Adamawa Plateau last month may today be flowing in the Sanaga River, evaporating from a maize field, or slowly percolating toward an aquifer tens of metres below the surface. Understanding this perpetual motion --- the \cle{hydrologic cycle} --- is the foundation of hydrogeology.

\courssub{Hydrogeology and Groundwater: Definitions and Importance}

\begin{definition}[Hydrogeology]
\cle{Hydrogeology} is the science of the \textbf{origin}, \textbf{circulation}, \textbf{distribution} and \textbf{quality} of underground water. It studies how water enters the ground, how it is stored and transmitted by rocks, and how it reappears at the surface in springs, rivers and wells.
\end{definition}

\begin{definition}[Groundwater]
\cle{Groundwater} is all water found beneath the Earth's surface, occupying pores, fractures and cavities in soils and rocks. The part that saturates the rock completely (below the \cle{water table}) is the water that wells, boreholes and springs can exploit.
\end{definition}

\begin{attention}[Groundwater is not an underground river]
A common misconception is that groundwater flows in vast underground rivers or lakes like those at the surface. In reality, except in karst limestone terrains, groundwater occupies the \textbf{tiny pores and cracks between mineral grains} and moves extremely slowly --- often only a few metres per year. Think of a wet sponge, not a buried pipe.
\end{attention}

\textbf{Why groundwater matters for Cameroon.}\par
Groundwater is a strategic resource for our country:

\begin{itemize}
\item \textbf{Drinking water:} in most towns of the North and Far North (Maroua, Garoua, Ngaoundéré), and in countless villages of the West and North-West, boreholes and wells are the main or only source of potable water, especially during the dry season when surface streams dry up.
\item \textbf{Agriculture:} irrigation of market gardens around Yaoundé, Bafoussam and in the Logone floodplain depends heavily on shallow groundwater.
\item \textbf{Industry:} breweries, food-processing plants and mining operations in Douala and elsewhere pump groundwater for their processes.
\item \textbf{Security:} unlike rivers, groundwater is naturally filtered and protected from rapid pollution and from evaporation, making it a reliable reserve in times of drought.
\end{itemize}

\courssub{The Global Hydrologic Cycle}

\textbf{Observation first.} After a heavy rain in Limbe, puddles disappear within hours under the sun; wet clothes dry on the line; clouds form over the ocean and drift toward Mount Cameroon, where they release rain. Nothing is lost, nothing is created: water simply changes state and place. This endless circulation is the \cle{hydrologic cycle} (or water cycle).

\begin{definition}[The hydrologic cycle]
The \cle{hydrologic cycle} is the continuous, closed circulation of water between the oceans, the atmosphere and the continents, driven by solar energy and gravity. It involves changes of state (evaporation, condensation) and transfers (precipitation, runoff, infiltration, groundwater flow).
\end{definition}

\textbf{The stages of the cycle, step by step:}

\begin{enumerate}
\item \textbf{Evaporation:} solar energy transforms liquid water from oceans, lakes, rivers and wet soils into water vapour. The oceans supply the great majority of atmospheric moisture.
\item \textbf{Transpiration:} plants absorb soil water through their roots and release vapour through their leaves. Evaporation and transpiration together are called \cle{evapotranspiration} (E). A dense forest like the Congo Basin rainforest returns enormous quantities of water to the atmosphere.
\item \textbf{Condensation:} rising vapour cools at altitude and condenses into droplets, forming clouds.
\item \textbf{Precipitation (P):} when droplets grow heavy enough, they fall as rain, hail or (rarely in Cameroon, on high peaks) snow. The seaward flank of Mount Cameroon is among the wettest places in Africa, receiving several metres of rain per year.
\item \textbf{Runoff (R):} on slopes and impermeable ground, part of the rain flows over the surface into streams and rivers, which carry it back to the ocean.
\item \textbf{Infiltration (I):} another part soaks into the soil.
\item \textbf{Percolation:} the infiltrated water migrates downward through the unsaturated zone until it reaches and recharges the saturated zone --- the groundwater reservoir.
\item \textbf{Groundwater flow and return to the ocean:} groundwater moves slowly through permeable rocks toward rivers (which it feeds during the dry season), springs and ultimately the sea, closing the cycle.
\end{enumerate}

\begin{definition}[Infiltration versus percolation]
\cle{Infiltration} is the \textbf{entry} of water into the soil at the surface --- the crossing of the air--soil boundary. \cle{Percolation} is the subsequent \textbf{downward migration} of that water through the unsaturated (vadose) zone toward the water table, under gravity. Infiltration happens at the surface; percolation happens beneath it.
\end{definition}

\begin{popfigure}
\begin{center}
\begin{tikzpicture}[scale=1.0, font=\small]
% Ocean
\fill[popblueL] (0,0) rectangle (5,-1.2);
\draw[popblue] (0,0) -- (5,0);
\node[popdark] at (2.5,-0.6) {Ocean};
% Land
\fill[popgreenL] (5,-1.2) -- (5,0) -- (7,0.9) -- (9,0.2) -- (11,0.6) -- (13,0) -- (13,-1.2) -- cycle;
\draw[popdark] (5,0) -- (7,0.9) -- (9,0.2) -- (11,0.6) -- (13,0);
% Water table (dashed)
\draw[popblue, dashed, thick] (5,-0.5) -- (13,-0.5);
\node[popblue, anchor=west] at (9.2,-0.78) {water table};
% Groundwater body
\fill[popblueL, opacity=0.6] (5,-0.5) rectangle (13,-1.2);
% Clouds
\draw[popmuted, thick] (2.2,3.4) ellipse (1.1 and 0.45);
\draw[popmuted, thick] (8.6,3.6) ellipse (1.3 and 0.5);
\node[popmuted] at (2.2,3.4) {cloud};
\node[popmuted] at (8.6,3.6) {cloud};
% Sun
\draw[popgold, thick] (0.4,4.3) circle (0.35);
\foreach \a in {0,45,...,315} \draw[popgold] (0.4,4.3) ++(\a:0.4) -- ++(\a:0.25);
% Evaporation arrows
\draw[->, poporange, very thick] (2.0,0.2) -- (2.0,2.6) node[midway, right, popdark]{evaporation};
\draw[->, poporange, thick] (10.2,0.85) -- (10.2,2.4) node[midway, right, popdark]{transpiration};
% Condensation label
\node[popmuted] at (5.2,4.1) {condensation};
\draw[->, popmuted, thick] (3.4,3.5) -- (7.0,3.6);
% Precipitation
\foreach \x in {8.0,8.5,9.0,9.5} \draw[->, popblue, thick] (\x,3.1) -- (\x,2.2);
\node[popblue] at (8.75,1.95) {precipitation};
% Runoff
\draw[->, popblue, very thick] (8.2,0.55) .. controls (7.2,0.35) .. (6.0,0.05) node[midway, above, popdark]{runoff};
% Infiltration / percolation
\draw[->, poppurple, very thick] (11.0,0.55) -- (11.0,-0.45) node[midway, right, popdark, align=left]{infiltration,\\percolation};
% Groundwater flow
\draw[->, poppurple, very thick] (11.5,-0.9) -- (5.4,-0.9) node[midway, below, popdark]{groundwater flow to ocean};
\end{tikzpicture}
\end{center}
\legende{The hydrologic cycle: solar-driven evaporation and transpiration feed clouds; precipitation returns water to the land, where it divides between runoff, infiltration and percolation; groundwater slowly flows back to the ocean.}
\end{popfigure}

\courssub{Water Reservoirs and Residence Times}

The total amount of water on Earth is essentially constant, but it is distributed among reservoirs of very different sizes, and water spends very different lengths of time in each. The \cle{residence time} is the average time a water molecule spends in a reservoir before moving to another.

\begin{center}
\begin{tabularx}{\linewidth}{|l|X|X|}
\hline
\rowcolor{popblueL} \textbf{Reservoir} & \textbf{Approximate share of Earth's water} & \textbf{Typical residence time} \\
\hline
Oceans and seas (salt water) & about 97\% of all water & thousands of years \\
\hline
Ice caps and glaciers & about 2\% (most of the fresh water) & tens to thousands of years \\
\hline
Groundwater & about 0.6--1\% (the bulk of \emph{liquid} fresh water) & days to thousands of years \\
\hline
Lakes and rivers & less than 0.02\% & days (rivers) to years (lakes) \\
\hline
Atmosphere & about 0.001\% & about 9--10 days \\
\hline
\end{tabularx}
\end{center}

\begin{popfigure}
\begin{center}
\begin{tikzpicture}
\begin{axis}[
    ybar, bar width=22pt,
    width=0.85\linewidth, height=6cm,
    ymin=0, ymax=100,
    ylabel={Share of Earth's water (\%)},
    symbolic x coords={Oceans, Ice, Groundwater, Surface fresh, Atmosphere},
    xtick=data,
    x tick label style={rotate=20, anchor=east, font=\footnotesize},
    ytick={0,20,40,60,80,100},
    nodes near coords, nodes near coords style={font=\footnotesize},
    every node near coord/.append style={anchor=west, rotate=90},
]
\addplot[fill=popblue, draw=popdark] coordinates {(Oceans,97) (Ice,2) (Groundwater,0.9) (Surface fresh,0.02) (Atmosphere,0.001)};
\end{axis}
\end{tikzpicture}
\end{center}
\legende{Distribution of Earth's water. Oceans dominate; among fresh water, groundwater far exceeds all lakes and rivers combined, which is why it is such a vital resource.}
\end{popfigure}

\begin{aretenir}[Key points on reservoirs]
\begin{itemize}
\item About \textbf{97\%} of Earth's water is salt water in the oceans; only about \textbf{3\%} is fresh.
\item Of that fresh water, most is locked in \textbf{ice}; the largest accessible \emph{liquid} fresh-water store is \textbf{groundwater}.
\item Rivers and the atmosphere hold tiny fractions but recycle very fast; groundwater recycles slowly, so over-pumping or polluting it has long-lasting consequences.
\end{itemize}
\end{aretenir}

\courssub{The Water Balance: $P = E + R + I$}

When rain falls on a catchment, every drop must go somewhere. It can return to the atmosphere by evapotranspiration, run off the surface into rivers, or infiltrate into the ground. This simple accounting gives the water balance.

\begin{propriete}[The water balance equation]
For a given area and period, precipitation is shared between evapotranspiration, runoff and infiltration:
\[ P = E + R + I \]
where $P$ = precipitation, $E$ = evapotranspiration, $R$ = surface runoff, $I$ = infiltration (all expressed in millimetres of water over the period). Only the term $I$ can recharge groundwater. (Proof not examinable; the equation is a direct statement of conservation of mass.)
\end{propriete}

\textbf{Seasonal variation in tropical climates.} In Cameroon's humid South (Yaoundé, Douala), rainfall is abundant and spread over much of the year, so $I$ remains significant for many months and aquifers recharge regularly. In the Sudano-Sahelian North (Maroua, Kousséri), a short rainy season concentrates $P$ into a few months, while potential evapotranspiration is very high all year; $E$ dominates the balance, $R$ is brief and violent (flash floods), and $I$ is small --- which is why groundwater there is precious and boreholes are essential.

\begin{exoresolu}[Water balance for a catchment near Bafoussam]
Statement: A small catchment on the Western Highlands receives an annual rainfall $P = 1800$ mm. Measurements and estimates give annual evapotranspiration $E = 1100$ mm and surface runoff $R = 450$ mm. (a) Calculate the annual infiltration $I$. (b) What percentage of the rainfall recharges the groundwater?
\tcblower
Solution:
\begin{align*}
P &= E + R + I \\
I &= P - E - R = 1800 - 1100 - 450 = 250\ \text{mm}
\end{align*}
(b) The recharge percentage is
\[ \frac{I}{P} \times 100 = \frac{250}{1800} \times 100 \approx 13.9\ \%. \]
Interpretation: only about 14\% of the rain actually feeds the aquifer. Even in a humid highland climate, most rainfall returns to the atmosphere or runs off --- a good reason to protect recharge zones from deforestation and erosion.
\end{exoresolu}

\courssub{The Sources (Origins) of Groundwater}

Where does the water in an aquifer come from? Hydrogeologists recognise four origins.

\begin{definition}[Sources of groundwater]
\begin{itemize}
\item \cle{Meteoric water}: water from rain and snow that infiltrates the soil and percolates to the water table. This is \textbf{by far the dominant source} of exploitable groundwater worldwide and in Cameroon.
\item \cle{Connate water} (fossil water): seawater or lake water trapped in the pores of sediments at the time of their deposition and burial, e.g. saline water preserved in deep sedimentary formations of the Douala basin or the Benue trough. It is often highly mineralised and unfit for drinking.
\item \cle{Juvenile water} (magmatic water): water released from magma during cooling and volcanic activity, reaching the surface for the first time. It contributes to hot springs and fumaroles, for example along the Cameroon Volcanic Line, but its volume is very small.
\item \cle{Vadose water}: water held temporarily in the \textbf{unsaturated zone} above the water table (soil moisture available to plants). It is groundwater in the broad sense but is not part of the permanent saturated reserve.
\end{itemize}
\end{definition}

\begin{attention}[Common mistake: ``most groundwater is fossil'']
Many students believe that aquifers are filled with ancient, fossil water trapped long ago. In fact, \textbf{nearly all exploitable groundwater is of meteoric origin}: it is rain that infiltrated recently (years to centuries ago) and is continually renewed by recharge. Connate water is usually too saline to use, and juvenile water is negligible in quantity. Never write in an exam that groundwater is ``mostly fossil''.
\end{attention}

\courssub{What Controls How Much Rain Infiltrates?}

Two neighbouring plots can receive the same storm yet recharge their aquifers very differently. The share of rainfall that becomes infiltration depends on four main factors:

\begin{itemize}
\item \textbf{Climate:} gentle, prolonged rain infiltrates well; short, violent storms (typical of the Far North) exceed the soil's intake capacity and mostly run off. High temperatures raise $E$, reducing the water available for $I$.
\item \textbf{Vegetation cover:} forest and grassland slow runoff, protect the soil from crusting, and their roots open pathways for water. Deforested, baked soils in the Sudano-Sahelian zone crust over and shed water like a roof.
\item \textbf{Slope:} on steep slopes (flanks of Mount Cameroon, the Bamenda escarpment) water runs off rapidly and infiltration time is short; on gentle plains water ponds and infiltrates.
\item \textbf{Soil type:} sandy, porous soils (e.g. parts of the Logone plain) admit water readily; clayey or compacted lateritic soils are nearly impermeable and favour runoff.
\end{itemize}

\begin{exemplebox}[Humid South versus Sudano-Sahelian North]
Around Ebolowa in the humid South, abundant and well-distributed rain falls on forested, gentle terrain: infiltration is high and shallow wells rarely dry up. Around Maroua, the same annual logic fails: a few intense storms fall on sparsely vegetated, often clay-crusted soils, so most water runs off to the Mandara streams or evaporates, and only deep boreholes tapping regional aquifers provide a reliable supply.
\end{exemplebox}

\begin{methode}[GCE exam tip: sketching the hydrologic cycle in under 5 minutes]
\begin{enumerate}
\item Draw the sea on the left and a simple land profile (with one hill) on the right; add a dashed line for the water table.
\item Add a sun, one or two clouds, and rain arrows falling on the land.
\item Draw and \textbf{label} the arrows: evaporation (up from the sea), transpiration (up from vegetation), condensation (in the cloud), precipitation (down), runoff (along the surface to the sea), infiltration/percolation (down to the water table) and groundwater flow (below ground toward the sea).
\item Check that every label is spelled correctly --- examiners penalise unlabelled or mislabelled arrows.
\end{enumerate}
Practise this sketch until you can reproduce it neatly in less than five minutes; it is a classic GCE question.
\end{methode}

\begin{savaistu}[Did you know?]
Lake Nyos, in the North-West Region, reminds us that groundwater and deep Earth fluids are connected: gases of magmatic origin dissolved in deep waters accumulated in the lake and were catastrophically released in 1986. Studying underground water and gases is therefore not only about water supply --- it is also a matter of public safety.
\end{savaistu}

\begin{exercice}[Water balance and recharge]
A catchment near Garoua receives $P = 900$ mm of rain per year. Evapotranspiration is estimated at $E = 720$ mm and infiltration at $I = 60$ mm.
\begin{enumerate}
\item Calculate the annual runoff $R$.
\item Express infiltration as a percentage of precipitation.
\item Give two reasons, linked to the local climate and soils, why infiltration is so much lower here than near Bafoussam.
\end{enumerate}
\end{exercice}

\begin{corrige}[Exercise 1]
\begin{enumerate}
\item From $P = E + R + I$: $R = P - E - I = 900 - 720 - 60 = 120$ mm.
\item $\dfrac{I}{P} \times 100 = \dfrac{60}{900} \times 100 \approx 6.7\ \%$.
\item The short rainy season concentrates rain into violent storms that exceed the soil's infiltration capacity, so runoff dominates; and the sparse vegetation plus clay-crusted, baked soils of the Sudano-Sahelian zone prevent water from entering the ground, while very high evapotranspiration removes most of the rest.
\end{enumerate}
\end{corrige}

\begin{aretenir}[To remember from this section]
\begin{itemize}
\item Hydrogeology studies the origin, circulation, distribution and quality of underground water; groundwater fills pores and fractures, not underground rivers.
\item The hydrologic cycle links oceans, atmosphere and land through evaporation, transpiration, condensation, precipitation, runoff, infiltration and percolation.
\item Oceans hold about 97\% of Earth's water; groundwater is the largest accessible liquid fresh-water reservoir.
\item The water balance $P = E + R + I$ shows that only the infiltrated fraction recharges aquifers; this fraction depends on climate, vegetation, slope and soil type.
\item Almost all exploitable groundwater is meteoric (infiltrated rain); connate and juvenile waters are minor and often unusable.
\end{itemize}
\end{aretenir}

\coursec{Underground Water Circulation and Accumulation}

In the previous section we followed water through the hydrologic cycle and saw that part of the rain falling on Cameroonian soil does not run off into rivers but \emph{infiltrates} into the ground. A natural question then arises: where does this infiltrated water go, how does it move in the dark beneath our feet, and why can a well dug in one village yield abundant water while another, a few hundred metres away, stays dry? Answering these questions is the object of this section. We shall first describe the vertical arrangement of subsurface water, then study how groundwater circulates from highlands to lowlands, and finally establish the conditions under which water accumulates to form an exploitable underground reservoir.

\courssub{Vertical Zonation of Subsurface Water}

\textbf{An observation to start with.} Anyone who has watched a well being dug near Buea or Bamenda has noticed the following sequence: the first metres of excavated material are damp but no water flows; deeper down, the walls become wetter and wetter; finally, at a certain depth, water begins to seep in and accumulate at the bottom of the hole. This simple observation reveals that the subsurface is divided into two great zones.

\begin{definition}[Zone of aeration and zone of saturation]
The \cle{zone of aeration} (or \emph{unsaturated zone}, also called the \emph{vadose zone}) is the upper part of the subsurface in which the pores of the soil and rock contain \textbf{both air and water}. Beneath it lies the \cle{zone of saturation} (or \emph{phreatic zone}), in which \textbf{all pores, cracks and fissures are completely filled with water}. The water held in the zone of saturation is called \emph{groundwater} strictly speaking, or phreatic water.
\end{definition}

Within the zone of aeration, water exists in three physical states, distinguished by the forces that hold it:

\begin{itemize}
\item \textbf{Hygroscopic water}: a very thin film of water bound tightly to the surface of soil grains by molecular attraction. It cannot move under gravity and cannot be absorbed by plant roots.
\item \textbf{Capillary water}: water held in the narrow spaces between grains by \emph{capillary forces} (surface tension). It can move slowly upward or sideways and is the main source of moisture for plants.
\item \textbf{Gravitational water} (vadose water in transit): water that, after rainfall, percolates \emph{downward} through the pores under the pull of gravity, on its way to the saturation zone. It is temporary: given time, it drains away.
\end{itemize}

Just above the top of the saturation zone lies the \cle{capillary fringe}: a transition belt in which water is drawn \emph{upward} from the saturated zone by capillarity, rather as paraffin rises in the wick of a lamp. The pores there are essentially full of water, but the water is held under tension and will not flow freely into a well. The thickness of the fringe depends on grain size: it is thin (a few centimetres) in coarse sands and gravels, but may reach a metre or more in fine silts and clays, because capillary rise is greater in narrower pores.

\courssub{The Water Table}

\begin{definition}[Water table]
The \cle{water table} (or \emph{phreatic surface}) is the upper surface of the zone of saturation, defined as the level at which the pressure of the groundwater is exactly equal to \emph{atmospheric pressure}. Below the water table, water pressure exceeds atmospheric pressure and increases with depth; above it, in the aeration zone, water is held at pressures below atmospheric.
\end{definition}

A well that reaches the saturation zone fills with water up to the level of the water table: the water level observed in a well is therefore a direct measurement of the local position of the water table.

\textbf{Morphology: a subdued replica of the topography.} The water table is \emph{not} flat. Because groundwater flows from high areas toward low areas, the water table is generally higher beneath hills and lower beneath valleys: it \emph{follows the relief, but in a softened, subdued way}. Under a hill the water table may lie tens of metres below the surface; near a river it approaches the surface and may even intersect it, which is precisely why rivers flow. Where the water table cuts the land surface, water emerges: springs, marshes, lakes and rivers are the visible outcrops of the water table.

\textbf{Seasonal fluctuations.} The water table is not fixed in time. In Cameroon, with its marked alternation of rainy and dry seasons, the water table \emph{rises during the rainy season} (abundant recharge by infiltrating rain) and \emph{falls during the dry season} (drainage toward rivers continues while recharge stops). The difference between the high (wet-season) and low (dry-season) levels may reach several metres. This is why some shallow wells that yield water in October dry up completely in February: they were dug only down to the wet-season water table.

\textbf{Perched water tables.} A local lens of impermeable clay lying \emph{within} the zone of aeration can arrest percolating water and hold a small, isolated body of saturated ground \emph{above} the main water table. This is called a \cle{perched water table}. It explains why a shallow spring may flow halfway up a hillside even though the regional water table lies far deeper. Perched bodies are usually small and may dry up in the dry season.

\begin{attention}[Common mistake]
Do not confuse the water table with the \emph{top of the zone of aeration} (that top is simply the ground surface!) and do not assume the water table is horizontal. The water table is the \emph{top of the zone of saturation}; it is a curved surface, a subdued replica of the relief, and it fluctuates seasonally. Also remember that the capillary fringe lies \emph{above} the water table but is saturated: saturation alone does not define the water table --- the criterion is \emph{pressure equal to atmospheric pressure}.
\end{attention}

\begin{popfigure}
\begin{tikzpicture}[scale=1.0, font=\small]
% ground surface
\draw[thick, popdark] (0,3.2) .. controls (2,4.4) and (4,4.6) .. (6,3.4) .. controls (8,2.6) and (10,2.4) .. (12,2.8);
% fill aeration
\fill[popgold!25] (0,3.2) .. controls (2,4.4) and (4,4.6) .. (6,3.4) .. controls (8,2.6) and (10,2.4) .. (12,2.8)
 -- (12,1.7) .. controls (10,1.3) and (8,1.5) .. (6,2.2) .. controls (4,3.1) and (2,2.9) .. (0,1.9) -- cycle;
% capillary fringe
\fill[popblue!35] (0,1.9) .. controls (2,2.9) and (4,3.1) .. (6,2.2) .. controls (8,1.5) and (10,1.3) .. (12,1.7)
 -- (12,1.25) .. controls (10,0.85) and (8,1.05) .. (6,1.75) .. controls (4,2.65) and (2,2.45) .. (0,1.45) -- cycle;
% saturation zone
\fill[popblue!75] (0,1.45) .. controls (2,2.45) and (4,2.65) .. (6,1.75) .. controls (8,1.05) and (10,0.85) .. (12,1.25)
 -- (12,-1.2) -- (0,-1.2) -- cycle;
% water table (wet season)
\draw[very thick, popblue] (0,1.9) .. controls (2,2.9) and (4,3.1) .. (6,2.2) .. controls (8,1.5) and (10,1.3) .. (12,1.7);
% water table (dry season)
\draw[very thick, poporange, dashed] (0,1.45) .. controls (2,2.45) and (4,2.65) .. (6,1.75) .. controls (8,1.05) and (10,0.85) .. (12,1.25);
% impermeable bedrock
\fill[pattern=north east lines, pattern color=popmuted] (0,-1.2) rectangle (12,-1.9);
\draw[thick, popmuted] (0,-1.2) -- (12,-1.2);
% labels
\node[popdark, align=center] at (3,5.0) {Ground surface (topography)};
\node[align=center] at (9.3,3.6) {Zone of aeration\\ (air + water in pores)};
\node[align=center, popblue!60!black] at (9.6,0.2) {Zone of saturation\\ (pores full of water)};
\node[align=center, popmuted] at (6,-1.55) {Impermeable bedrock (substratum)};
% arrows and annotations
\draw[->, thick, popblue] (11.3,2.6) -- (11.3,2.0);
\node[popblue, align=center] at (11.3,2.85) {Water table (wet season)};
\draw[->, thick, poporange] (11.3,0.6) -- (11.3,1.1);
\node[poporange, align=center] at (11.3,0.35) {Water table (dry season)};
\node[align=center, popblue!50!black] at (2.2,2.15) {Capillary fringe};
% well
\draw[thick, popdark] (4.6,4.55) -- (4.6,1.0) (5.0,4.55) -- (5.0,1.0);
\fill[popblue!75] (4.6,1.0) rectangle (5.0,2.7);
\node[align=center, popdark] at (4.8,5.0) {Well};
\draw[->, thick, popdark] (5.3,2.4) -- (5.0,2.55);
\node[align=center, popdark, right] at (5.3,2.35) {water level in well = water table};
\end{tikzpicture}
\legende{Cross-section showing the vertical zonation of subsurface water: zone of aeration, capillary fringe, water table (wet- and dry-season positions) and zone of saturation resting on an impermeable substratum. Note that the water table follows the relief in a subdued manner.}
\end{popfigure}

\courssub{Movement of Groundwater: Recharge, Transit and Discharge}

\textbf{Intuition.} Water in the ground obeys the same master as water on the surface: \emph{gravity}. Just as rainwater runs downhill on a slope, groundwater percolates from places where the water table is high (beneath hills and highlands) toward places where it is low (beneath valleys, rivers and ultimately the sea). The driving force is the difference of \emph{hydraulic head} between two points. For the slow flows considered here, the hydraulic head at a point is, to an excellent approximation, simply the \emph{altitude of the water table} at that point (measured above sea level): water flows from points of higher head to points of lower head.

This allows us to divide any landscape into three functional areas:

\begin{itemize}
\item \textbf{Recharge areas}: highlands and permeable outcrops where infiltration exceeds drainage; water enters the ground and the water table is relatively high (e.g.\ the flanks of Mount Cameroon or the Bamenda Highlands during the rainy season).
\item \textbf{Transit areas}: the intermediate zones through which groundwater flows slowly, following curved paths that dive down and then rise again toward the outlets.
\item \textbf{Discharge areas}: valleys, river beds, lake shores, springs and the sea floor, where groundwater leaves the ground and rejoins the surface water network.
\end{itemize}

\begin{definition}[Hydraulic gradient]
The \cle{hydraulic gradient} $i$ is the slope of the water table between two points: the difference of hydraulic head $\Delta h$ divided by the horizontal distance $L$ between the points:
\[
i = \frac{\Delta h}{L}
\]
It is a dimensionless number (e.g.\ $i = 0.01$ means the water table drops $1$ m over $100$ m). Groundwater flows in the direction of \emph{decreasing head}, and the greater the gradient, the faster the flow.
\end{definition}

\textbf{Flow paths are curved, not straight.} Because recharge occurs over the whole hill while discharge is concentrated at the river, flow lines curve: they descend beneath the recharge area, travel nearly horizontally at depth, and rise toward the discharge zone. Flow lines are always \emph{perpendicular to the contours of the water table} (lines of equal head), just as a stream flows perpendicular to topographic contour lines.

\begin{popfigure}
\begin{tikzpicture}[scale=1.0, font=\small]
% ground
\draw[thick, popdark] (0,4.2) .. controls (2,5.0) and (4,5.0) .. (6,3.4) .. controls (8,2.2) and (10,1.8) .. (12,1.9);
% river
\fill[popblue] (10.6,1.9) .. controls (11.0,2.15) and (11.4,2.15) .. (11.8,1.9) -- (11.8,1.55) .. controls (11.4,1.75) and (11.0,1.75) .. (10.6,1.55) -- cycle;
\node[popblue, align=center] at (11.2,2.6) {River (discharge)};
% water table
\draw[very thick, popblue] (0,2.9) .. controls (2,3.6) and (4,3.6) .. (6,2.6) .. controls (8,1.7) and (10,1.5) .. (11.2,1.7);
\node[popblue] at (1.6,3.15) {Water table};
% saturation fill
\fill[popblue!18] (0,2.9) .. controls (2,3.6) and (4,3.6) .. (6,2.6) .. controls (8,1.7) and (10,1.5) .. (11.2,1.7)
 -- (11.2,0.2) -- (0,0.2) -- cycle;
% flow lines
\draw[->, thick, poppurple] (1.5,4.35) -- (1.5,3.7);
\draw[thick, poppurple] (1.5,3.7) .. controls (1.6,2.2) and (3,0.8) .. (6,0.8) .. controls (9,0.8) and (10.6,1.2) .. (11.1,1.6);
\draw[->, thick, poppurple] (3.5,4.6) -- (3.5,3.9);
\draw[thick, poppurple] (3.5,3.9) .. controls (3.6,2.8) and (5,1.6) .. (7,1.5) .. controls (9,1.4) and (10.5,1.5) .. (11.0,1.65);
\draw[->, thick, poppurple] (5.5,4.05) -- (5.5,3.3);
\draw[thick, poppurple] (5.5,3.3) .. controls (6,2.6) and (8,2.0) .. (9.6,1.85) .. controls (10.3,1.8) and (10.8,1.75) .. (11.0,1.72);
% recharge label
\node[align=center, popgreen!60!black] at (2.6,5.5) {Recharge area (highland)};
\draw[->, popgreen!60!black] (2.6,5.25) -- (2.6,4.75);
% rain
\draw[->, popblue] (0.8,5.6) -- (0.8,4.9);
\draw[->, popblue] (4.6,5.7) -- (4.6,5.0);
\node[popblue] at (4.6,5.95) {rain};
% hydraulic gradient annotation
\draw[<->, thick, poporange] (1.5,3.62) -- (11.0,1.68);
\node[poporange, align=center] at (6.4,3.35) {head difference $\Delta h$ over distance $L$: $i=\Delta h / L$};
% labels
\node[align=center, popdark] at (6.0,0.35) {Transit zone: curved flow lines};
\end{tikzpicture}
\legende{Groundwater flow from a recharge hill to a discharge river. Flow lines (purple) leave the recharge area, curve downward and rise toward the river, always perpendicular to water-table contours. The hydraulic gradient $i=\Delta h/L$ drives the flow.}
\end{popfigure}

\textbf{The role of fractures, joints and karst.} In many rocks, water does not move through the pores between grains but along \emph{discontinuities}: joints and fractures in crystalline rocks (granites, gneisses of the Congo craton margin), cooling joints and fissures in the basalts of the Cameroon Volcanic Line, and --- most spectacularly --- \emph{karst conduits} in limestone. In limestone terrains, slightly acidic rainwater (charged with dissolved carbon dioxide) dissolves the rock along cracks, progressively enlarging them into channels, caves and underground rivers. Flow is then rapid and concentrated along these conduits, which explains the large springs and the sudden disappearance of streams into sinkholes typical of karst regions. Fractured volcanic rocks similarly channel groundwater along fissure networks, which is why well siting on the slopes of Mount Cameroon requires locating fracture zones.

\courssub{Accumulation of Groundwater: Conditions for a Reservoir}

Circulating water is of little use unless it can \emph{accumulate} in sufficient quantity to supply wells and springs. Two conditions must be met simultaneously:

\begin{enumerate}
\item A \textbf{porous and permeable rock body} must be present, capable of \emph{storing} water in its pores or fissures and of \emph{transmitting} it to a well (such a rock is called an \emph{aquifer} --- studied in detail in Section 4).
\item An \textbf{impermeable substratum} (a bed of clay, unfractured granite or massive rock) must underlie or bound the permeable body, preventing the water from draining away to greater depth.
\end{enumerate}

The geometry is exactly analogous to a sponge resting on a plate: the sponge stores and yields water; the plate stops it from escaping downward. Where the impermeable layer instead \emph{overlies} the aquifer, the water becomes confined under pressure (an \emph{artesian} situation, developed in the next section).

\courssub{Quantifying the Flow: Darcy's Law (Extension)}

The qualitative statement ``water flows faster where the gradient is steep and the rock is permeable'' was turned into a precise law by the French engineer Henry Darcy (1856), from experiments on water percolating through columns of sand.

\begin{propriete}[Darcy's law]
The discharge $Q$ (volume of water per unit time) through a porous medium is proportional to the cross-sectional area $S$ of flow and to the hydraulic gradient $i$:
\[
Q = K \cdot S \cdot i
\]
where $K$ is the \cle{hydraulic conductivity} (coefficient of permeability), a constant characterising the rock, expressed in $\mathrm{m\,s^{-1}}$. High $K$: gravels, sands, fractured basalt, karst limestone. Low $K$: clays, unfractured crystalline rocks. \emph{(Extension useful for GCE problem-solving; the formula is applied directly, no proof required.)}
\end{propriete}

\begin{exoresolu}[Applying Darcy's law]
Statement. A sand aquifer of hydraulic conductivity $K = 2\times 10^{-4}\ \mathrm{m\,s^{-1}}$ is $10$ m thick and $200$ m wide. The water table drops by $4$ m over a horizontal distance of $800$ m. Calculate the groundwater discharge through the aquifer section.
\tcblower
Solution. The cross-section of flow is $S = 10 \times 200 = 2000\ \mathrm{m^2}$. The hydraulic gradient is
\[
i = \frac{\Delta h}{L} = \frac{4}{800} = 0.005.
\]
Applying Darcy's law:
\[
Q = K \cdot S \cdot i = 2\times 10^{-4} \times 2000 \times 0.005 = 2\times 10^{-3}\ \mathrm{m^3\,s^{-1}}.
\]
The aquifer delivers $Q = 0.002\ \mathrm{m^3\,s^{-1}}$, i.e.\ $2$ litres per second --- enough to supply a small village water scheme.
\end{exoresolu}

\begin{methode}[Deducing groundwater flow direction from water-table contours]
Given a map with contours of the water table (lines of equal head, in metres):
\begin{enumerate}
\item Identify the \emph{high} contours (recharge zones, usually beneath hills) and the \emph{low} contours (discharge zones, usually along rivers).
\item Draw flow lines \emph{perpendicular to the contours}, since water moves along the steepest slope of the water table.
\item Orient each flow line from \emph{high head toward low head} (downslope of the water table).
\item To quantify, measure $\Delta h$ between two contours and the map distance $L$ (converted with the scale), then compute $i = \Delta h / L$.
\end{enumerate}
\end{methode}

\begin{savaistu}[Groundwater in Cameroon]
Much of rural Cameroon depends directly on groundwater: hand-dug wells and boreholes tapping weathered and fractured crystalline basement supply villages across the Adamawa Plateau and the Far North, while the deeply weathered volcanic soils of the West and North-West Regions give numerous small springs used for domestic water. Understanding where the water table lies, and how it falls in the dry season, is therefore not an academic luxury but a practical necessity for every water-supply project in the country.
\end{savaistu}

\begin{aretenir}[Key points]
\begin{itemize}
\item The subsurface comprises the \emph{zone of aeration} (pores with air + hygroscopic, capillary and gravitational water) and the \emph{zone of saturation} (pores full of water), separated by the \emph{water table}, where pressure equals atmospheric pressure; a \emph{capillary fringe} sits just above it.
\item The water table is a \emph{subdued replica of the topography} and \emph{fluctuates seasonally} (high in the rainy season, low in the dry season); a clay lens within the aeration zone may hold a \emph{perched water table}.
\item Groundwater flows by gravity from \emph{recharge} areas (highlands) through \emph{transit} zones to \emph{discharge} areas (valleys, rivers, sea), along curved flow lines perpendicular to water-table contours, driven by the hydraulic gradient $i = \Delta h / L$.
\item Fractures, joints and karst dissolution channels concentrate and accelerate flow, especially in limestone and volcanic terrains.
\item Accumulation requires a porous, permeable rock body resting on an \emph{impermeable substratum}.
\item Darcy's law $Q = K \cdot S \cdot i$ quantifies groundwater discharge.
\end{itemize}
\end{aretenir}

\begin{exercice}[Water table and flow]
On a topographic cross-section through a hill and a river valley, two observation wells $A$ (on the hill) and $B$ (near the river) show water levels at $812$ m and $796$ m altitude respectively. The horizontal distance $AB$ is $400$ m. (1) Compute the hydraulic gradient between $A$ and $B$. (2) State the direction of groundwater flow. (3) Explain why a well dug to $800$ m altitude near the river yields water while a well of the same depth on the hill remains dry.
\end{exercice}

\begin{corrige}[Exercise 1]
(1) $i = \Delta h / L = (812 - 796)/400 = 16/400 = 0.04$. (2) Flow is from $A$ toward $B$, i.e.\ from the hill (high head, recharge) toward the river (low head, discharge), perpendicular to the water-table contours. (3) The water table is a subdued replica of the relief: near the river it lies close to the surface (around $796$ m), so a well reaching $800$ m is below it and fills; on the hill the water table lies near $812$ m, far deeper, so a well bottom at $800$ m is still \emph{above} the water table, inside the zone of aeration, and stays dry.
\end{corrige}

\coursec{Springs, Rivers, Oceans and Water Wells}

In the previous section we saw how rainwater infiltrates, percolates through pores and fractures, and finally accumulates as groundwater stored in aquifers. But groundwater does not remain hidden forever. Wherever the water table meets the land surface, groundwater reappears: as a trickle on a hillside in Bamenda, as the dry-season flow of the River Wouri, or as water drawn from a village well in the Sahelian North. This section studies these \textbf{points of contact between groundwater and the surface world}: springs, rivers, the ocean, and wells.

\courssub{Springs: Natural Outflows of Groundwater}

\textbf{Intuition.} Walk along a valley side in the Western Highlands after the rainy season and you will often see water seeping out of the slope, feeding a small stream. Nobody dug a hole; the water simply ``comes out''. Why here and not elsewhere? Because at that precise spot, the underground water body has been \emph{intercepted} by the topography.

\begin{definition}[Spring]
A \cle{spring} is a natural emergence (outflow) of groundwater at the ground surface, occurring where the \cle{water table} (or a water-bearing fracture) intersects the topography, or where groundwater is forced to the surface along a geological structure such as an impermeable contact or a fault.
\end{definition}

\begin{propriete}[Conditions for a spring to exist]
Two conditions must be satisfied simultaneously:
\begin{enumerate}
\item an \textbf{aquifer} must be present and supplied by infiltration (recharge);
\item the \textbf{land surface must cut across the water table} or across a structure (impermeable contact, fault, joint) that channels groundwater to the surface.
\end{enumerate}
This is why springs are common on \textbf{valley sides, at the foot of escarpments, and along fault lines}, and rare on flat plateaux where the water table lies deep.
\end{propriete}

\courssub{Types of Springs}

The geometry of the aquifer and its boundaries determines the type of spring. Four types are examinable at GCE Advanced Level, plus two special cases.

\textbf{(a) Depression (overflow) spring.} The water table slopes gently and simply \emph{crosses} the ground surface in a depression or valley. Water seeps out over a broad zone. Such springs are often weak and may dry up in the dry season when the water table falls below the surface.

\textbf{(b) Contact spring.} A permeable aquifer overlies an \textbf{impermeable layer (aquiclude)} whose outcrop is cut by the topography. Water moving down through the aquifer is stopped by the aquiclude, flows along the contact, and emerges where the contact reaches the surface. Contact springs typically form a \emph{line of springs} along a hillside at the same stratigraphic level.

\textbf{(c) Fault spring.} A \textbf{fault} brings an aquifer into contact with an impermeable rock, or creates a fractured channel. Groundwater migrates along the fault plane and emerges along the fault line. Fault springs may be deep-sourced and therefore \emph{warm}, because the water has circulated at depth where the temperature is higher (geothermal gradient).

\textbf{(d) Artesian spring.} In a \textbf{confined aquifer}, water is under pressure: it is sandwiched between two aquicludes and recharged at a higher elevation. Where a fracture or fault reaches the confined aquifer, pressure drives the water up to the surface, sometimes as a vigorous jet. The spring flows because the \textbf{piezometric level} stands above the ground surface at that point.

\textbf{Special cases: karst and hot/mineral springs.} In limestone regions, dissolution creates caves and conduits; groundwater emerges suddenly as large \textbf{karst springs} (resurgences), often where an underground river meets the surface. \textbf{Hot (thermal) springs} occur where water circulates deeply, is heated (by the geothermal gradient or by cooling magma), and rises along faults; on its way up the water dissolves minerals and gases, giving \textbf{mineral springs}. Around \textbf{Mount Cameroon} and other centres of the Cameroon Volcanic Line, warm and mineral springs are associated with volcanic fractures.

\begin{popfigure}
\begin{center}
\begin{tikzpicture}[scale=0.86, font=\small]
% ---- Depression spring ----
\begin{scope}
\draw[fill=popblueL] (0,0) -- (0,1.1) .. controls (1.2,1.3) and (2.2,1.5) .. (3.4,1.2) -- (3.4,0) -- cycle;
\draw[thick, popink] (0,1.1) .. controls (1.2,1.3) and (2.2,1.5) .. (3.4,1.2);
\draw[dashed, popblue, thick] (0,1.35) .. controls (1.2,1.5) and (2.2,1.62) .. (3.4,1.4);
\draw[->, popblue, thick] (2.6,1.35) -- (2.9,1.2);
\node at (1.7,-0.35) {\textbf{Depression spring}};
\node[popblue] at (0.7,1.7) {\scriptsize water table};
\node[popblue] at (3.0,0.55) {\scriptsize spring};
\end{scope}
% ---- Contact spring ----
\begin{scope}[xshift=4.6cm]
\draw[fill=popblueL] (0,0) -- (0,1.6) -- (3.4,0.9) -- (3.4,0) -- cycle;
\draw[fill=gray!45] (0,0) -- (0,-0.7) -- (3.4,-0.7) -- (3.4,0) -- cycle;
\draw[thick, popink] (0,1.6) -- (3.4,0.9);
\draw[thick, popink] (0,0) -- (3.4,0);
\draw[->, popblue, thick] (1.2,1.0) -- (2.6,0.35);
\node at (1.7,-1.1) {\textbf{Contact spring}};
\node at (1.1,0.75) {\scriptsize aquifer};
\node at (1.7,-0.35) {\scriptsize aquiclude};
\node[popblue] at (3.15,0.5) {\scriptsize spring};
\end{scope}
% ---- Fault spring ----
\begin{scope}[xshift=9.2cm]
\draw[fill=popblueL] (0,0) rectangle (1.5,1.7);
\draw[fill=gray!45] (1.7,0) rectangle (3.4,1.7);
\draw[thick, popink] (0,1.7) -- (3.4,1.7);
\draw[very thick, poporange] (1.5,-0.5) -- (1.75,1.7);
\draw[->, popblue, thick] (0.7,0.3) -- (1.35,0.6);
\draw[->, popblue, thick] (1.55,0.7) -- (1.68,1.55);
\node at (1.7,-0.9) {\textbf{Fault spring}};
\node at (0.75,1.0) {\scriptsize aquifer};
\node at (2.6,1.0) {\scriptsize impermeable};
\node[poporange] at (2.6,0.1) {\scriptsize fault};
\node[popblue] at (1.15,1.95) {\scriptsize spring};
\end{scope}
% ---- Artesian spring ----
\begin{scope}[xshift=0cm, yshift=-4.4cm]
\draw[fill=gray!45] (0,1.5) -- (13.0,1.1) -- (13.0,0.7) -- (0,1.1) -- cycle;
\draw[fill=popblueL] (0,1.1) -- (13.0,0.7) -- (13.0,0.0) -- (0,0.4) -- cycle;
\draw[fill=gray!45] (0,0.4) -- (13.0,0.0) -- (13.0,-0.5) -- (0,-0.1) -- cycle;
\draw[thick, popink] (0,1.5) .. controls (4,1.9) and (8,1.6) .. (13.0,1.1);
\draw[very thick, poporange] (9.6,-0.4) -- (10.3,1.45);
\draw[->, popblue, thick] (2.0,0.75) -- (5.5,0.55);
\draw[->, popblue, thick] (8.6,0.45) -- (9.75,0.7);
\draw[->, popblue, thick] (9.95,0.85) -- (10.2,1.35);
\node at (2.2,1.32) {\scriptsize aquiclude};
\node at (3.2,0.72) {\scriptsize confined aquifer};
\node at (6.5,0.28) {\scriptsize aquiclude};
\node[poporange] at (11.2,0.35) {\scriptsize fault};
\node[popblue] at (10.9,1.6) {\scriptsize artesian spring};
\node at (6.5,-0.85) {\textbf{Artesian spring}};
\end{scope}
\end{tikzpicture}
\end{center}
\legende{Cross-sections of the four main spring types. Blue: aquifer/water; grey: aquiclude; orange: fault; dashed line: water table.}
\end{popfigure}

\begin{attention}[Exam tip (GCE): identifying spring types]
A frequent structured question shows a labelled cross-section and asks you to \textbf{name the type of spring and justify your answer}. Your justification must quote the geometry: e.g.\ ``contact spring, because the aquifer rests on an aquiclude whose outcrop is cut by the slope''. Always mention \emph{both} the aquifer and the structure that forces the water out.
\end{attention}

\courssub{Groundwater and Rivers: Effluent and Influent Streams}

\textbf{Intuition.} In the heart of the dry season in Yaoundé, when no rain has fallen for weeks, some streams still flow. Where does that water come from? Not from the sky — from the ground. The river is being \emph{fed by the aquifer}.

\begin{definition}[Effluent and influent rivers]
An \cle{effluent (gaining) river} is one whose channel lies \textbf{below the water table}, so that groundwater flows \emph{into} the river. An \cle{influent (losing) river} is one whose channel lies \textbf{above the water table}, so that river water seeps \emph{down into} the aquifer.
\end{definition}

The sustained dry-weather flow of a river, supplied almost entirely by groundwater seeping from the banks and bed, is called the \cle{base flow}. Base flow explains why permanent rivers in Cameroon keep flowing through the dry season: the aquifer acts as a natural reservoir that releases water slowly. When the water table drops below the river bed (prolonged drought, over-pumping nearby), an effluent river can become influent, or even dry up completely. Conversely, influent rivers are typical of \textbf{arid regions and of permeable terrain} (sands, karst limestones), where rivers lose water to the ground and may even disappear underground.

\begin{popfigure}
\begin{center}
\begin{tikzpicture}[scale=0.9, font=\small]
% Effluent
\begin{scope}
\draw[fill=popblueL] (0,0) -- (0,1.4) -- (6,1.4) -- (6,0) -- cycle;
\draw[fill=white] (2.2,1.4) .. controls (2.6,0.9) and (3.4,0.9) .. (3.8,1.4) -- cycle;
\draw[thick, popink] (0,1.4) -- (2.2,1.4) .. controls (2.6,0.9) and (3.4,0.9) .. (3.8,1.4) -- (6,1.4);
\draw[dashed, popblue, thick] (0,1.15) -- (2.2,1.15);
\draw[dashed, popblue, thick] (3.8,1.15) -- (6,1.15);
\draw[->, popblue, thick] (1.2,0.7) -- (2.45,0.95);
\draw[->, popblue, thick] (4.8,0.7) -- (3.55,0.95);
\node[popblue] at (0.9,1.32) {\scriptsize water table};
\node at (3,0.35) {\scriptsize river};
\node at (3,-0.4) {\textbf{Effluent (gaining) river}};
\end{scope}
% Influent
\begin{scope}[xshift=8cm]
\draw[fill=popblueL] (0,0) -- (0,0.55) -- (6,0.55) -- (6,0) -- cycle;
\draw[fill=white] (2.2,1.4) .. controls (2.6,0.9) and (3.4,0.9) .. (3.8,1.4) -- cycle;
\draw[thick, popink] (0,1.4) -- (2.2,1.4) .. controls (2.6,0.9) and (3.4,0.9) .. (3.8,1.4) -- (6,1.4);
\draw[dashed, popblue, thick] (0,0.55) -- (6,0.55);
\draw[->, popblue, thick] (2.7,0.85) -- (2.2,0.35);
\draw[->, popblue, thick] (3.3,0.85) -- (3.8,0.35);
\node[popblue] at (5.0,0.75) {\scriptsize water table};
\node at (3,1.05) {\scriptsize river};
\node at (3,-0.4) {\textbf{Influent (losing) river}};
\end{scope}
\end{tikzpicture}
\end{center}
\legende{Left: an effluent river — the water table is higher than the river bed and groundwater feeds the river. Right: an influent river — the river bed lies above the water table and river water recharges the aquifer.}
\end{popfigure}

\courssub{Groundwater and the Ocean: Submarine Discharge and Saltwater Intrusion}

Along the \textbf{Douala–Kribi coastal plain}, aquifers in sands and sedimentary rocks extend beneath the sea. Two phenomena occur at this groundwater–ocean interface.

\textbf{Submarine groundwater discharge.} Where the water table on land stands above sea level, fresh groundwater flows seaward and can emerge \emph{beneath the sea} as submarine springs — an invisible loss (or potential resource) of fresh water.

\textbf{Saltwater intrusion.} Fresh water is less dense than salt water, so in a coastal aquifer fresh groundwater \emph{floats} on denser salt water that has penetrated from the sea. The position of the fresh–salt interface follows from a simple pressure balance. Consider a column of fresh water of total height $(h + z)$, where $h$ is the height of the water table above sea level and $z$ the depth of the interface below sea level. At the interface, the pressure of the fresh column must equal the pressure of the adjacent sea-water column of height $z$:
\[
\rho_f\,(h+z) = \rho_s\, z
\quad\Longrightarrow\quad
z = \frac{\rho_f}{\rho_s - \rho_f}\, h .
\]
With $\rho_f \approx 1.00\ \mathrm{g\,cm^{-3}}$ for fresh water and $\rho_s \approx 1.025\ \mathrm{g\,cm^{-3}}$ for sea water, the ratio $\rho_f/(\rho_s - \rho_f) = 1.00/0.025 = 40$, giving the \textbf{Ghyben–Herzberg relation}:
\[
z \approx 40\,h .
\]
In words: \textbf{for every 1 metre of fresh water standing above sea level, about 40 metres of fresh water extend below sea level}. The danger is that pumping lowers $h$; the interface then rises about 40 times as much, and salt water can reach the well — a serious risk for coastal boreholes around Douala and Kribi.

\begin{exoresolu}[Applying the Ghyben–Herzberg relation]
A coastal well near Kribi taps an unconfined sandy aquifer. The water table stands $1.5\ \mathrm{m}$ above sea level. (a) Estimate the depth of the fresh–salt water interface below sea level. (b) If over-pumping lowers the water table by $0.5\ \mathrm{m}$, what happens to the interface?
\tcblower
\textbf{(a)} Using $z \approx 40\,h$:
\[
z = 40 \times 1.5 = 60\ \mathrm{m\ below\ sea\ level.}
\]
\textbf{(b)} The new height is $h' = 1.5 - 0.5 = 1.0\ \mathrm{m}$, so
\[
z' = 40 \times 1.0 = 40\ \mathrm{m.}
\]
The interface has \textbf{risen by $20\ \mathrm{m}$} for only $0.5\ \mathrm{m}$ of drawdown — illustrating why even modest over-pumping of coastal wells can cause salinisation.
\end{exoresolu}

\courssub{Water Wells: Tapping the Aquifer}

\begin{definition}[Water well]
A \cle{well} is an artificial excavation or borehole sunk into the ground to reach the water table (or a confined aquifer) and extract groundwater for domestic, agricultural or industrial use.
\end{definition}

\textbf{Types of wells.}
\begin{itemize}
\item \textbf{Hand-dug wells:} wide shafts (about $1$–$1.5\ \mathrm{m}$ diameter), usually less than $20\ \mathrm{m}$ deep, common in rural Cameroon; cheap but easily polluted at the surface.
\item \textbf{Drilled boreholes (tube wells):} narrow (about $0.1$–$0.3\ \mathrm{m}$), drilled by machine to tens or hundreds of metres; they can reach deep or confined aquifers and are easier to protect.
\item \textbf{Artesian wells:} boreholes that penetrate a \textbf{confined aquifer}. The water rises in the casing up to the piezometric level; if that level stands \emph{above the ground surface}, the well \textbf{flows naturally} without any pump (a \emph{flowing artesian well}).
\end{itemize}

\textbf{Main parts of a borehole.} The \textbf{casing} is a pipe lining the hole to prevent collapse and block shallow contaminated water; the \textbf{screen} is a slotted section placed opposite the aquifer to let water in; the \textbf{gravel pack} (clean coarse sand/gravel around the screen) filters the incoming water; the \textbf{pump} lifts the water to the surface.

\begin{definition}[Cone of depression]
The \cle{cone of depression} is the funnel-shaped lowering of the water table (or piezometric surface) that develops around a pumped well, because water is removed faster near the well than it can flow in from the aquifer.
\end{definition}

\begin{popfigure}
\begin{center}
\begin{tikzpicture}[scale=1.0, font=\small]
% Aquifer block
\draw[fill=popblueL] (0,0) rectangle (10,3.2);
\draw[thick, popink] (0,3.2) -- (10,3.2);
\node at (1.3,0.4) {\scriptsize aquifer};
% Static water table
\draw[dashed, popblue, thick] (0,2.6) -- (3.6,2.6);
\draw[dashed, popblue, thick] (6.4,2.6) -- (10,2.6);
\node[popblue] at (1.6,2.85) {\scriptsize static level};
% Cone of depression
\draw[popdark, thick] (3.6,2.6) .. controls (4.3,1.6) and (4.6,1.3) .. (4.85,1.2);
\draw[popdark, thick] (6.4,2.6) .. controls (5.7,1.6) and (5.4,1.3) .. (5.15,1.2);
\node[popdark] at (8.3,1.6) {\scriptsize cone of depression};
% Pumping level
\draw[poporange, thick] (4.55,1.2) -- (5.45,1.2);
\node[poporange] at (7.6,1.05) {\scriptsize pumping level};
% Well
\draw[fill=gray!60] (4.85,-0.4) rectangle (4.95,3.6);
\draw[fill=gray!60] (5.05,-0.4) rectangle (5.15,3.6);
\draw[fill=white] (4.95,-0.4) rectangle (5.05,3.6);
% Screen
\draw[popgreen, very thick] (4.85,0.0) -- (4.85,0.9);
\draw[popgreen, very thick] (5.15,0.0) -- (5.15,0.9);
\node[popgreen] at (6.6,0.35) {\scriptsize screen + gravel pack};
% Pump and arrows
\draw[fill=popink] (4.85,1.6) rectangle (5.15,1.95);
\node at (6.3,2.2) {\scriptsize pump};
\draw[->, popblue, very thick] (5.0,2.0) -- (5.0,3.4);
\draw[->, popblue, thick] (3.6,0.5) -- (4.7,0.5);
\draw[->, popblue, thick] (6.4,0.5) -- (5.3,0.5);
% Drawdown arrow
\draw[<->, poppink, thick] (4.35,1.2) -- (4.35,2.6);
\node[poppink] at (3.0,1.9) {\scriptsize drawdown $s$};
\node at (4.0,3.5) {\scriptsize casing};
\end{tikzpicture}
\end{center}
\legende{A pumped borehole: casing, screen, static level, pumping level, drawdown $s$ and the cone of depression in the surrounding water table.}
\end{popfigure}

\courssub{Well Hydraulics in Simple Terms}

When a well is pumped, the water level inside falls from the \textbf{static level} (rest level) to a lower \textbf{pumping level}. The difference is the \cle{drawdown} $s$:
\[
s = \text{static level} - \text{pumping level} \quad (\text{in metres}).
\]
If the well yields a discharge $Q$ (e.g.\ in $\mathrm{m^3\,h^{-1}}$), we define the \cle{specific capacity}:
\[
\text{Specific capacity} = \frac{Q}{s},
\]
the yield obtained per metre of drawdown. A high specific capacity indicates a productive, permeable aquifer; a low value warns of a poor aquifer or a badly designed well.

\begin{exemplebox}[Drawdown and specific capacity]
A borehole has a static level at $12\ \mathrm{m}$ below ground. Pumped at $Q = 20\ \mathrm{m^3\,h^{-1}}$, the level stabilises at $17\ \mathrm{m}$. Then $s = 17 - 12 = 5\ \mathrm{m}$ and the specific capacity is $Q/s = 20/5 = 4\ \mathrm{m^3\,h^{-1}\,m^{-1}}$ — a reasonably productive well.
\end{exemplebox}

\textbf{Risks of over-pumping.} If abstraction exceeds recharge:
\begin{itemize}
\item the water table falls and shallow wells and springs \textbf{dry up};
\item in compressible sediments (clays, silts), loss of water pressure can cause \textbf{land subsidence};
\item in coastal areas, the fresh–salt interface rises and \textbf{saltwater intrusion} renders wells unusable.
\end{itemize}

\begin{methode}[Siting a village well: geological and sanitary checklist]
\begin{enumerate}
\item \textbf{Locate the aquifer:} use a geological map, existing wells, spring lines and vegetation indicators to estimate the \textbf{depth to the water table} — shallow enough to be reached, deep enough not to dry up seasonally.
\item \textbf{Check aquifer thickness and permeability:} a thick, coarse-grained or fractured aquifer gives a sustainable yield; thin or clayey layers give little water.
\item \textbf{Confirm recharge:} the catchment upslope must receive enough rain to replenish what is pumped.
\item \textbf{Respect sanitary distances:} site the well \textbf{upslope and far from pollution sources} — pit latrines, refuse dumps, animal pens, fuel stores (a rule of thumb is at least $30\ \mathrm{m}$, and more in fractured rock).
\item \textbf{Plan protection:} raised apron, cover, and casing sealed at the top to stop surface water running back in.
\end{enumerate}
\end{methode}

\begin{attention}[Common mistake: ``deep means safe'']
A deep well is \textbf{not} automatically a safe well. Depth alone does not guarantee quality: if the casing is cracked, the apron broken, or a latrine sits upslope, contaminated water can still enter — and deep wells can even draw in \emph{naturally} mineralised or saline water. Safety depends on \textbf{protection and siting}, not only on depth.
\end{attention}

\begin{exercice}[Spring identification and well siting]
A cross-section shows a sandstone aquifer resting on shale. On the valley side, water emerges exactly along the sandstone–shale boundary. (a) Name the type of spring and justify. (b) A village wants a well $200\ \mathrm{m}$ from this spring: state two geological and two sanitary criteria to check.
\end{exercice}

\begin{corrige}[Exercise 1]
\textbf{(a)} Contact spring: the permeable sandstone overlies the impermeable shale, groundwater flows along the contact, and emerges where the topography cuts that contact. \textbf{(b)} Geological: confirm the water table is shallow enough and the sandstone thick/permeable enough for a sustainable yield. Sanitary: site the well upslope of latrines/refuse dumps and at a safe distance (at least about $30\ \mathrm{m}$), with a sealed casing and apron.
\end{corrige}

\begin{aretenir}[Key points]
\begin{itemize}
\item A \textbf{spring} is a natural outflow of groundwater where the water table or a structure (contact, fault) meets the surface; know the four types: depression, contact, fault, artesian.
\item \textbf{Effluent rivers} are fed by groundwater (base flow); \textbf{influent rivers} recharge the aquifer.
\item In coastal aquifers, fresh water floats on salt water; by \textbf{Ghyben–Herzberg}, $z \approx 40h$ — pumping can trigger \textbf{saltwater intrusion}.
\item A \textbf{well} taps the aquifer; pumping creates a \textbf{cone of depression} and a \textbf{drawdown} $s$; specific capacity $= Q/s$.
\item Over-pumping causes drying up, subsidence and salinisation; good \textbf{siting and protection} matter more than depth alone.
\end{itemize}
\end{aretenir}

\coursec{Rock Properties Controlling Groundwater Flow; Aquifers and Aquicludes}

In the previous section we followed groundwater to its natural outlets: springs, rivers and, ultimately, the ocean, and we saw how a well allows us to tap this hidden resource. But one question was left hanging: why does water accumulate abundantly in the sands of the Douala coastal plain, while a few kilometres away a borehole drilled into massive granite may yield almost nothing? The answer lies in two fundamental rock properties — \cle{porosity} and \cle{permeability} — and in the way permeable and impermeable layers are stacked underground. Mastering these notions is essential: they explain where groundwater lives, how it moves, and how it can be exploited.

\courssub{Porosity: the Storage Capacity of Rocks}

\textbf{Intuition.}\par
Take a bucket and fill it with gravel. The gravel looks ``full'', yet you can still pour a surprising amount of water into it before it overflows. The water occupies the empty spaces — the \emph{voids} or \emph{pores} — between the grains. Every rock, from loose beach sand to the hardest granite, contains such voids, and the proportion of voids in the rock determines how much water it can \emph{store}.

\begin{definition}[Porosity]
The \cle{porosity} of a rock, denoted $n$, is the ratio of the volume of voids (pores, cracks, cavities) to the total volume of the rock, usually expressed as a percentage:
\[ n = \frac{V_{\text{voids}}}{V_{\text{total}}} \times 100\ \% \]
where $V_{\text{voids}}$ is the volume of empty space and $V_{\text{total}}$ the bulk volume of the rock sample (grains + voids).
\end{definition}

\textbf{Primary and secondary porosity.}\par
Two great families of pores must be distinguished:
\begin{itemize}
\item \cle{Primary porosity} (or \emph{intergranular} porosity) is created at the same time as the rock itself. It consists of the spaces left between sedimentary grains during deposition. It is typical of gravels, sands, sandstones and clays.
\item \cle{Secondary porosity} develops \emph{after} the rock has formed, through fracturing (joints and faults in granite, basalt or quartzite) or through dissolution (\emph{karst} porosity: caves, fissures and channels widened by the dissolution of limestone by acidified rainwater). Secondary porosity is the only reason why otherwise compact rocks such as granite or massive limestone can hold water at all.
\end{itemize}

\textbf{Factors controlling porosity.}\par
\begin{itemize}
\item \textbf{Sorting (grain size distribution).} Well-sorted sediments (grains of similar size) have high porosity. Poorly sorted sediments have low porosity because small grains fill the voids between large ones. Note that grain size \emph{alone} does not control porosity: a well-sorted clay and a well-sorted gravel can have similar porosities.
\item \textbf{Packing and grain shape.} Loosely packed, rounded grains leave more void space than tightly packed or angular, interlocking grains.
\item \textbf{Cementation.} In sandstones, calcite or silica cement precipitates in the pores and glues the grains together, reducing porosity — sometimes drastically.
\item \textbf{Fracturing and dissolution.} These processes \emph{increase} porosity in compact rocks by creating secondary voids.
\end{itemize}

\begin{center}
\begin{tabularx}{\linewidth}{|l|X|l|}
\hline
\rowcolor{popblueL} \textbf{Rock / sediment} & \textbf{Dominant porosity type} & \textbf{Typical porosity (\%)} \\
\hline
Gravel (well sorted) & Primary, intergranular & 25 -- 40 \\
\hline
Sand & Primary, intergranular & 25 -- 50 \\
\hline
Sandstone & Primary, reduced by cement & 5 -- 30 \\
\hline
Clay & Primary, very small pores & 35 -- 60 \\
\hline
Limestone & Secondary (karst) + primary & 1 -- 20 \\
\hline
Granite (unfractured / fractured) & Secondary (fractures) & $<1$ / up to 5 \\
\hline
\end{tabularx}
\end{center}

\begin{attention}[Porous does not mean permeable]
The most frequent mistake at the GCE is to equate \cle{porosity} with \cle{permeability}. Clay is the classic counter-example: it has a \emph{very high} porosity (up to 60\ \%) yet water passes through it extremely slowly because its pores are microscopic and poorly connected. Always separate the two ideas: porosity measures \emph{storage}; permeability measures \emph{transmission}.
\end{attention}

\courssub{Permeability: the Transmission Capacity of Rocks}

\textbf{Intuition.}\par
Pour water onto a heap of gravel: it disappears instantly. Pour the same water onto a slab of compact clay: it spreads out and stays at the surface. Both materials contain pores, but in the gravel the pores are large and interconnected, so water flows through; in the clay the pores are so small that molecular attraction holds the water back. The property that differs is permeability.

\begin{definition}[Permeability]
\cle{Permeability} is the capacity of a rock to let water (or any fluid) flow through its interconnected voids under a pressure gradient. A rock is said to be \emph{permeable} if water circulates through it easily (gravel, sand, fractured basalt, karstified limestone) and \emph{impermeable} if it does not (clay, shale, unfractured granite).
\end{definition}

Permeability depends on:
\begin{itemize}
\item the \textbf{size of the pores} — large pores (gravel, sand) transmit water far better than micropores (clay);
\item the \textbf{interconnection of the pores} — isolated cavities (as in some vesicular lavas) store water but do not transmit it;
\item the \textbf{presence of fractures or karst channels} — these can give a compact rock an excellent secondary permeability.
\end{itemize}

\begin{propriete}[Darcy's law (qualitative form)]
The discharge $Q$ of groundwater through a permeable bed is proportional to the cross-sectional area $A$ and to the \emph{hydraulic gradient} $i$ (the slope of the water table or piezometric surface, $i = \dfrac{\Delta h}{L}$):
\[ Q = K \times A \times i \]
where $K$ is the \cle{hydraulic conductivity}, a measure of the rock's permeability. Gravel and clean sand have high $K$; clay has a $K$ millions of times smaller. (Only the qualitative reading — flow is faster through more permeable rocks and down steeper hydraulic gradients — is required at this level.)
\end{propriete}

\begin{popfigure}
\begin{center}
\begin{tikzpicture}[scale=0.92]
% axes
\draw[->,thick,popink] (0,0) -- (11,0);
\draw[->,thick,popink] (0,0) -- (0,5.6);
\node[rotate=90,font=\small,popink] at (-0.7,2.8) {Relative value (qualitative scale)};
% bars: pairs porosity (blue) / permeability (orange)
\fill[popblue]  (0.3,0) rectangle (0.9,3.5);
\fill[poporange](1.0,0) rectangle (1.6,5.0);
\node[font=\small,align=center] at (0.95,-0.35) {Gravel};
\fill[popblue]  (2.3,0) rectangle (2.9,4.0);
\fill[poporange](3.0,0) rectangle (3.6,4.4);
\node[font=\small,align=center] at (2.95,-0.35) {Sand};
\fill[popblue]  (4.3,0) rectangle (4.9,2.0);
\fill[poporange](5.0,0) rectangle (5.6,2.6);
\node[font=\small,align=center] at (4.95,-0.35) {Sandstone};
\fill[popblue]  (6.3,0) rectangle (6.9,5.2);
\fill[poporange](7.0,0) rectangle (7.6,0.4);
\node[font=\small,align=center] at (6.95,-0.35) {Clay};
\fill[popblue]  (8.3,0) rectangle (8.9,0.3);
\fill[poporange](9.0,0) rectangle (9.6,0.3);
\node[font=\small,align=center] at (8.95,-0.35) {Granite};
% legend
\fill[popblue] (3.6,5.9) rectangle (4.0,6.2); \node[right,font=\small] at (4.05,6.05) {Porosity};
\fill[poporange] (5.9,5.9) rectangle (6.3,6.2); \node[right,font=\small] at (6.35,6.05) {Permeability};
\end{tikzpicture}
\end{center}
\legende{Comparison of porosity and permeability for common rocks. Note clay: the highest porosity but almost no permeability; note gravel and sand: both properties high — ideal aquifer materials.}
\end{popfigure}

\textbf{Specific yield and specific retention.}\par
If porosity tells us how much water a rock \emph{contains}, it does not tell us how much we can \emph{extract}. When a saturated rock is drained by gravity (for example by pumping a well), part of the water remains trapped in the pores by capillary forces and molecular attraction.
\begin{itemize}
\item \cle{Specific yield} is the volume of water that drains freely by gravity, expressed as a percentage of the total rock volume. It is high in coarse materials (gravel, sand: typically 20--35\ \%).
\item \cle{Specific retention} is the volume of water retained against gravity, as a percentage of the total volume. It is high in fine materials (clay), where capillary forces are strong.
\end{itemize}
The two quantities add up to the porosity:
\[ n = \text{specific yield} + \text{specific retention} \]
This is why a clay holding 50\ \% of water may deliver almost nothing to a well, while a sand of porosity 30\ \% delivers most of its stored water.

\begin{exemplebox}[Worked example: how much water can a sand deposit deliver?]
A saturated sand layer has a volume of $2\,000\ \mathrm{m^3}$, a porosity of 30\ \% and a specific retention of 5\ \%. Total stored water: $V_{\text{water}} = 0.30 \times 2000 = 600\ \mathrm{m^3}$. Water extractable by gravity (specific yield $= 30 - 5 = 25\ \%$): $V_{\text{yield}} = 0.25 \times 2000 = 500\ \mathrm{m^3}$. The remaining $100\ \mathrm{m^3}$ stays bound in the pores.
\end{exemplebox}

\courssub{Aquifers and Aquicludes}

\begin{definition}[Aquifer and aquiclude]
An \cle{aquifer} is a permeable rock formation that \emph{stores} and \emph{transmits} groundwater in quantities large enough to be exploited (wells, springs). Typical aquifers: sands and gravels, porous sandstones, fractured volcanic rocks, karstified limestones.\par
An \cle{aquiclude} is a formation that may \emph{store} water but does \emph{not transmit} it in appreciable amounts (clay, shale). An \cle{aquifuge} neither stores nor transmits water (massive unfractured granite). An \cle{aquitard} is an intermediate case: a semi-permeable layer (silty clay, fractured shale) that transmits water only very slowly, retarding flow between aquifers.
\end{definition}

\textbf{Types of aquifers.}\par
\begin{itemize}
\item \textbf{Unconfined (water-table) aquifer.} The aquifer is bounded above only by unsaturated ground; its upper surface is the \emph{water table}, free to rise and fall with recharge and pumping. The water in a well dug into it stands at the water-table level.
\item \textbf{Confined aquifer.} The aquifer is sandwiched between two impermeable layers (aquicludes). The water it contains is under pressure; its level in a borehole rises above the top of the aquifer, defining an imaginary \emph{piezometric surface}.
\item \textbf{Perched aquifer.} A small lens of impermeable material (e.g. a clay bed) lying \emph{above} the main water table traps a local, limited body of groundwater. Perched aquifers are common in volcanic terrains where clayey weathered horizons alternate with lava flows; they often dry up seasonally.
\item \textbf{Artesian aquifer.} A special case of the confined aquifer in which the pressure is sufficient for water to rise \emph{above the ground surface} in a well, sometimes flowing freely.
\end{itemize}

\begin{propriete}[Conditions for artesian flow]
Artesian flow (water rising spontaneously at the surface) requires four conditions, all visible on a geological cross-section:
\begin{enumerate}
\item a \textbf{permeable aquifer layer} (e.g. sandstone) \textbf{confined between two aquicludes} (e.g. clay or shale);
\item layers that are \textbf{inclined} (a synclinal or monoclinal structure), forming an \emph{artesian basin};
\item a \textbf{recharge area} where the aquifer crops out at \emph{higher altitude}, so that rainwater enters it;
\item a \textbf{piezometric level above the ground surface} at the well site, so that hydrostatic pressure pushes the water up and out.
\end{enumerate}
(The explanation — pressure transmitted through the confined aquifer from the high recharge zone — is examinable as a labelled cross-section.)
\end{propriete}

\begin{popfigure}
\begin{center}
\begin{tikzpicture}[scale=1.0]
% ground surface
\draw[thick,popink] (0,3.5) -- (2,3.5) -- (3.2,4.6) -- (4.4,3.5) -- (11,3.5);
% upper aquiclude
\fill[popgold!60] (0,2.6) -- (2,2.6) -- (3.2,3.7) -- (4.4,2.6) -- (11,2.6)
 -- (11,2.2) -- (4.4,2.2) -- (3.2,3.3) -- (2,2.2) -- (0,2.2) -- cycle;
% confined aquifer
\fill[popblue!55] (0,2.2) -- (2,2.2) -- (3.2,3.3) -- (4.4,2.2) -- (11,2.2)
 -- (11,1.2) -- (4.4,1.2) -- (3.2,2.3) -- (2,1.2) -- (0,1.2) -- cycle;
% lower aquiclude
\fill[popgold!60] (0,1.2) -- (2,1.2) -- (3.2,2.3) -- (4.4,1.2) -- (11,1.2)
 -- (11,0.4) -- (0,0.4) -- cycle;
% labels of layers
\node[font=\small,popdark] at (9.6,2.4) {Aquiclude (clay)};
\node[font=\small,white] at (9.6,1.7) {Confined aquifer (sandstone)};
\node[font=\small,popdark] at (9.6,0.8) {Aquiclude (shale)};
% recharge area
\draw[->,thick,popblue] (3.2,5.3) -- (3.2,4.0);
\node[font=\small,align=center,popblue] at (3.2,5.7) {Recharge area\\(aquifer outcrop)};
\draw[->,thick,popblue] (2.9,4.4) -- (3.15,3.6);
% piezometric surface
\draw[dashed,thick,popgreen] (3.2,4.6) -- (11,4.6);
\node[font=\small,popgreen,align=center] at (8.0,4.9) {Piezometric surface};
% flowing well
\draw[thick,popink] (7.0,3.5) -- (7.0,1.6);
\draw[thick,popink] (7.15,3.5) -- (7.15,1.6);
\draw[->,very thick,popblue] (7.07,1.7) -- (7.07,4.3);
\draw[->,thick,popblue] (7.07,4.3) .. controls (7.4,4.5) .. (7.8,4.1);
\node[font=\small,align=center,popink] at (7.2,3.0) [right,xshift=2pt]{Flowing\\artesian well};
% non-flowing well
\draw[thick,popink] (9.9,3.5) -- (9.9,1.6);
\draw[thick,popink] (10.05,3.5) -- (10.05,1.6);
\draw[popblue,thick] (9.9,4.0) -- (10.05,4.0);
\node[font=\small,align=center,popink] at (10.6,3.0) {Water rises but\\does not flow};
% rain
\node[font=\small,popmuted] at (1.2,4.1) {Rain};
\draw[->,popmuted] (1.2,3.9) -- (1.2,3.55);
\end{tikzpicture}
\end{center}
\legende{Cross-section of an artesian basin: an inclined sandstone aquifer confined between two aquicludes is recharged at high altitude. Where the piezometric surface lies above the ground, wells flow naturally.}
\end{popfigure}

\begin{methode}[Identifying the aquifer type on a geological cross-section]
\begin{enumerate}
\item \textbf{Locate the permeable layer} (sand, sandstone, fractured lava, limestone) — this is the potential aquifer.
\item \textbf{Look for confining layers above it.} If the aquifer is overlain directly by soil/unsaturated rock, it is \emph{unconfined}: mark the water table. If it is capped by clay or shale, it is \emph{confined}: draw the piezometric surface (the level water reaches in boreholes), which lies \emph{above} the top of the aquifer.
\item \textbf{Check the geometry.} If the confined aquifer is inclined and crops out at higher altitude (recharge zone), compare the piezometric surface with the ground level at the well: above ground $\Rightarrow$ \emph{flowing artesian} well; between aquifer top and ground $\Rightarrow$ \emph{non-flowing artesian} (pumping needed).
\item \textbf{Check for perched conditions:} a small clay lens above the regional water table with a local saturated zone on top indicates a \emph{perched aquifer}.
\end{enumerate}
\end{methode}

\courssub{Cameroonian Examples}

\begin{savaistu}[Aquifers of Cameroon]
Cameroon offers a textbook illustration of every aquifer type:
\begin{itemize}
\item The \textbf{sandy coastal aquifers of the Douala basin}: thick, poorly consolidated sands and gravels of the coastal sedimentary plain form highly porous and permeable unconfined to semi-confined aquifers that supply much of the city's water — but their shallow water table makes them vulnerable to pollution and, near the coast, to salt-water intrusion.
\item The \textbf{volcanic aquifers of the West highlands} (Bambouto, Manengouba, Cameroon Volcanic Line): water circulates through fractures and scoriaceous (vesicular) horizons between basaltic lava flows. Alternation of lavas and impermeable weathered clays produces perched aquifers and the famous high-altitude springs of the Western Region.
\item The \textbf{basement fractured aquifers of the Centre and Adamawa} (Yaoundé region, Adamawa plateau): the ancient crystalline rocks (granites, gneisses) of the Congo craton margin have almost no primary porosity; groundwater is stored and transmitted only in the \emph{weathered zone} and in fracture networks — a typical secondary-porosity aquifer, where borehole success depends on intersecting fractures.
\end{itemize}
\end{savaistu}

\begin{exoresolu}[Worked exercise: choosing a borehole site]
\textbf{Statement.} A village on the Adamawa plateau needs a borehole. The local geology shows (a) massive unfractured granite, (b) a fractured granite zone, and (c) a clay layer. State, with reasons, which horizon should be targeted and what rock properties justify the choice.
\tcblower
\textbf{Solution.} The target is horizon (b), the fractured granite. Massive granite (a) is an aquifuge: porosity below 1\ \%, no interconnected voids, so it neither stores nor transmits water. Clay (c) is an aquiclude: high porosity but extremely low permeability, so it stores water without releasing it to a well. Fractured granite (b) possesses secondary (fracture) porosity and permeability: the interconnected cracks store groundwater and allow it to flow towards the borehole. It therefore behaves as a usable fractured basement aquifer.
\end{exoresolu}

\begin{aretenir}[Key points]
\begin{itemize}
\item Porosity $n = \dfrac{V_{\text{voids}}}{V_{\text{total}}} \times 100\ \%$: the storage property. Primary (intergranular) vs secondary (fracture, karst).
\item Permeability: the transmission property; requires large, interconnected voids. Clay = high porosity, low permeability. Flow rate increases with permeability and with the hydraulic gradient (Darcy's law, $Q = K A i$).
\item Specific yield + specific retention = porosity; only the specific yield can be pumped.
\item Aquifer = stores \emph{and} transmits; aquiclude = stores but does not transmit; aquifuge = neither; aquitard = transmits very slowly.
\item Aquifer types: unconfined (water table), confined (piezometric surface), perched, artesian. Artesian flow needs a confined, inclined aquifer, high recharge and a piezometric level above ground.
\end{itemize}
\end{aretenir}

\begin{exercice}[Porosity, yield and aquifer identification]
\begin{enumerate}
\item A sandstone sample of total volume $500\ \mathrm{cm^3}$ contains $75\ \mathrm{cm^3}$ of voids. Calculate its porosity.
\item The same sandstone has a specific retention of 4\ \%. What volume of water (per $\mathrm{m^3}$ of rock) can be extracted by pumping?
\item On a cross-section, a sand layer lies between two shale beds and dips from a high plateau towards a plain. A borehole in the plain encounters water that rises to the surface. Name the aquifer type and list the four conditions it illustrates.
\end{enumerate}
\end{exercice}

\begin{corrige}[Exercise 1]
\begin{enumerate}
\item $n = \dfrac{75}{500}\times 100 = 15\ \%$.
\item Specific yield $= 15 - 4 = 11\ \%$, i.e. $0.11\ \mathrm{m^3}$ of extractable water per $\mathrm{m^3}$ of rock ($110\ \mathrm{L\,m^{-3}}$).
\item A flowing artesian aquifer. Conditions: (i) permeable aquifer confined between two aquicludes (shales); (ii) inclined layers; (iii) recharge at the high plateau outcrop; (iv) piezometric level above the ground surface at the borehole.
\end{enumerate}
\end{corrige}

Having understood how rocks store and transmit water, and how aquifers are structured, we can now ask the final practical question of this chapter: is this groundwater \emph{fit to drink}? That is the subject of the next section, devoted to water quality and groundwater contamination.

\coursec{Water Quality and Groundwater Contamination}

In the previous sections we established that the \cle{porosity} and \cle{permeability} of rocks determine where groundwater is stored and how fast it moves, and that the geological succession of \cle{aquifers} and \cle{aquicludes} controls whether a water body is unconfined or confined. But knowing that an aquifer \emph{contains} water is not enough: the essential question for any community is whether that water is \emph{safe to drink}. A well sunk into the Douala sedimentary basin or into the weathered basement of the Adamawa Plateau may yield abundant water that is nevertheless dangerous. This final section therefore examines the \cle{quality} of groundwater, the ways in which it becomes \cle{contaminated}, and the measures used to protect it — a topic of enormous practical importance in Cameroon, where millions of people depend on wells, boreholes and springs.

\courssub{Water Quality Criteria}

\begin{definition}[Potable water]
\cle{Potable water} is water that meets accepted physical, chemical and bacteriological standards and is therefore safe for human consumption without risk of disease or long-term poisoning.
\end{definition}

Water quality is judged under three headings.

\textbf{1. Physical criteria.} These are the properties we can sense directly:
\begin{itemize}
\item \cle{Colour}: potable water should be colourless. A brown or yellow tint suggests dissolved organic matter (humic acids) or iron oxides.
\item \cle{Turbidity}: cloudiness caused by suspended clay, silt or organic particles (measured in NTU — Nephelometric Turbidity Units). Turbid water has often entered the well recently from the surface, carrying pollution with it; it also shields pathogens from disinfection.
\item \cle{Taste and odour}: good water is tasteless and odourless. A rotten-egg smell indicates hydrogen sulphide (\ce{H2S}) from sulphate reduction; a salty taste indicates dissolved chlorides or sodium.
\item \cle{Temperature}: groundwater normally has a fairly constant temperature (close to mean annual air temperature). A sudden change may indicate direct hydraulic connection with surface water.
\end{itemize}

\textbf{2. Chemical criteria.} Groundwater dissolves minerals from the rocks through which it flows, so it always contains dissolved substances. The key parameters and typical WHO guideline values are:

\begin{center}
\begin{tabularx}{\linewidth}{|l|X|c|}\hline
\rowcolor{popblueL}
\textbf{Parameter} & \textbf{Significance / Origin} & \textbf{WHO Guideline Value} \\ \hline
\cle{pH} & Acidic water corrodes pipes, mobilises toxic metals (Pb, Cd); alkaline water tastes bitter. & 6.5 -- 8.5 \\ \hline
\cle{Hardness} & Caused by \ce{Ca^{2+}} and \ce{Mg^{2+}} (carbonate = temporary; non-carbonate = permanent). Not hazardous but causes scaling, excessive soap use. & No health-based limit (taste \textasciitilde 200--500 mg/L \ce{CaCO3}) \\ \hline
\cle{Total Dissolved Solids (TDS)} & Sum of all dissolved ions. High TDS gives unpleasant taste, laxative effect. & \textless 600 mg/L (palatable); \textgreater 1200 mg/L unacceptable \\ \hline
\cle{Nitrates (\ce{NO3-})} & Mainly from latrines, animal waste, fertilisers. Oxidised form of nitrogen. Causes methaemoglobinaemia (``blue baby'' syndrome) in infants. & 50 mg/L (as \ce{NO3-}) \\ \hline
\cle{Iron (\ce{Fe})} & Common in basement/volcanic aquifers (reducing conditions). Stains laundry, metallic taste, promotes iron bacteria. & 0.3 mg/L (aesthetic) \\ \hline
\cle{Fluoride (\ce{F-})} & Natural from fluorite, apatite, volcanic rocks. Beneficial \textless 0.5 mg/L (dental health); excess causes dental/skeletal fluorosis. & 1.5 mg/L \\ \hline
\cle{Chloride (\ce{Cl-})} & Indicator of salinity / salt-water intrusion. Taste threshold \textasciitilde 250 mg/L. & 250 mg/L (taste) \\ \hline
\cle{Arsenic (\ce{As})} & Geogenic in some volcanic/sedimentary basins. Highly toxic, carcinogenic. & 10 $\mu$ g/L \\ \hline
\end{tabularx}
\end{center}

\textbf{3. Bacteriological criteria.} This is the most critical group for health. Water is tested for \cle{indicator organisms} rather than every pathogen:
\begin{itemize}
\item \cle{Total coliforms}: bacteria found in soil, vegetation, and faeces. Their presence indicates possible contamination.
\item \cle{Faecal coliforms / \emph{Escherichia coli}}: specific to warm-blooded animal intestines. \emph{Definitive proof of recent faecal contamination}.
\item \cle{Pathogens}: bacteria (\emph{Vibrio cholerae}, \emph{Salmonella typhi}), viruses (hepatitis A, rotavirus), protozoa (\emph{Giardia}, \emph{Cryptosporidium}).
\end{itemize}
Potable water should contain \emph{zero} faecal coliforms (\emph{E. coli}) in a $100\ \mathrm{cm^3}$ sample.

\begin{aretenir}[Potability classes (WHO framework)]
\begin{itemize}
\item \cle{Potable}: clear, colourless, odourless, acceptable taste and pH, all chemical parameters below guideline values, zero faecal coliforms/100 mL — safe to drink directly.
\item \cle{Treatable}: fails one or more criteria but can be made safe by simple household treatment (boiling, filtration, chlorination, SODIS).
\item \cle{Polluted}: contains dangerous levels of pathogens or toxic chemicals (arsenic, nitrate \textgreater 100 mg/L, hydrocarbons); unfit for consumption without major treatment, sometimes unusable.
\end{itemize}
\end{aretenir}

\begin{attention}[Clear water is not necessarily safe]
The most dangerous mistake is to judge water by its appearance. Pathogens, nitrates, arsenic, fluoride and dissolved hydrocarbons are \emph{invisible}: a crystal-clear well can transmit cholera or slowly poison a child with nitrates. Only laboratory analysis or proper sanitary inspection can establish potability. Conversely, slightly coloured water (e.g. from harmless iron or humic acids) may be perfectly safe after settling.
\end{attention}

\courssub{Sources of Groundwater Contamination}

Contamination may be \emph{anthropogenic} (caused by human activities) or \emph{natural (geogenic)}.

\textbf{Anthropogenic sources} (point and non-point):
\begin{itemize}
\item \cle{Pit latrines and septic tanks}: the most widespread source in Cameroonian towns and villages. Liquid effluent percolates downward carrying bacteria, viruses, nitrates, phosphates and pharmaceutical residues toward the water table.
\item \cle{Refuse dumps (landfills)}: rainwater leaches rotting waste, producing a toxic liquid (\cle{leachate}) rich in organic matter (high BOD/COD), ammonia, nitrates, heavy metals (Pb, Cd, Hg) and persistent organic pollutants.
\item \cle{Agricultural fertilisers and pesticides}: nitrates and phosphates not taken up by crops leach into the soil; persistent pesticides (organochlorines, atrazine) and herbicides can reach aquifers, especially in intensively farmed areas such as the banana/tea plantations of the South-West (Mungo, Tiko) and West (Dschang, Bafoussam) Regions.
\item \cle{Industrial effluents}: untreated liquid waste from breweries, tanneries, abattoirs, textile factories, and agro-industries discharged onto the ground or into streams (e.g. Wouri estuary industrial zone).
\item \cle{Fuel stations and garages}: leaking underground storage tanks (LUSTs) and spilled hydrocarbons (BTEX: benzene, toluene, ethylbenzene, xylene); even few litres of petrol render thousands of cubic metres of water undrinkable (benzene limit 10 $\mu$ g/L).
\item \cle{Mining and quarrying}: artisanal gold mining (batan) in the East and North exposes sulphide rocks to oxidation, generating Acid Mine Drainage (AMD) — acidic, metal-laden runoff.
\end{itemize}

\textbf{Natural (geogenic) contamination} occurs without any human activity:
\begin{itemize}
\item \cle{Salinization of coastal aquifers}: around Douala, Kribi, Limbe and Tiko, excessive pumping lowers the freshwater hydraulic head, allowing denser sea water to migrate inland into the aquifer (\emph{salt-water intrusion}). The Ghyben-Herzberg principle states that for every metre of freshwater drawdown above sea level, the salt-water interface rises approximately 40 metres below sea level.
\item \cle{Excess iron and manganese}: in reducing conditions of deep basement aquifers (North, Adamawa, East) or volcanic rocks, prolonged contact dissolves \ce{Fe^{2+}} and \ce{Mn^{2+}}.
\item \cle{Excess fluoride}: in some volcanic terrains (Cameroon Volcanic Line) and granitic basement (Adamawa, North), fluorine-bearing minerals (fluorite, apatite, micas) raise \ce{F-} above 1.5 mg/L, causing dental/skeletal fluorosis.
\item \cle{Arsenic}: locally elevated in volcanic rocks and geothermal areas.
\item \cle{Radon gas (\ce{^{222}Rn})}: radioactive decay product of uranium in granites/volcanics; dissolves in groundwater, inhalation risk during domestic use (showering).
\end{itemize}

\courssub{Pathways of Contamination and Aquifer Vulnerability}

\begin{definition}[Aquifer vulnerability]
The \cle{vulnerability} of an aquifer is the ease with which a pollutant introduced at the ground surface can reach the water table. It depends on the depth to the water table, the permeability and thickness of the overlying materials (unsaturated zone), and the presence or absence of a protective aquiclude.
\end{definition}

Pollutants reach groundwater by several routes:
\begin{enumerate}
\item \cle{Direct infiltration}: rain dissolves pollutants at the surface (dumps, agricultural fields, spills) and carries them straight down through permeable soil and rock fractures.
\item \cle{Leakage through the unsaturated zone (vadose zone)}: contaminants from latrines, septic tanks, and leaking sewers percolate slowly downward. The soil filters suspended solids and retains some bacteria/viruses by adsorption; soil microbes degrade biodegradable organics. However, a shallow water table (\textless 10 m), coarse sand, lateritic gravel, or fissured basalt offers little protection — travel times can be days to weeks.
\item \cle{River--aquifer exchange}: where a polluted river is \emph{influent} (losing stream, feeding the aquifer), contaminants pass from the river into the groundwater body. Common in dry season when river stage drops below water table.
\item \cle{Preferential flow paths}: macropores (root holes, animal burrows, desiccation cracks), fractures, faults, and unsealed boreholes act as rapid conduits bypassing the soil matrix filtration.
\item \cle{Well-head pollution}: dirty buckets/ropes, surface run-off entering around a poorly sealed well apron, animals drinking at the well, or flooding of the well head introduce contamination directly into the water column.
\end{enumerate}

The \cle{protective cover} — the soil layer and especially any overlying \cle{aquiclude} (clay, shale, lateritic duricrust) — is the aquifer's natural shield. A \cle{confined aquifer} beneath a thick, continuous clay layer is well protected: pollutants cannot easily cross the impermeable barrier. An \cle{unconfined aquifer} with a shallow water table beneath sandy, permeable soil or fissured rock is highly \cle{vulnerable}: a latrine dug a few metres away can contaminate it within days or weeks.

\begin{popfigure}
\begin{tikzpicture}[scale=0.9, font=\small, every node/.style={text width=2.2cm, align=center}]
% Ground surface
\draw[thick, popink] (0,0) -- (14,0);
% Unsaturated zone (left: permeable; right: clay layer)
\fill[popgold!25] (0,0) rectangle (7,-2.8);
\fill[popgold!25] (7,-0) rectangle (14,-1.2);
\fill[popmuted!60] (7,-1.2) rectangle (14,-2.2); % Clay aquiclude
\fill[popgold!25] (7,-2.2) rectangle (14,-2.8);
% Saturated zone
\fill[popblueL] (0,-2.8) rectangle (7,-5.5);
\fill[popblueL] (7,-2.8) rectangle (14,-5.5);
% Water table line (left, unconfined)
\draw[thick, popblue, dashed] (0,-2.8) -- (7,-2.8);
\node[popblue, anchor=west, text width=3cm] at (0.2,-2.5) {Water table (unconfined)};
% Confined piezometric level
\draw[thick, popblue, dashed] (7,-1.5) -- (14,-1.5);
\node[popblue, anchor=west, text width=3cm] at (7.2,-1.2) {Piezometric level (confined)};
% Labels
\node[popink, align=center] at (3.5,-4.2){Unconfined aquifer\\(vulnerable)};
\node[popink, align=center] at (10.5,-4.2){Confined aquifer\\(protected)};
\node[popmuted, anchor=west] at (7.2,-1.7) {Clay aquiclude};
% Latrine (left)
\draw[thick, popdark, fill=popdark!10] (1,0) rectangle (2,-2);
\node[popdark, anchor=north] at (1.5,0.3) {Pit latrine};
% Refuse dump (left)
\draw[thick, poporange, fill=poporange!30] (3.5,0) .. controls (3.8,0.5) and (4.7,0.5) .. (5,0) -- cycle;
\node[poporange, anchor=south] at (4.25,0.3) {Refuse dump};
% Contamination arrows (vulnerable side)
\draw[-latex, very thick, poppink] (1.5,-2.1) -- (1.5,-3.5);
\draw[-latex, very thick, poppink] (4.25,-0.1) -- (4.25,-3.5);
% Shallow well (vulnerable)
\draw[thick, popink] (5.8,0) -- (5.8,-3.8) (6.4,0) -- (6.4,-3.8);
\draw[thick, popink] (5.8,0) -- (5.8,0.3) -- (6.4,0.3) -- (6.4,0);
\node[popink, anchor=south] at (6.1,0.5) {Shallow well};
\draw[-latex, very thick, poppink] (4.8,-3.5) -- (5.7,-3.8);
% Protected borehole (right)
\draw[thick, popgreen] (10.5,0) -- (10.5,-5.2) (11.1,0) -- (11.1,-5.2);
\draw[thick, popgreen, fill=popgreen!20] (10.3,0) rectangle (11.3,0.4);
\node[popgreen, anchor=south] at (10.8,0.6) {Sealed borehole};
% Latrine on right (blocked)
\draw[thick, popdark, fill=popdark!10] (12.5,0) rectangle (13.5,-1);
\node[popdark, anchor=north] at (13,0.3) {Latrine};
% Blocked arrow at aquiclude
\draw[-latex, very thick, poppink] (13,-1.1) -- (13,-2.15);
\draw[very thick, poppink, decorate, decoration={markings, mark=at position 0.5 with {\node[poppink, rotate=90] {\texttimes};}}] (12.5,-2.2) -- (13.5,-2.2);
\node[poppink, anchor=south, text width=2.5cm] at (13,0.5) {Plume blocked by aquiclude};
\end{tikzpicture}
\legende{Contamination pathways. Left: a latrine and refuse dump overlie a shallow unconfined aquifer in permeable laterite; pollutants infiltrate freely to the water table and migrate to a nearby shallow well. Right: a thick clay aquiclude blocks downward leakage, protecting the confined aquifer tapped by a sealed borehole; the latrine plume is stopped at the aquiclude.}
\end{popfigure}

\textbf{Self-purification (natural attenuation) and its limits.} As water percolates through the unsaturated zone, several processes reduce contamination:
\begin{itemize}
\item \cle{Filtration}: suspended particles, protozoan cysts (\emph{Giardia}, \emph{Cryptosporidium}) and some bacteria are physically strained out.
\item \cle{Adsorption}: viruses, heavy metals and hydrophobic organics adhere to clay minerals, iron oxides and organic matter.
\item \cle{Biodegradation}: aerobic and anaerobic soil microbes decompose organic matter, ammonium (nitrification), and some pesticides.
\item \cle{Dilution and dispersion}: mixing with cleaner groundwater reduces concentrations.
\end{itemize}
This capacity is \emph{real but limited}: it fails when the soil is thin, coarse-grained, fissured, or waterlogged; when the water table is shallow (\textless 5--10 m); or when the pollutant load exceeds the soil's assimilative capacity. Crucially, \emph{dissolved substances such as nitrates, chlorides, arsenic, fluoride and many hydrocarbons are not removed by filtration} and persist in the groundwater.

\courssub{Comparing Protected and Unprotected Water Points}

Field surveys in Cameroonian urban and peri-urban areas (Yaoundé, Douala, Bafoussam, Maroua) consistently show that wells/boreholes with a proper apron, cover, casing seal and adequate distance from pollution sources have far lower nitrate and coliform levels than unprotected traditional wells. The chart below illustrates this typical contrast (illustrative data based on hydrogeological literature).

\begin{popfigure}
\begin{tikzpicture}
\begin{axis}[
ybar, bar width=1.0cm,
width=12cm, height=6.5cm,
ylabel={Concentration / Count (relative units)},
symbolic x coords={Nitrates (mg/L), Faecal coliforms (CFU/100mL)},
xtick=data,
ymin=0, ymax=100,
legend style={at={(0.5,1.05)}, anchor=south, legend columns=2, font=\small},
nodes near coords, every node near coord/.append style={font=\small},
tick label style={font=\small},
]
\addplot[fill=popgreen!70, draw=popgreen] coordinates {(Nitrates (mg/L),15) (Faecal coliforms (CFU/100mL),5)};
\addplot[fill=poppink!70, draw=poppink] coordinates {(Nitrates (mg/L),75) (Faecal coliforms (CFU/100mL),90)};
\legend{Protected borehole / well, Unprotected traditional well}
\end{axis}
\end{tikzpicture}
\legende{Illustrative comparison of nitrate and faecal coliform levels in protected versus unprotected water points. Protection (sanitary seal, apron, cover, distance \textgreater 30 m from latrines) reduces contamination dramatically.}
\end{popfigure}

\courssub{Protection, Treatment and Sustainable Management}

\begin{methode}[Sanitary inspection of a well or borehole (GCE fieldwork standard)]
A systematic inspection checks the following critical points:
\begin{enumerate}
\item \textbf{Apron}: Is there a watertight concrete apron (\textgreater 1.5 m radius) around the well-head, sloping outward (\textgreater 1\%) so that spilled water and surface run-off drain away?
\item \textbf{Cover and seal}: Is the well covered with a tight-fitting lid? Is the casing sealed above ground level (minimum 0.3--0.5 m) so that surface water, animals, dust and dirty buckets/ropes cannot enter?
\item \textbf{Drainage}: Is waste water (from drawing water) channelled away into a soakaway or drain at a safe distance (\textgreater 5 m), not ponding near the well?
\item \textbf{Distance from pollution sources}: As a rule of thumb, a water point should be at least \SI{30}{m} from any pit latrine, septic tank, refuse dump, animal pen, or fuel storage, and preferably \emph{upslope} (hydraulically upgradient) from them.
\item \textbf{Water-table depth and protective cover}: The deeper the water table and the thicker the low-permeability soil/rock cover, the safer the source.
\item \textbf{Condition of surroundings}: Is the area fenced? Is there evidence of open defecation, solid waste dumping, or chemical storage nearby?
\end{enumerate}
\end{methode}

\textbf{Sanitary protection zones.} Around a production well or borehole, concentric protection zones are established (based on travel-time criteria):
\begin{itemize}
\item \cle{Immediate protection zone (Zone I)}: Radius 10--20 m. Fenced, concreted, strictly forbidden access. Prevents direct well-head pollution and vandalism.
\item \cle{Inner protection zone (Zone II)}: Defined by the 50-day (or 100-day) travel time of groundwater to the well. No latrines, dumps, livestock, agriculture, fuel tanks, or industries allowed. Size depends on hydraulic gradient and permeability (typically 100--500 m radius).
\item \cle{Outer protection zone (Zone III)}: The entire recharge area (catchment) of the aquifer. Land-use planning controls potentially polluting activities (zoning regulations, best agricultural practices, effluent treatment enforcement).
\end{itemize}

\textbf{Water treatment basics (household / community level).} Where potability is doubtful, simple treatment makes water safe:
\begin{itemize}
\item \cle{Boiling}: A rolling boil for at least 1 minute (3 minutes above 2000 m altitude) kills virtually all pathogens (bacteria, viruses, protozoa). Does \emph{not} remove dissolved chemicals (nitrates, fluoride, arsenic, heavy metals) and concentrates them slightly.
\item \cle{Chlorination}: Careful dosing with sodium hypochlorite (bleach, typically 2.6\% active chlorine) or chlorine tablets (e.g. Aquatabs) to achieve 0.2--0.5 mg/L free residual chlorine after 30 min contact. Effective against bacteria and viruses; less effective against \emph{Cryptosporidium}. Requires clear water (low turbidity).
\item \cle{Filtration}: Slow sand filters (community) or household ceramic/candle filters remove turbidity, protozoa and most bacteria. Do not remove viruses or dissolved chemicals unless combined with activated carbon or specific media.
\item \cle{SODIS (Solar Water Disinfection)}: Filling transparent PET bottles and exposing to full sunlight for 6 hours (UV-A + heat). Low-cost, effective for bacteria/viruses in clear water. Promoted in rural Cameroon.
\end{itemize}

\textbf{Sustainable exploitation.} Pumping must not exceed the aquifer's long-term average natural recharge; over-pumping (mining groundwater) lowers the water table, dries up shallow wells and springs, reduces baseflow to rivers, and in coastal areas (Douala, Kribi, Limbe, Tiko) draws sea water into the aquifer, causing irreversible salinization. Managed Aquifer Recharge (MAR) — directing stormwater or treated wastewater into infiltration basins — is a growing strategy to augment recharge.

\begin{exoresolu}[Assessing contamination risk in a Cameroonian town]
\textbf{Statement.} In a quarter of Bafoussam (West Region), a family digs a shallow traditional well reaching the water table at \SI{6}{m} depth in weathered, permeable lateritic soil (saprolite). A pit latrine stands \SI{8}{m} away, slightly upslope (hydraulically upgradient). A neighbour proposes instead a borehole drilled to \SI{40}{m}, tapping a fractured basalt aquifer overlain by a \SI{15}{m} thick clay layer (weathered basalt/alterite), with a sealed concrete apron and sanitary seal. (a) Assess the contamination risk of each water point. (b) Advise the family.

\tcblower
\textbf{Solution.} (a) \emph{Shallow traditional well}: 
\begin{itemize}
\item \textbf{Hydrogeological setting}: Unconfined aquifer in permeable laterite (high hydraulic conductivity). Water table shallow (\SI{6}{m}) $\Rightarrow$ thin unsaturated zone, short travel time.
\item \textbf{Protective cover}: None (no aquiclude). Laterite offers some filtration but limited adsorption for viruses/nitrates.
\item \textbf{Pollution source}: Pit latrine at \SI{8}{m} (well inside the \SI{30}{m} minimum), \emph{upslope} $\Rightarrow$ effluent plume flows directly toward the well.
\item \textbf{Pathways}: Direct infiltration of latrine effluent through unsaturated zone; well-head pollution likely (no apron/cover mentioned).
\item \textbf{Risk}: \textbf{Very high}. Water almost certainly contains faecal coliforms, high nitrates, and possibly pathogens even if clear.
\end{itemize}
\emph{Deep borehole}:
\begin{itemize}
\item \textbf{Hydrogeological setting}: Confined (or semi-confined) aquifer in fractured basalt. Piezometric level likely above the aquiclude.
\item \textbf{Protective cover}: \SI{15}{m} thick clay aquiclude (low permeability) $\Rightarrow$ excellent barrier to vertical leakage. Travel time through clay: years to decades.
\item \textbf{Construction}: Sealed apron and sanitary seal prevent well-head pollution.
\item \textbf{Risk}: \textbf{Low}. Main residual risk is contamination via the borehole annulus if grouting is defective, or cross-contamination during drilling.
\end{itemize}
(b) \emph{Advice}: 
\begin{enumerate}
\item Immediately stop drinking untreated water from the shallow well. Treat by boiling or chlorination if no alternative exists.
\item Invest in the protected borehole; verify driller's log confirms clay thickness and proper grouting/sealing.
\item Enforce a sanitary protection zone around the borehole: no new latrines/dumps within \SI{30}{m} (ideally \SI{50}{m} upslope), relocate existing latrine if possible.
\item Conduct baseline water quality analysis (coliforms, nitrates, iron, fluoride, pH) before commissioning and repeat quarterly.
\item Educate household on safe water handling (clean containers, no dipping).
\end{enumerate}
\end{exoresolu}

\begin{savaistu}[Groundwater and public health in Cameroon]
Cholera outbreaks recorded in Douala, Yaoundé, Maroua and Bamenda have repeatedly been epidemiologically linked to contaminated shallow wells and springs, especially in densely built informal quarters where pit latrines stand \textless 10 m from water points. Conversely, the rapid growth of professional borehole drilling — from Yaoundé's fractured basement to the Far North's sedimentary aquifers — demonstrates how understanding aquifer protection literally saves lives. In coastal Douala, the \emph{Société des Eaux du Cameroun} and hydrogeologists must limit pumping rates and monitor chloride trends to slow the advance of the salt-water wedge into the Wouri and Mungo aquifers. In the Lake Nyos region, groundwater monitoring also tracks dissolved \ce{CO2} hazards.
\end{savaistu}

\begin{exercice}[Essay practice (GCE Advanced Level style)]
A rapidly growing Cameroonian town obtains its water from two sources: (i) shallow hand-dug wells in a sandy unconfined aquifer, and (ii) deep boreholes in a confined sandstone aquifer protected by a thick clay layer. 
\begin{enumerate}
\item Describe and explain the differences in contamination risk between the two sources.
\item Identify the likely pollutants in each case and their pathways to the water points.
\item Propose a comprehensive management plan to secure the town's water supply.
\end{enumerate}
\end{exercice}

\begin{corrige}[Exercise 1]
\emph{Expected structured answer:}
\begin{enumerate}
\item \textbf{Risk contrast}: 
\begin{itemize}
\item Source (i) Unconfined sandy aquifer: High permeability, shallow water table, no protective aquiclude $\Rightarrow$ \textbf{High vulnerability}. Pollutants infiltrate rapidly (days/weeks) through unsaturated zone. Well-head pollution adds risk.
\item Source (ii) Confined sandstone aquifer: Clay aquiclude blocks vertical leakage $\Rightarrow$ \textbf{Low vulnerability}. Long travel times (years) allow attenuation. Main risk is well-head pollution or faulty borehole seals.
\end{itemize}
\item \textbf{Likely pollutants \& pathways}:
\begin{itemize}
\item Source (i): Faecal coliforms/\emph{E. coli} (latrines $\rightarrow$ infiltration), Nitrates (latrines, dumps, fertilisers $\rightarrow$ leaching), Hydrocarbons (garages $\rightarrow$ infiltration), Salinity (if coastal $\rightarrow$ lateral intrusion). Pathways: Direct infiltration, preferential flow, well-head entry.
\item Source (ii): Natural geogenic (Fluoride, Iron, Arsenic from rock-water interaction). Anthropogenic only if borehole seal fails (surface run-off $\rightarrow$ annulus) or deep injection wells exist. Pathway: Borehole conduit.
\end{itemize}
\item \textbf{Management plan}:
\begin{itemize}
\item \emph{Immediate}: Sanitary inspection of all wells; enforce \SI{30}{m} lateral distance / upslope rule; construct concrete aprons, covers, drainage for all wells.
\item \emph{Protection zones}: Delineate Zone I (fenced), Zone II (50-day travel time, land-use restrictions), Zone III (recharge area planning) for borehole fields.
\item \emph{Pollution source control}: Relocate/upgrade latrines to sealed septic tanks/soakaways; line refuse dumps; regulate agro-chemicals; enforce effluent treatment for industries.
\item \emph{Monitoring}: Quarterly sampling for coliforms, nitrates, chlorides, pH, conductivity at representative points.
\item \emph{Treatment}: Chlorination at borehole heads; household boiling/filtration for shallow well users during transition.
\item \emph{Sustainable yield}: Pumping tests to determine safe yield; install water meters; limit abstraction to \textless 70\% recharge; monitor water levels and chloride in coastal boreholes to detect salt-water intrusion early.
\item \emph{Institutional}: Water safety plans (WHO WSP approach); community water committees; capacity building for local technicians.
\end{itemize}
\end{enumerate}
A clear, structured answer linking \emph{aquifer type $\rightarrow$ vulnerability $\rightarrow$ human activities $\rightarrow$ specific pollutants $\rightarrow$ pathways $\rightarrow$ targeted protection measures} earns full credit.
\end{corrige}

\begin{attention}[Exam tip (GCE Advanced Level)]
Expect an essay or structured question asking you to \emph{link aquifer type, human activities and contamination risk} in a named Cameroonian context (e.g. shallow wells in a Douala/Bafoussam/Maroua quarter, coastal salinization near Kribi/Limbe, fluoride in Adamawa volcanic zone). Structure your answer: 
\begin{enumerate}
\item Define potability (physical, chemical, bacteriological) and aquifer vulnerability.
\item Describe the hydrogeological setting (cross-section sketch: unconfined vs confined, protective cover, water table/piezometric surface, flow direction).
\item Identify pollution sources (point/non-point, anthropogenic/natural) and pathways (infiltration, leakage, river exchange, well-head).
\item Propose protection (sanitary zones, \SI{30}{m} rule, sealing, land-use control), treatment (boiling, chlorination, filtration), and sustainable management (recharge balance, monitoring).
\end{enumerate}
A labelled cross-section sketch \emph{always} strengthens the answer and is often explicitly required.
\end{attention}

\begin{aretenir}[Key points of the section]
\begin{itemize}
\item Water quality is judged by \cle{physical} (colour, turbidity, taste, temperature), \cle{chemical} (pH, hardness, TDS, nitrates, iron, fluoride, chloride, arsenic) and \cle{bacteriological} (coliforms, \emph{E. coli}, pathogens) criteria; potable water satisfies all three (WHO guidelines).
\item Clear water is not necessarily safe: pathogens, nitrates, arsenic, fluoride are invisible.
\item Major contamination sources: \emph{Anthropogenic} — pit latrines, refuse dumps (leachate), agriculture (nitrates, pesticides), industry, fuel leaks, mining (AMD); \emph{Natural} — coastal salt-water intrusion, excess Fe/Mn/F/As/Rn from rock-water interaction.
\item Contamination reaches aquifers by: direct infiltration, leakage through unsaturated zone, river-aquifer exchange, preferential flow paths, well-head pollution.
\item \cle{Vulnerability} depends on depth to water table and protective cover; unconfined shallow aquifers are highly vulnerable, confined aquifers beneath aquicludes are protected.
\item Protection combines: sanitary protection zones (I, II, III), \SI{30}{m} minimum distance rule, sealed aprons/covers, source control, water quality monitoring, household treatment (boiling, chlorination, filtration, SODIS), and sustainable pumping $\le$ recharge.
\end{itemize}
\end{aretenir}

\newpage

\coursec{Integration activity}

\courssub{Aims of the activity}

This integration activity closes the chapter on hydrogeology by asking you to act as a consulting hydrogeologist. Working in groups, you will use everything studied in this chapter to solve a concrete, realistic problem: \textbf{where should a borehole be drilled to supply a quarter of Bafoussam with safe drinking water, and how should it be protected?}

By the end of the activity you should be able to:

\begin{itemize}
\item \textbf{Read and interpret} a simplified geological cross-section (topography, rock units, water table);
\item \textbf{Identify} the \cle{aquifer} and the \cle{aquiclude}, and state the \cle{type of aquifer} (unconfined, confined or perched) with justification;
\item \textbf{Predict} the direction of \cle{groundwater flow} from the shape of the water table and the topography;
\item \textbf{Apply} the concepts of \cle{porosity}, \cle{permeability} and \cle{specific yield} to compare candidate drilling sites;
\item \textbf{Predict} the formation of a \cle{cone of depression} around a pumped well and explain one risk of \cle{over-pumping};
\item \textbf{Assess contamination risks} (latrines, refuse dump, stream) and apply \cle{sanitary distance} criteria;
\item \textbf{Propose} realistic protection measures for a borehole in a Cameroonian urban setting;
\item \textbf{Communicate} a reasoned technical choice on a poster and defend it in a class debate.
\end{itemize}

\begin{aretenir}[Skills mobilised]
Hydrologic cycle (recharge); sources of groundwater; infiltration and percolation; rock properties (porosity, permeability); aquifer types; water table and piezometric surface; flow direction; wells, drawdown and cone of depression; springs and streams; groundwater contamination and protection.
\end{aretenir}

\courssub{Situation}

\textbf{Bafoussam, capital of the West Region of Cameroon, lies on the Bamenda--Bafoussam volcanic plateau at about \SI{1400}{m} altitude.} The quarter of \emph{Ndiangdam}, on the southern edge of the town, has grown rapidly. Its 8\,000 inhabitants depend on a few hand-dug wells that dry up between December and February, and on a small stream that is visibly polluted. The quarter's water committee has raised funds to drill \textbf{one borehole equipped with a hand pump}, and has hired your team of young hydrogeologists to recommend the best site.

A geological reconnaissance survey has produced the simplified cross-section below, running roughly north--south across the quarter.

\begin{popfigure}
\begin{tikzpicture}[x=1cm,y=0.55cm,font=\small]
% ground surface
\draw[thick,popink] (0,6) -- (2,6.4) -- (4,6.2) -- (6,6.8) -- (8,6.5) -- (10,5.6) -- (12,4.8) -- (14,4.4);
% basalt fill
\fill[popblueL] (0,0) -- (14,0) -- (14,4.4) -- (12,4.8) -- (10,5.6) -- (8,6.5) -- (6,6.8) -- (4,6.2) -- (2,6.4) -- (0,6) -- cycle;
% clay layer
\fill[popgold!60] (0,0) -- (14,0) -- (14,1.2) -- (10,1.4) -- (6,1.6) -- (0,1.8) -- cycle;
% fractures in basalt
\draw[popdark] (1.5,1.9) -- (1.2,5.9); \draw[popdark] (3.5,1.8) -- (3.8,6.2);
\draw[popdark] (6.5,1.7) -- (6.2,6.7); \draw[popdark] (9,1.5) -- (9.3,5.5);
\draw[popdark] (11.5,1.4) -- (11.2,4.7); \draw[popdark] (13,1.3) -- (13.3,4.5);
% water table (dashed)
\draw[dashed,thick,popblue] (0,4.6) -- (2,4.9) -- (4,4.7) -- (6,5.2) -- (8,4.9) -- (10,4.1) -- (12,3.4) -- (14,3.1);
\node[popblue,anchor=west] at (11.2,3.9) {\small water table};
% stream
\draw[popblue,thick] (13.4,4.35) -- (13.8,4.2) -- (14,4.4);
\node[popblue,align=center] at (13.3,5.3) {\small stream\\ \small (Mifi tributary)};
% latrines
\draw[popink,fill=white] (2.6,6.35) rectangle (3.4,7.15);
\node[align=center] at (3,7.8) {\small latrines};
% refuse dump
\fill[popmuted] (7.4,6.55) -- (8.6,6.55) -- (8.3,7.3) -- (7.7,7.3) -- cycle;
\node[align=center] at (8,7.9) {\small refuse dump};
% candidate sites
\draw[thick,poporange] (2.9,6.4) -- (2.9,0.4);
\node[poporange,anchor=south] at (2.9,0.4) {\textbf{A}};
\draw[thick,poporange] (7.9,6.5) -- (7.9,0.4);
\node[poporange,anchor=south] at (7.9,0.4) {\textbf{B}};
\draw[thick,poporange] (12.4,4.7) -- (12.4,0.4);
\node[poporange,anchor=south] at (12.4,0.4) {\textbf{C}};
% labels
\node[align=center] at (4.5,3.2) {\small fractured basalt\\ \small (weathered at top)};
\node at (4.5,0.8) {\small massive clay};
% scale
\draw[<->] (0,-1.2) -- (14,-1.2) node[midway,below]{\small 1400 m (horizontal, not to scale vertically)};
\node[anchor=east] at (-0.2,6) {\small N};
\node[anchor=west] at (14.2,4.4) {\small S};
\end{tikzpicture}
\legende{Simplified geological cross-section of the Ndiangdam quarter (Bafoussam). Vertical scale exaggerated. A, B and C are the three candidate borehole sites.}
\end{popfigure}

\textbf{Data collected by the survey team}

\begin{center}
\begin{tabularx}{\linewidth}{|l|X|}
\hline
\rowcolor{popblueL} \textbf{Feature} & \textbf{Description} \\
\hline
Climate & Tropical highland; rainy season March--October (about \SI{1800}{mm/year}); marked dry season November--February. \\
\hline
Soil cover & Thin (0.5--2 m) lateritic soil, moderately permeable. \\
\hline
Fractured basalt & 15--40 m thick; low intergranular porosity but \textbf{well fractured and weathered near the top}, giving good secondary permeability. \\
\hline
Massive clay & At least 30 m thick (never crossed by any well); very low permeability; no borehole in the region has found water below it. \\
\hline
Water table & Follows the topography in a subdued form; depth below surface: about 8 m at A, 7 m at B, 5 m at C; lies \textbf{entirely within the basalt}. \\
\hline
Stream & Small tributary of the River Mifi at the southern (low) end; flows all year but is polluted downstream of the quarter. \\
\hline
Latrines & Cluster of pit latrines near site A (about 15 m away). \\
\hline
Refuse dump & Uncontrolled household dump on the hilltop, about 20 m from site B. \\
\hline
Distances & A to B: 500 m; B to C: 600 m; latrines to B: 500 m; dump to C: 600 m. \\
\hline
Demand & About \SI{40}{m^3/day} for the quarter (hand pump, moderate pumping rate). \\
\hline
\end{tabularx}
\end{center}

The water committee reminds you that the borehole must (i) \textbf{strike enough water} to supply the quarter even in the dry season, (ii) deliver \textbf{safe water}, and (iii) be \textbf{affordable} (drilling cost increases strongly with depth).

\courssub{Tasks}

Work in groups of 4--6. Answer the questions in order; each one prepares the next. Prepare a poster showing your cross-section annotations, your chosen site and your arguments.

\begin{enumerate}
\item \textbf{Recharge.} Using the hydrologic cycle, explain how the groundwater stored in the basalt of Ndiangdam is recharged. Why is the dry season (November--February) the critical period for the quarter's wells?
\item \textbf{Aquifer and aquiclude.} Identify the aquifer and the aquiclude on the cross-section. Justify each choice using the rock properties (porosity, permeability) of fractured basalt and of massive clay.
\item \textbf{Aquifer type.} Is this aquifer unconfined, confined or perched? Justify your answer by referring to the position of the water table and the absence or presence of an impermeable cover.
\item \textbf{Flow direction.} Reproduce the cross-section and draw arrows showing the direction of groundwater flow. In which general direction does the water move, and towards which natural discharge point? Which law or principle (relation between flow and the slope of the water table) are you applying?
\item \textbf{Choosing the site.} Complete a comparison table for sites A, B and C using the following criteria: depth to water (drilling cost), expected yield (thickness of saturated basalt), distance from pollution sources (latrines, dump), and position relative to groundwater flow (upstream or downstream of the pollution sources). Conclude with your chosen site.
\item \textbf{Cone of depression.} Sketch the water table around your chosen borehole during pumping. Name the feature formed, define \emph{drawdown}, and explain why the pump must not be placed below the water table's lowest dry-season level.
\item \textbf{Over-pumping.} Predict and explain \textbf{one} risk if the community later replaces the hand pump by a powerful motorised pump (think of neighbouring hand-dug wells and of the stream).
\item \textbf{Protection.} Propose \textbf{three} concrete protection measures for the borehole (think of the immediate surroundings of the well head, of the pollution sources, and of the community's habits).
\item \textbf{Poster and debate.} Present your reasoned choice in 5 minutes. Be ready to defend your site against the choices of the other groups: a good hydrogeologist must justify \emph{and} anticipate objections.
\end{enumerate}

\begin{attention}[Common traps]
\begin{itemize}
\item Do not confuse \cle{porosity} (storage) with \cle{permeability} (transmission): clay is porous but almost impermeable.
\item The water table is \emph{not} flat: it is a subdued replica of the topography. Groundwater flows from high to low hydraulic head, i.e.\ roughly from hill to valley, \emph{not} necessarily in the direction of the nearest well.
\item A site close to a pollution source may still be safe if it lies \emph{upstream} of it relative to groundwater flow --- but the reverse is dangerous.
\item ``Deeper'' is not automatically ``better'': below the clay there is no water, and drilling cost rises with depth.
\end{itemize}
\end{attention}

\courssub{Model answer}

\textbf{1. Recharge.} The groundwater of Ndiangdam is meteoric in origin: it comes from the abundant rainfall of the rainy season (March--October). Part of the rain runs off to the stream, part evaporates or is transpired by vegetation, and part \cle{infiltrates} through the thin lateritic soil, then \cle{percolates} down through the weathered and fractured basalt until it reaches the zone of saturation, raising the water table. During the dry season there is almost no recharge, while the population continues to pump: the water table falls, the shallow hand-dug wells (which barely reach the saturated zone) dry up, and only a properly designed borehole reaching well below the dry-season water table can keep supplying the quarter.

\textbf{2. Aquifer and aquiclude.} The \cle{aquifer} is the \textbf{fractured basalt}: although its intergranular (primary) porosity is low, its fractures and weathered upper zone give it a high \cle{secondary porosity} and a good \cle{permeability}, so it both stores and transmits water --- exactly as along the Cameroon Volcanic Line, where fractured lavas are the main aquifers. The \cle{aquiclude} is the \textbf{massive clay}: it has a high porosity but extremely low permeability, so it neither transmits water nor yields it to wells; it forms the impermeable floor of the aquifer.

\textbf{3. Aquifer type.} The aquifer is \cle{unconfined} (a water-table aquifer): its upper surface is the water table, which is in direct contact with the unsaturated zone above and receives recharge directly from rainfall infiltrating through the soil. There is no impermeable layer above the basalt, so it cannot be confined; and since the saturated zone rests on a continuous clay floor and is connected across the whole quarter, it is not a perched aquifer either.

\textbf{4. Flow direction.} The water table slopes gently from north to south, mirroring the topography. Groundwater therefore flows \textbf{from the northern hills towards the southern valley}, perpendicular to the water-table contours and down the hydraulic gradient, and discharges naturally into the Mifi tributary (the stream is the \cle{discharge zone}; this is why it flows all year --- it is fed by the water table in the dry season). This applies the principle that groundwater moves from points of high hydraulic head to points of low hydraulic head (Darcy's law: flow is proportional to the hydraulic gradient).

\textbf{5. Choosing the site.}

\begin{center}
\begin{tabularx}{\linewidth}{|l|X|X|X|}
\hline
\rowcolor{popblueL} \textbf{Criterion} & \textbf{Site A} & \textbf{Site B} & \textbf{Site C} \\
\hline
Depth to water / cost & 8 m: deep, costly & 7 m: moderate & 5 m: cheapest \\
\hline
Saturated thickness / yield & Good & Good & Good (valley) \\
\hline
Pollution source nearby & Latrines 15 m away & Refuse dump 20 m away & None nearby \\
\hline
Position relative to flow & Upstream of B and C, but \textbf{downstream of nothing and next to latrines} & \textbf{Downstream of the dump and of the latrines}: receives their leachate & \textbf{Downstream of the whole quarter}: receives all leachates \\
\hline
Verdict & Rejected: sanitary distance not respected & Rejected: polluted inflow & Rejected: polluted inflow \\
\hline
\end{tabularx}
\end{center}

Strict application of the criteria shows that \textbf{none of the three sites is perfect}, and the group must argue. The most defensible professional choice is \textbf{site A, provided the latrines are relocated}: it lies \emph{upstream} of the dump and of the general flow through the quarter, has a good saturated thickness, and its only defect (the latrines 15 m away) can be corrected by moving them at least 50 m away and downhill of the borehole --- far cheaper than treating contaminated water. Sites B and C cannot be made safe: no affordable measure can stop leachate from the dump (B) or from the whole quarter (C) migrating with the flow towards the well. Choosing C because it is ``cheapest and closest to the stream'' is the classic beginner's error: the cheapest borehole is the one that does not have to be abandoned.

\textbf{6. Cone of depression.} When the pump works, water is removed faster than it flows in, so the water table sinks around the borehole into a funnel-shaped surface called the \cle{cone of depression}; the vertical lowering at the well is the \cle{drawdown}. The sketch shows the dashed water table dipping towards the borehole on both sides. The pump intake must sit below the lowest dry-season water table, with a safety margin, otherwise the well will pump air during drawdown in February.

\textbf{7. Over-pumping.} If a powerful motorised pump is installed, the cone of depression will deepen and widen until it reaches the neighbouring hand-dug wells, \textbf{drying them up} (well interference), and may reverse the hydraulic gradient near the valley so that the stream stops being fed by the aquifer and instead \textbf{polluted stream water is drawn into the borehole} (induced recharge). Either effect would be a social and sanitary disaster.

\textbf{8. Protection measures.} Three examples among many acceptable answers:
\begin{itemize}
\item Build a \textbf{concrete apron and drainage channel} around the well head, fence it, and ban animals, washing and waste disposal within at least 15--30 m (sanitary perimeter);
\item \textbf{Relocate the pit latrines} to at least 50 m from the borehole and downhill of it, and close or move the refuse dump away from the recharge zone;
\item Set up a \textbf{water committee} that monitors the water level, limits pumping to the quarter's real needs (about \SI{40}{m^3/day}), and has the water analysed (bacteriological and nitrate tests) at least once a year.
\end{itemize}

\begin{aretenir}[What the examiner expects]
A correct aquifer/aquiclude identification with \emph{properties} as justification; flow arrows down the water-table slope towards the stream; a site choice argued with \emph{flow direction} (upstream/downstream), not just distance; a labelled cone of depression; and protection measures that are concrete and local. Rigorously argued reasoning is worth more than the ``right'' letter.
\end{aretenir}

\newpage

\coursec{Methods and worked examples}

\courssub{Methods and Worked Examples}

This section presents the essential quantitative and qualitative techniques required for the GCE Advanced Level Hydrogeology paper. Each method is broken down into clear, examinable steps, followed by a fully worked GCE-style question. Master these reflexes: they are the difference between a pass and a distinction.

\begin{methode}[Water Budget Calculation and Effective Rainfall Estimation]
\textbf{Objective:} Quantify the components of the hydrologic cycle at a catchment scale to estimate groundwater recharge ($R$).
\textbf{Key Formula:} $P = ET + R_o + R + \Delta S$ (where $P$ = Precipitation, $ET$ = Evapotranspiration, $R_o$ = Surface Runoff, $R$ = Recharge/Infiltration, $\Delta S$ = Change in Soil Moisture Storage). For long-term annual averages, $\Delta S \approx 0$.
\textbf{Steps:}
\begin{enumerate}
    \item \textbf{List knowns:} Mean annual $P$ (mm/yr), $ET$ (mm/yr) -- use Thornthwaite or Penman-Monteith if data given, $R_o$ (mm/yr) from stream gauging.
    \item \textbf{Apply the balance:} $R = P - ET - R_o$ (assuming steady state).
    \item \textbf{Calculate Recharge Coefficient:} $k_r = R / P \times 100\%$.
    \item \textbf{Interpret geology:} High $k_r$ on fractured basalts (Mount Cameroon flanks) or sandstone (Mamfe Basin); low $k_r$ on clay/shale (Logbaba) or lateritic crusts.
    \item \textbf{Check units:} Convert all to mm/yr or m$^3$/yr (multiply by catchment area $A$ in m$^2$: Volume = $R(\text{m}) \times A$).
\end{enumerate}
\textbf{Examiner's Reflex:} If $\Delta S$ is not negligible (e.g., drought year), state the assumption explicitly. Never confuse $ET$ (actual) with $PET$ (potential).
\end{methode}

\begin{exoresolu}[Catchment Water Balance on the Slopes of Mount Cameroon]
A catchment on the south-eastern flank of Mount Cameroon (area $A = 45\ \mathrm{km^2}$) receives a mean annual rainfall $P = 3200\ \mathrm{mm}$. The mean annual potential evapotranspiration ($PET$) is $1450\ \mathrm{mm}$. The actual evapotranspiration ($ET$) is estimated at $85\%$ of $PET$. The gauged stream at the outlet shows a mean annual discharge $Q = 1.85\ \mathrm{m^3\,s^{-1}}$. Assume negligible change in storage ($\Delta S = 0$).
\begin{enumerate}
    \item Calculate the actual evapotranspiration $ET$ in mm/yr.
    \item Calculate the surface runoff $R_o$ in mm/yr.
    \item Determine the annual groundwater recharge $R$ (in mm/yr and m$^3$/yr).
    \item Compute the recharge coefficient $k_r$. Comment on the value in the context of the volcanic geology.
\end{enumerate}
\tcblower
\textbf{Detailed Solution}
\begin{enumerate}
    \item \textbf{Actual Evapotranspiration ($ET$):}
    \[ ET = 0.85 \times PET = 0.85 \times 1450 = \mathbf{1232.5\ \mathrm{mm/yr}} \]
    \item \textbf{Surface Runoff ($R_o$):}
    Convert discharge $Q$ to depth over catchment area $A$.
    \[ A = 45\ \mathrm{km^2} = 45 \times 10^6\ \mathrm{m^2} \]
    Annual volume \[ V_{ro} = Q \times \text{seconds in year} = 1.85 \times (365.25 \times 24 \times 3600) = 1.85 \times 31\,557\,600 \approx 58\,381\,560\ \mathrm{m^3/yr} \]
    \[ R_o = \frac{V_{ro}}{A} = \frac{58\,381\,560}{45 \times 10^6} = 1.297\ \mathrm{m/yr} = \mathbf{1297\ \mathrm{mm/yr}} \]
    \item \textbf{Groundwater Recharge ($R$):}
    Water Balance: $P = ET + R_o + R$
    \[ R = P - ET - R_o = 3200 - 1232.5 - 1297 = \mathbf{670.5\ \mathrm{mm/yr}} \]
    Volume $V_R = R(\text{m}) \times A = 0.6705 \times 45 \times 10^6 = \mathbf{30.17 \times 10^6\ \mathrm{m^3/yr}}$ (or $30.2\ \mathrm{Mm^3/yr}$).
    \item \textbf{Recharge Coefficient ($k_r$):}
    \[ k_r = \frac{R}{P} \times 100\% = \frac{670.5}{3200} \times 100\% \approx \mathbf{20.95\%} \]
    \textbf{Comment:} A recharge coefficient of $\sim 21\%$ is moderate to high. The flank of Mount Cameroon consists of young, highly permeable basaltic lava flows, scoria cones, and pyroclastic deposits with well-developed fracture networks and lava tubes. This allows rapid infiltration despite high rainfall intensity. The value is lower than theoretical maximum due to high $ET$ (dense tropical forest) and significant surface runoff during extreme events on steep slopes.
\end{enumerate}
\end{exoresolu}

\begin{methode}[Darcy's Law and Aquifer Hydraulic Parameters]
\textbf{Objective:} Calculate specific discharge ($q$), average linear velocity ($v$), transmissivity ($T$), and storativity ($S$) for 1D flow.
\textbf{Key Formulas:}
\begin{itemize}
    \item Darcy: $q = -K \frac{dh}{dl}$ (specific discharge, m/s). $Q = K A \frac{dh}{dl}$ (volumetric flow).
    \item Seepage velocity: $v = \frac{q}{n_e} = \frac{K}{n_e} \frac{dh}{dl}$ ($n_e$ = effective porosity).
    \item Confined Aquifer: $T = K b$ (m$^2$/s), $S = S_s b$ (dimensionless).
    \item Unconfined Aquifer: $T = K h$ (varies with $h$), $S \approx S_y$ (specific yield).
\end{itemize}
\textbf{Steps:}
\begin{enumerate}
    \item \textbf{Identify flow dimension:} 1D horizontal? Radial to well? Vertical leakage?
    \item \textbf{Extract $K$, $b$, $n_e$, $dh/dl$:} Check units (m/s, m, m/m). Convert $K$ from m/day if needed ($1\ \mathrm{m/day} = 1.157 \times 10^{-5}\ \mathrm{m/s}$).
    \item \textbf{Calculate $T$ and $S$:} $T = Kb$ (confined). For unconfined, use saturated thickness $h$ at the point of interest.
    \item \textbf{Compute $v$:} Use $n_e$ (not total porosity $n$). Typical $n_e$: Sand $0.2-0.3$, Fractured basalt $0.01-0.1$, Sandstone $0.1-0.2$.
    \item \textbf{Travel time:} $t = L / v$ for a distance $L$.
\end{enumerate}
\textbf{Examiner's Trap:} Using total porosity $n$ instead of effective porosity $n_e$ for velocity. Confusing $K$ (hydraulic conductivity, m/s) with $T$ (transmissivity, m$^2$/s). Forgetting the minus sign in Darcy (flow from high to low head).
\end{methode}

\begin{exoresolu}[Darcy Flow in the Douala Basin Sands]
A confined aquifer in the Douala Basin (Eocene sands) has a thickness $b = 40\ \mathrm{m}$ and hydraulic conductivity $K = 2.5 \times 10^{-4}\ \mathrm{m/s}$. Two piezometers $500\ \mathrm{m}$ apart show a head difference $\Delta h = 2.5\ \mathrm{m}$ (Piezometer A higher than B). Effective porosity $n_e = 0.22$.
\begin{enumerate}
    \item Calculate the Transmissivity $T$ and Storativity $S$ (given specific storage $S_s = 1.5 \times 10^{-5}\ \mathrm{m^{-1}}$).
    \item Determine the hydraulic gradient $i$ and its direction.
    \item Compute the specific discharge $q$ and the average linear velocity $v$ (in m/day).
    \item How long (in years) for a conservative tracer to travel from A to B?
\end{enumerate}
\tcblower
\textbf{Detailed Solution}
\begin{enumerate}
    \item \textbf{Aquifer Parameters:}
    \[ T = K \times b = (2.5 \times 10^{-4}) \times 40 = \mathbf{0.01\ \mathrm{m^2/s}} \quad (\text{or } 864\ \mathrm{m^2/day}) \]
    \[ S = S_s \times b = (1.5 \times 10^{-5}) \times 40 = \mathbf{6.0 \times 10^{-4}} \quad (\text{dimensionless, typical confined}) \]
    \item \textbf{Hydraulic Gradient ($i$):}
    \[ i = \frac{\Delta h}{L} = \frac{2.5}{500} = \mathbf{0.005\ \mathrm{m/m}} \]
    Direction: From A (high head) to B (low head).
    \item \textbf{Flow Velocities:}
    Specific discharge (Darcy flux):
    \[ q = K i = (2.5 \times 10^{-4}) \times 0.005 = \mathbf{1.25 \times 10^{-6}\ \mathrm{m/s}} \]
    Average linear velocity:
    \[ v = \frac{q}{n_e} = \frac{1.25 \times 10^{-6}}{0.22} = 5.68 \times 10^{-6}\ \mathrm{m/s} \]
    Convert to m/day: $v = 5.68 \times 10^{-6} \times 86400 = \mathbf{0.491\ \mathrm{m/day}}$.
    \item \textbf{Travel Time ($t$):}
    \[ t = \frac{L}{v} = \frac{500\ \mathrm{m}}{0.491\ \mathrm{m/day}} = 1018\ \mathrm{days} \]
    \[ t = \frac{1018}{365.25} \approx \mathbf{2.79\ \mathrm{years}} \]
\end{enumerate}
\end{exoresolu}

\begin{methode}[Well Hydraulics: Thiem Equation for Steady Radial Flow]
\textbf{Objective:} Determine $T$, $K$, drawdown ($s$), or radius of influence ($R$) for a pumping well in equilibrium.
\textbf{Key Formulas (Steady State):}
\begin{itemize}
    \item \textbf{Confined:} $Q = 2\pi T \frac{h_2 - h_1}{\ln(r_2/r_1)} \implies s_1 - s_2 = \frac{Q}{2\pi T} \ln\frac{r_2}{r_1}$.
    \item \textbf{Unconfined:} $Q = \pi K \frac{h_2^2 - h_1^2}{\ln(r_2/r_1)}$ (Dupuit assumptions). Approx: $s_1 - s_2 \approx \frac{Q}{2\pi T} \ln\frac{r_2}{r_1}$ if $s \ll h$.
    \item \textbf{Radius of Influence ($R$):} $s_w = \frac{Q}{2\pi T} \ln\frac{R}{r_w}$ (Confined). $R$ often estimated by Sichardt: $R = 3000 s_w \sqrt{K}$ ($s_w$ in m, $K$ in m/s).
\end{itemize}
\textbf{Steps:}
\begin{enumerate}
    \item \textbf{Identify aquifer type:} Confined (use $h$) or Unconfined (use $h^2$).
    \item \textbf{Plot/Select data:} Use observation wells at $r_1, r_2$ with drawdowns $s_1, s_2$. Or use well radius $r_w$ and $R$.
    \item \textbf{Calculate $T$:} $T = \frac{Q}{2\pi (s_1 - s_2)} \ln\frac{r_2}{r_1}$ (Confined).
    \item \textbf{Calculate $K$:} $K = T / b$.
    \item \textbf{Check assumptions:} Steady state? Fully penetrating? Homogeneous, isotropic? No recharge boundaries?
\end{enumerate}
\textbf{Examiner's Reflex:} If given drawdown in well ($s_w$) and $Q$ only, you \emph{must} estimate $R$ (Sichardt) to find $T$. State the assumption clearly. For unconfined, use $h^2$ difference unless $s < 0.1 h$.
\end{methode}

\begin{exoresolu}[Pumping Test in the Yaoundé Fractured Basement]
A well fully penetrates a confined fractured aquifer (thickness $b=30\ \mathrm{m}$) in the Yaoundé Pan-African granitoids. Pumping rate $Q = 15\ \mathrm{m^3/h}$. Two observation wells give:
\begin{itemize}
    \item OW1: $r_1 = 10\ \mathrm{m}$, $s_1 = 4.2\ \mathrm{m}$
    \item OW2: $r_2 = 50\ \mathrm{m}$, $s_2 = 1.8\ \mathrm{m}$
\end{itemize}
Well radius $r_w = 0.15\ \mathrm{m}$.
\begin{enumerate}
    \item Calculate Transmissivity $T$ (m$^2$/day) and Hydraulic Conductivity $K$ (m/s).
    \item Estimate the radius of influence $R$ using the Sichardt formula.
    \item Calculate the drawdown in the pumping well ($s_w$) assuming no well losses.
    \item If the measured drawdown in the well is $s_{w,meas} = 12.5\ \mathrm{m}$, calculate the well efficiency.
\end{enumerate}
\tcblower
\textbf{Detailed Solution}
\begin{enumerate}
    \item \textbf{Transmissivity ($T$):} Use Thiem (Confined) between OW1 and OW2 (eliminates well losses/R).
    \[
    Q = 15\ \mathrm{m^3/h} = \frac{15}{3600} = 4.167 \times 10^{-3}\ \mathrm{m^3/s}.
    \]
    Better to work in m$^3$/day: $Q = 15 \times 24 = \mathbf{360\ \mathrm{m^3/day}}$.
    \[ T = \frac{Q}{2\pi (s_1 - s_2)} \ln\frac{r_2}{r_1} = \frac{360}{2\pi (4.2 - 1.8)} \ln\frac{50}{10} \]
    \[ T = \frac{360}{2\pi \times 2.4} \times \ln(5) = \frac{360}{15.08} \times 1.609 = 23.87 \times 1.609 = \mathbf{38.4\ \mathrm{m^2/day}} \]
    \[ K = \frac{T}{b} = \frac{38.4}{30} = 1.28\ \mathrm{m/day} = \frac{1.28}{86400} = \mathbf{1.48 \times 10^{-5}\ \mathrm{m/s}} \]
    \item \textbf{Radius of Influence ($R$):} Sichardt: $R = 3000 s_w \sqrt{K}$. Need $s_w$ first? Iterative. Approximate $s_w$ using OW1.
    Drawdown at well (theoretical, no loss): $s_{w,th} = s_1 + \frac{Q}{2\pi T} \ln\frac{r_1}{r_w}$.
    \[ \frac{Q}{2\pi T} = \frac{360}{2\pi \times 38.4} = 1.492\ \mathrm{m} \]
    \[ s_{w,th} = 4.2 + 1.492 \times \ln(10/0.15) = 4.2 + 1.492 \times \ln(66.67) = 4.2 + 1.492 \times 4.20 = 4.2 + 6.26 = 10.46\ \mathrm{m} \]
    Now $R = 3000 \times 10.46 \times \sqrt{1.48 \times 10^{-5}} = 31380 \times 0.00385 = \mathbf{121\ \mathrm{m}}$.
    (Check: $s_w = \frac{Q}{2\pi T} \ln(R/r_w) = 1.492 \times \ln(121/0.15) = 1.492 \times 6.69 = 9.98\ \mathrm{m}$. Close enough).
    \item \textbf{Theoretical Well Drawdown ($s_{w,th}$):} $\mathbf{10.5\ \mathrm{m}}$ (using $R=121$ m) or $\mathbf{10.5\ \mathrm{m}}$ (using OW1 extrapolation).
    \item \textbf{Well Efficiency ($E_w$):}
    \[ E_w = \frac{s_{w,th}}{s_{w,meas}} \times 100\% = \frac{10.5}{12.5} \times 100\% = \mathbf{84\%} \]
    \textbf{Interpretation:} $16\%$ head loss due to turbulent flow in fractures near wellbore (skin effect) or partial penetration. Good efficiency for a fractured rock well.
\end{enumerate}
\end{exoresolu}

\begin{methode}[Spring Hydrograph Analysis and Recession Constants]
\textbf{Objective:} Characterize aquifer storage and geometry from spring discharge recession limb ($Q_t = Q_0 e^{-\alpha t}$).
\textbf{Key Concepts:}
\begin{itemize}
    \item Recession limb: $\ln Q_t = \ln Q_0 - \alpha t$. Plot $\ln Q$ vs $t$; slope $= -\alpha$.
    \item Multiple segments $\to$ multiple reservoirs (e.g., $\alpha_1$ fast: epikarst/fractures; $\alpha_2$ slow: matrix/conduit).
    \item Volume drained: $V = \int Q dt = Q_0 / \alpha$ (for one exponential).
    \item Maillet formula: $Q = Q_0 e^{-\alpha t}$. $\alpha$ (day$^{-1}$) = recession coefficient. $k = 1/\alpha$ = reservoir capacity (days).
    \item Baseflow separation: Graphical (straight line, fixed base length $N = A^{0.2}$) or Digital filter (Lyne-Hollick).
\end{itemize}
\textbf{Steps:}
\begin{enumerate}
    \item \textbf{Select recession segment:} After peak, no rain. Plot $\ln Q$ vs $t$.
    \item \textbf{Fit line(s):} Identify breaks in slope. Calculate $\alpha_i$ for each segment.
    \item \textbf{Interpret:} High $\alpha$ (steep) = low storage, quick flow (karst conduits, fractures). Low $\alpha$ (flat) = high storage, slow flow (porous matrix, deep aquifer).
    \item \textbf{Calculate reserves:} $V_i = Q_{0,i} / \alpha_i$ (instantaneous reserves at start of segment $i$).
    \item \textbf{Regime classification:} Permanent ($\alpha < 0.005$), Temporary ($\alpha > 0.02$), etc. (Mangin classification for karst).
\end{enumerate}
\textbf{Examiner's Trap:} Fitting a single line to a multi-segment recession. Using arithmetic $Q$ vs $t$ instead of semi-log. Confusing $\alpha$ (recession coeff) with $\alpha$ (dispersivity) -- context!
\end{methode}

\begin{exoresolu}[Recession Analysis of a Spring in the Mandara Mountains]
A spring draining a sandstone aquifer in the Mandara Mountains (Mokolo area) has the following recession discharge data after the rainy season:
\begin{center}
\begin{tabular}{|c|c|c|c|c|c|c|c|c|}
\hline
Time $t$ (days) & 0 & 10 & 20 & 30 & 40 & 50 & 60 & 70 \\
\hline
Discharge $Q$ (L/s) & 120 & 85 & 60 & 43 & 31 & 22 & 16 & 11 \\
\hline
\end{tabular}
\end{center}
\begin{enumerate}
    \item Plot $\ln Q$ vs $t$ (describe the plot) and determine if a single reservoir model fits.
    \item Calculate the recession coefficient $\alpha$ (day$^{-1}$) and the reservoir constant $k$ (days).
    \item Estimate the dynamic reserves (volume) drained between $t=0$ and $t=70$ days.
    \item Classify the spring regime based on $\alpha$.
\end{enumerate}
\tcblower
\textbf{Detailed Solution}
\begin{enumerate}
    \item \textbf{Semi-log Plot Analysis:}
    Compute $\ln Q$:
    $t=0: \ln 120 = 4.787$; $10: \ln 85 = 4.443$; $20: \ln 60 = 4.094$; $30: \ln 43 = 3.761$;
    $40: \ln 31 = 3.434$; $50: \ln 22 = 3.091$; $60: \ln 16 = 2.773$; $70: \ln 11 = 2.398$.
    Plotting $\ln Q$ vs $t$ yields a \textbf{straight line} (check differences: $\Delta \ln Q / \Delta t \approx -0.034$ constant). A single exponential model fits well.
    \item \textbf{Recession Coefficient ($\alpha$):}
    Slope $= -\alpha$. Using endpoints:
    \[ \alpha = -\frac{\ln Q_{70} - \ln Q_0}{70 - 0} = -\frac{2.398 - 4.787}{70} = \frac{2.389}{70} = \mathbf{0.0341\ \mathrm{day^{-1}}} \]
    Reservoir constant (time constant):
    \[ k = \frac{1}{\alpha} = \frac{1}{0.0341} = \mathbf{29.3\ \mathrm{days}} \]
    \item \textbf{Dynamic Reserves ($V$):}
    Volume drained from $t=0$ to $\infty$: $V_{total} = Q_0 / \alpha$.
    $Q_0 = 120\ \mathrm{L/s} = 0.12\ \mathrm{m^3/s}$.
    \[ V_{total} = \frac{0.12}{0.0341 / 86400} = \frac{0.12 \times 86400}{0.0341} = \frac{10368}{0.0341} \approx \mathbf{304\,000\ \mathrm{m^3}} \]
    Volume drained between $t=0$ and $t=70$:
    \[ V_{0-70} = V_{total} (1 - e^{-\alpha \times 70}) = 304\,000 \times (1 - e^{-2.389}) = 304\,000 \times (1 - 0.092) = \mathbf{276\,000\ \mathrm{m^3}} \]
    (Alternatively: Trapezoidal integration on $Q(t)$ table gives similar result).
    \item \textbf{Regime Classification:}
    $\alpha = 0.034\ \mathrm{day^{-1}}$. Mangin (Karst) limits: Temporary ($\alpha > 0.02$), Seasonal ($0.005 < \alpha < 0.02$), Permanent ($\alpha < 0.005$).
    For sandstone (porous/fractured), similar logic applies. $\alpha = 0.034$ indicates a \textbf{Temporary / Seasonal regime with relatively low storage capacity} (reserves last ~1 month). Typical for fractured sandstone with limited matrix porosity in a semi-arid climate (Mandara).
\end{enumerate}
\end{exoresolu}

\begin{methode}[Hydrochemical Facies Classification (Piper / Schoeller) and Quality Indices]
\textbf{Objective:} Classify water type, identify mixing/contamination, and assess suitability (drinking/irrigation).
\textbf{Key Tools:}
\begin{itemize}
    \item \textbf{Piper Diagram:} Trilinear (cations + anions) $\to$ Central diamond. Facies: Ca-HCO$_3$ (recharge), Na-HCO$_3$ (cation exchange), Na-Cl (saline/evaporative), Ca-SO$_4$ (gypsum/pollution).
    \item \textbf{Schoeller Diagram:} Semi-log plot of meq/L vs ions. Parallel lines = same origin. Crossing lines = mixing/pollution.
    \item \textbf{Indices:} SAR $= \frac{Na^+}{\sqrt{(Ca^{2+}+Mg^{2+})/2}}$ (irrigation). \%Na, RSC, Hardness (TH), TDS.
    \item \textbf{WHO Standards:} NO$_3^- < 50$ mg/L, F$^- < 1.5$ mg/L, Fe $< 0.3$ mg/L, TDS $< 1000$ mg/L, E. coli $= 0$.
\end{itemize}
\textbf{Steps:}
\begin{enumerate}
    \item \textbf{Convert mg/L to meq/L:} $\text{meq/L} = \frac{\text{mg/L}}{\text{Eq. weight}}$ (Ca: 20, Mg: 12.15, Na: 23, K: 39.1, HCO$_3$: 61, Cl: 35.45, SO$_4$: 48, NO$_3$: 62).
    \item \textbf{Check Electroneutrality:} $\sum \text{cations} \approx \sum \text{anions}$ (Error $< 5\%$).
    \item \textbf{Plot Piper:} %meq in each triangle. Read facies in diamond.
    \item \textbf{Calculate Indices:} SAR, \%Na, TH (as CaCO$_3$ = 2.5 Ca + 4.1 Mg meq/L).
    \item \textbf{Interpret:} Evolution along flow path? Seawater intrusion (Cl/Br, Na/Cl)? Anthropogenic (NO$_3$, K, SO$_4$)?
\end{enumerate}
\textbf{Examiner's Reflex:} Always verify charge balance first. If error $> 10\%$, data is unreliable. For Piper, % are based on meq, not mg/L. SAR uses meq/L.
\end{methode}

\begin{exoresolu}[Water Quality Assessment for Irrigation in the Logone Valley]
A groundwater sample from a shallow well in the Logone alluvial plain (Far North) gives (mg/L):
Ca$^{2+}=80$, Mg$^{2+}=35$, Na$^+=120$, K$^+=10$, HCO$_3^-=280$, Cl$^-=150$, SO$_4^{2-}=60$, NO$_3^-=45$.
\begin{enumerate}
    \item Convert to meq/L and verify charge balance error (CBE).
    \item Determine the hydrochemical facies (Piper classification).
    \item Calculate SAR, \%Na, and Total Hardness (TH as mg/L CaCO$_3$).
    \item Assess suitability for irrigation (Wilcox/US Salinity Diagram logic) and drinking (WHO NO$_3$ limit).
\end{enumerate}
\tcblower
\textbf{Detailed Solution}
\begin{enumerate}
    \item \textbf{Conversion to meq/L (Eq. wts: Ca=20, Mg=12.15, Na=23, K=39.1, HCO$_3$=61, Cl=35.45, SO$_4$=48, NO$_3$=62):}
    \begin{center}
    \begin{tabular}{|l|c|c|c|c|}
    \hline
    Ion & mg/L & Eq. wt & meq/L & \%meq \\
    \hline
    Ca$^{2+}$ & 80 & 20.0 & 4.00 & 34.5 \\
    Mg$^{2+}$ & 35 & 12.15 & 2.88 & 24.8 \\
    Na$^+$ & 120 & 23.0 & 5.22 & 45.0 \\
    K$^+$ & 10 & 39.1 & 0.26 & 2.2 \\
    \hline
    \textbf{Sum Cat.} & & & \textbf{12.36} & \textbf{100} \\
    \hline
    HCO$_3^-$ & 280 & 61.0 & 4.59 & 37.5 \\
    Cl$^-$ & 150 & 35.45 & 4.23 & 34.6 \\
    SO$_4^{2-}$ & 60 & 48.0 & 1.25 & 10.2 \\
    NO$_3^-$ & 45 & 62.0 & 0.73 & 6.0 \\
    \hline
    \textbf{Sum An.} & & & \textbf{10.80} & \textbf{100} \\
    \hline
    \end{tabular}
    \end{center}
    CBE $= \frac{|\sum Cat - \sum An|}{\sum Cat + \sum An} \times 100 = \frac{|12.36 - 10.80|}{23.16} \times 100 = \mathbf{6.7\%}$.
    \textbf{Acceptable} ($< 10\%$), slight cation excess (possible unmeasured Fe$^{2+}$, or HCO$_3$ titration error).
    \item \textbf{Facies (Piper):}
    Cations: Na+K (47\%) > Ca (34.5\%) > Mg (24.8\%) $\to$ \textbf{Na-Ca} type (or Sodic-Calcic).
    Anions: HCO$_3$ (37.5\%) $\approx$ Cl (34.6\%) > SO$_4$ (10.2\%) $\to$ \textbf{HCO$_3$-Cl} type (Bicarbonate-Chloride).
    \textbf{Overall Facies: Na-Ca-HCO$_3$-Cl} (Mixed type, typical of evaporation/concentration in semi-arid alluvium or mixing with deeper saline water).
    \item \textbf{Indices:}
    SAR $= \frac{Na^+}{\sqrt{(Ca^{2+}+Mg^{2+})/2}} = \frac{5.22}{\sqrt{(4.00+2.88)/2}} = \frac{5.22}{\sqrt{3.44}} = \frac{5.22}{1.85} = \mathbf{2.82}$.
    \%Na $= \frac{Na+K}{Sum Cat} \times 100 = \frac{5.48}{12.36} \times 100 = \mathbf{44.3\%}$.
    TH (meq/L) $= Ca + Mg = 4.00 + 2.88 = 6.88\ \mathrm{meq/L}$.
    TH (mg/L CaCO$_3$) $= 6.88 \times 50 = \mathbf{344\ \mathrm{mg/L}}$ (Hard water).
    \item \textbf{Assessment:}
    \textbf{Irrigation:} SAR = 2.82 (Low), \%Na = 44\% (Medium). On US Salinity Diagram (EC $\approx$ TDS/0.64; TDS $\approx \sum$mg/L $\approx 780$ mg/L $\to$ EC $\approx 1220\ \mu$S/cm = Medium Salinity C2). Class \textbf{C2-S1} (Medium Salinity, Low Sodium). \textbf{Good for most crops on permeable soils.} Warning: High Cl (4.23 meq/L) may affect chloride-sensitive crops (onion, citrus).
    \textbf{Drinking:} NO$_3^- = 45\ \mathrm{mg/L}$. WHO limit = 50 mg/L. \textbf{Borderline/Exceeds if analytical error positive.} Likely anthropogenic (latrines, fertilizer). Requires treatment or source protection. TH = 344 mg/L $> 300$ mg/L (High hardness, scaling, taste issues).
\end{enumerate}
\end{exoresolu}

\begin{methode}[Saltwater Intrusion: Ghyben-Herzberg Relation and Interface Position]
\textbf{Objective:} Estimate the depth of the freshwater-saltwater interface below sea level and the critical pumping limit for coastal wells.
\textbf{Key Formula (Static, Sharp Interface):}
\[ z = \frac{\rho_f}{\rho_s - \rho_f} h \approx 40 h \]
($z$ = depth of interface below sea level, $h$ = freshwater head above sea level, $\rho_f=1.00$, $\rho_s=1.025$ g/cm$^3$).
\textbf{Critical Drawdown:} To avoid upconing, max drawdown $s_{max} \approx h/3$ (or keep interface below screen).
\textbf{Steps:}
\begin{enumerate}
    \item \textbf{Measure/Estimate $h$:} Water level in well (m above MSL).
    \item \textbf{Calculate $z$:} $z = 40 h$. Interface toe inland distance $L \approx \frac{K}{Q} \frac{h^2}{2}$ (steady state, 1D).
    \item \textbf{Pumping limit:} For a well at distance $x$ from coast, max allowable drawdown $s_{crit} \approx h - \frac{z_{screen}}{40}$.
    \item \textbf{Monitor:} Cl$^-$ concentration, EC logs. Dispersion zone $\to$ transition zone thickness $\delta \approx \frac{0.1 v}{\alpha_L}$ (often 1-10 m).
\end{enumerate}
\textbf{Examiner's Trap:} Using $z=40h$ for pumping conditions (dynamic) -- it's only valid for static equilibrium. Forgetting that $h$ is measured \emph{above mean sea level}, not ground level. Confusing $\rho_s - \rho_f = 0.025$ with $0.025/1.00 = 1/40$.
\end{methode}

\begin{exoresolu}[Coastal Aquifer Management in Douala (Wouri Estuary)]
A monitoring well in the Douala Quaternary sands (unconfined, $K=20\ \mathrm{m/day}$) is located $500\ \mathrm{m}$ inland from the Wouri Estuary. The static water level in the well is $h = 1.2\ \mathrm{m}$ above mean sea level (MSL). The well screen is at $-15\ \mathrm{m}$ MSL (i.e., 15 m below sea level). Aquifer base is at $-40\ \mathrm{m}$ MSL.
\begin{enumerate}
    \item Using Ghyben-Herzberg, calculate the static depth of the saltwater interface below MSL at the well location.
    \item Is the well screen currently in freshwater or saltwater?
    \item Estimate the maximum allowable drawdown ($s_{max}$) in this well to prevent the interface from rising to the top of the screen (assume sharp interface).
    \item If the well pumps at $Q = 50\ \mathrm{m^3/h}$, estimate the steady-state drawdown at the well using Thiem (unconfined, $R=1000\ \mathrm{m}$, $r_w=0.2\ \mathrm{m}$). Is it safe?
\end{enumerate}
\tcblower
\textbf{Detailed Solution}
\begin{enumerate}
    \item \textbf{Static Interface Depth ($z$):}
    \[ z = 40 \times h = 40 \times 1.2 = \mathbf{48\ \mathrm{m}\ \text{below MSL}} \]
    \item \textbf{Screen Position:} Screen at $-15\ \mathrm{m}$ MSL. Interface at $-48\ \mathrm{m}$ MSL.
    Since $-15 > -48$ (screen is \emph{above} the interface), the screen is in the \textbf{Freshwater zone}. Safe statically.
    \item \textbf{Critical Drawdown ($s_{max}$):}
    We need the interface to stay at or below the screen top ($-15\ \mathrm{m}$ MSL).
    Required head $h_{crit}$ such that $z_{crit} = 15\ \mathrm{m}$.
    \[ h_{crit} = \frac{z_{crit}}{40} = \frac{15}{40} = 0.375\ \mathrm{m}\ \text{above MSL} \]
    Current head $h = 1.2\ \mathrm{m}$.
    Max allowable drawdown $s_{max} = h - h_{crit} = 1.2 - 0.375 = \mathbf{0.825\ \mathrm{m}}$.
    \item \textbf{Pumping Drawdown Check:}
    $Q = 50\ \mathrm{m^3/h} = 1200\ \mathrm{m^3/day}$.
    Unconfined Thiem (approx using $h_{avg}$): $Q = \pi K \frac{h_R^2 - h_w^2}{\ln(R/r_w)}$.
    $h_R = h = 1.2\ \mathrm{m}$ (head at $R$). Saturated thickness at coast $\approx$ sea level? No, $h_R$ is head at radius of influence. Assume $h_R \approx 1.2$ m (water table flat at $R$).
    Wait: Unconfined aquifer base at $-40$ m. Saturated thickness at well initially $b_0 = 1.2 - (-40) = 41.2\ \mathrm{m}$.
    At radius $R=1000$ m, assume $h_R = 1.2$ m (head above MSL).
    Let's use heads above base (datum = base = -40 m).
    $H_R = 41.2\ \mathrm{m}$. $H_w = ?$
    $Q = \pi K \frac{H_R^2 - H_w^2}{\ln(R/r_w)}$
    $1200 = \pi \times 20 \times \frac{41.2^2 - H_w^2}{\ln(1000/0.2)} = 62.83 \times \frac{1697 - H_w^2}{8.52} = 7.37 (1697 - H_w^2)$
    $162.8 = 1697 - H_w^2 \implies H_w^2 = 1534 \implies H_w = 39.2\ \mathrm{m}$.
    Drawdown $s = H_R - H_w = 41.2 - 39.2 = \mathbf{2.0\ \mathrm{m}}$.
    \textbf{Comparison:} $s_{calc} = 2.0\ \mathrm{m} > s_{max} = 0.825\ \mathrm{m}$.
    \textbf{Conclusion:} Pumping at $50\ \mathrm{m^3/h}$ causes \textbf{excessive drawdown}, risking saltwater upconing into the screen. Rate must be reduced to $\approx 20\ \mathrm{m^3/h}$ or well relocated inland.
\end{enumerate}
\end{exoresolu}

\begin{methode}[Vertical Electrical Sounding (VES) Interpretation for Aquifer Delineation]
\textbf{Objective:} Interpret Schlumberger VES curves to determine layer resistivities ($\rho$), thicknesses ($h$), and lithology/hydrogeology.
\textbf{Key Concepts:}
\begin{itemize}
    \item Apparent resistivity $\rho_a = K \frac{\Delta V}{I}$, $K = \pi \frac{(AB/2)^2 - (MN/2)^2}{MN} \approx \pi \frac{(AB/2)^2}{MN}$ for $AB \gg MN$.
    \item Curve Types (3 layers): K ($\rho_1 > \rho_2 < \rho_3$), H ($\rho_1 < \rho_2 > \rho_3$), A ($\rho_1 > \rho_2 > \rho_3$), Q ($\rho_1 < \rho_2 < \rho_3$).
    \item Hydrogeological mapping: Low $\rho$ ($< 50\ \Omega$m) $\to$ Clay/Shale/Saline water. Medium $\rho$ ($50-500\ \Omega$m) $\to$ Sand/Sandstone (Freshwater). High $\rho$ ($> 500\ \Omega$m) $\to$ Basement/Granite/Limestone (Dry/Fractured) or Freshwater in clean sand.
    \item Equivalence/Suppression: Thin layers ($h < 0.1 \times$ depth) may be suppressed. Use auxiliary points / joint inversion.
\end{itemize}
\textbf{Steps (Manual 2-Layer Master Curve / Software):}
\begin{enumerate}
    \item \textbf{Plot $\rho_a$ vs $AB/2$ (log-log).}
    \item \textbf{Identify curve type} (K, H, A, Q) $\to$ gives resistivity contrast sequence.
    \item \textbf{Match to Master Curves} (2-layer) or use software (IPI2Win, WinGLink) for $n$-layer inversion.
    \item \textbf{Extract Model:} $\rho_1, h_1, \rho_2, h_2, \dots, \rho_n$.
    \item \textbf{Geological Correlation:} Match $\rho_i, h_i$ to borehole logs (lithology, water strikes). Define Aquifer (permeable, $\rho$ medium) vs Aquiclude (clay, $\rho$ low) vs Bedrock (high $\rho$).
    \item \textbf{Transverse Resistance ($T_r = \sum \rho_i h_i$) and Longitudinal Conductance ($S_l = \sum h_i/\rho_i$)} for aquifer protection mapping.
\end{enumerate}
\textbf{Examiner's Reflex:} Always state the ambiguity (Equivalence). "The interpreted thickness $h_2$ has an equivalence range..." Correlate with \emph{at least one} borehole log. Never interpret VES in isolation.
\end{methode}

\begin{exoresolu}[VES Interpretation for Borehole Siting in the Benue Basin (Garoua)]
A Schlumberger VES (VES-07) near Garoua yields the following apparent resistivity curve (field data smoothed):
\begin{center}
\begin{tabular}{|c|c|c|c|c|c|c|c|c|c|}
\hline
AB/2 (m) & 1.5 & 3 & 6 & 10 & 15 & 25 & 40 & 60 & 100 \\
\hline
$\rho_a$ ($\Omega$m) & 120 & 110 & 95 & 85 & 70 & 45 & 30 & 25 & 22 \\
\hline
\end{tabular}
\end{center}
A nearby borehole log shows: 0-3 m: Lateritic topsoil; 3-18 m: Clayey sand; 18-35 m: Medium-coarse sand (Water strike at 22 m); 35-40 m: Clay; $>40$ m: Sandstone bedrock.
\begin{enumerate}
    \item Identify the VES curve type (K, H, A, Q) and the number of layers resolved.
    \item Propose a 4-layer quantitative model ($\rho_i, h_i$) consistent with the borehole log.
    \item Calculate the Transverse Resistance ($T_r$) and Longitudinal Conductance ($S_l$) of the aquifer layer (Layer 3).
    \item Determine the recommended borehole depth and screen setting.
\end{enumerate}
\tcblower
\textbf{Detailed Solution}
\begin{enumerate}
    \item \textbf{Curve Type:} $\rho_a$ decreases monotonically from 120 to 22 $\Omega$m. $\rho_1 > \rho_2 > \rho_3 > \rho_4$. \textbf{Type A (Descending)}. 4 distinct slope changes $\to$ \textbf{4 layers resolved} (possibly 5 if basement seen).
    \item \textbf{Quantitative Model (Constrained by Borehole):}
    \begin{center}
    \begin{tabular}{|c|c|c|c|c|}
    \hline
    Layer & Lithology (Log) & Thickness $h$ (m) & Depth Range (m) & Resistivity $\rho$ ($\Omega$m) \\
    \hline
    1 & Lateritic Topsoil & 3 & 0-3 & 120 (High, dry/ferruginous) \\
    2 & Clayey Sand & 15 & 3-18 & 40 (Low, clay conductive) \\
    3 & \textbf{Aquifer: Clean Sand} & 17 & 18-35 & \textbf{150 (Medium-High, Freshwater)} \\
    4 & Clay / Weathered Bedrock & 5 & 35-40 & 20 (Low, clay) \\
    5 & Sandstone Bedrock & $\infty$ & $>40$ & 500 (High, consolidated) \\
    \hline
    \end{tabular}
    \end{center}
    \textit{Note:} The field curve drops to 22 $\Omega$m at AB/2=100m. The model $\rho_4=20$ fits the asymptote. The high $\rho_3=150$ creates the "shoulder" around AB/2=15-25m (transition from Layer 2 to 3). The drop at large AB/2 is Layer 4/5.
    \item \textbf{Aquifer Parameters (Layer 3):}
    $h_3 = 17\ \mathrm{m}$, $\rho_3 = 150\ \Omega$m.
    Transverse Resistance: $T_r = \rho_3 \times h_3 = 150 \times 17 = \mathbf{2550\ \Omega\,m^2}$.
    Longitudinal Conductance: $S_l = h_3 / \rho_3 = 17 / 150 = \mathbf{0.113\ \mathrm{S}}$ (Siemens).
    \textbf{Significance:} High $T_r$ ($> 1000\ \Omega\,m^2$) indicates a \textbf{good aquifer potential} (thick, resistive). Low $S_l$ indicates low vulnerability to surface contamination (protected by Layer 2 clayey cap).
    \item \textbf{Borehole Design:}
    Target: Layer 3 (18-35 m).
    \textbf{Depth:} Drill to $35\ \mathrm{m}$ (base of aquifer / top of clay Layer 4).
    \textbf{Screen Setting:} $18\ \mathrm{m}$ to $35\ \mathrm{m}$ (Full penetration of aquifer).
    \textbf{Casing:} $0-18\ \mathrm{m}$ (Seal through Layer 1 \& 2).
    \textbf{Gravel Pack:} Annulus opposite screen.
    \textbf{Sanitary Seal:} Top 3-5 m cement grout.
\end{enumerate}
\end{exoresolu}

\begin{methode}[Groundwater Contamination: Advection-Dispersion and Retardation]
\textbf{Objective:} Estimate contaminant plume migration velocity and arrival time at a receptor (well/spring).
\textbf{Key Formulas (1D, Uniform Flow):}
\begin{itemize}
    \item Average linear velocity: $v = \frac{K i}{n_e}$.
    \item Retardation Factor: $R_f = 1 + \frac{\rho_b K_d}{n_e}$ ($\rho_b$ = bulk density, $K_d$ = distribution coeff).
    \item Contaminant Velocity: $v_c = \frac{v}{R_f}$.
    \item Longitudinal Dispersivity: $\alpha_L \approx 0.1 \times \text{plume length}$ (or $1-10\ \mathrm{m}$ for field scale).
    \item Concentration (Continuous source, 1D): $C(x,t) = \frac{C_0}{2} \left[ \text{erfc}\left(\frac{x - v_c t}{2\sqrt{D_L t}}\right) + \exp\left(\frac{v_c x}{D_L}\right) \text{erfc}\left(\frac{x + v_c t}{2\sqrt{D_L t}}\right) \right]$.
    \item Approx. Arrival Time (Advection dominant): $t_{arrival} \approx \frac{x R_f}{v}$.
\end{itemize}
\textbf{Steps:}
\begin{enumerate}
    \item \textbf{Determine Flow Field:} $K, i, n_e \to v$.
    \item \textbf{Determine Sorption:} $K_d$ (lab/field), $\rho_b \approx 1.6-2.0\ \mathrm{g/cm^3}$. Calculate $R_f$. Conservative tracer ($Cl^-$, $Br^-$): $R_f=1$. Heavy metals/Organics: $R_f > 1$.
    \item \textbf{Calculate $v_c = v / R_f$.}
    \item \textbf{Estimate Arrival Time:} $t = x / v_c$.
    \item \textbf{Assess Dispersion:} $D_L = \alpha_L v + D^*$ (molecular diffusion negligible usually). Plume spreading $\sigma_x \approx \sqrt{2 D_L t}$.
\end{enumerate}
\textbf{Examiner's Trap:} Using Darcy velocity $q$ instead of $v$ for transport. Forgetting $R_f$ for reactive contaminants. Assuming $K_d$ is constant (linear isotherm) when it may be pH/redox dependent.
\end{methode}

\begin{exoresolu}[Nitrate Plume Migration in the Yaoundé Alluvial Aquifer]
A pit latrine leak introduces nitrate ($NO_3^-$, conservative, $R_f=1$) into a shallow unconfined aquifer (alluvial sand) in a Yaoundé suburb.
Aquifer properties: $K = 10\ \mathrm{m/day}$, $i = 0.002$, $n_e = 0.25$.
A community well is located $x = 150\ \mathrm{m}$ downgradient.
Longitudinal dispersivity $\alpha_L = 5\ \mathrm{m}$.
\begin{enumerate}
    \item Calculate the groundwater seepage velocity $v$ and the nitrate velocity $v_c$.
    \item Estimate the arrival time of the nitrate front at the well (advection only).
    \item Calculate the longitudinal dispersion coefficient $D_L$.
    \item Estimate the time for the concentration at the well to reach $50\%$ of the source concentration ($C/C_0 = 0.5$) using the 1D advection-dispersion equation (ADE) approximation for a continuous source. (Use $C/C_0 = 0.5 \text{erfc}(\frac{x-v_c t}{2\sqrt{D_L t}})$ approx for large $Pe$).
\end{enumerate}
\tcblower
\textbf{Detailed Solution}
\begin{enumerate}
    \item \textbf{Velocities:}
    \[ v = \frac{K i}{n_e} = \frac{10 \times 0.002}{0.25} = \frac{0.02}{0.25} = \mathbf{0.08\ \mathrm{m/day}} \]
    Nitrate is conservative ($R_f = 1$), so $v_c = v = \mathbf{0.08\ \mathrm{m/day}}$.
    \item \textbf{Advection Arrival Time ($t_{adv}$):}
    \[ t_{adv} = \frac{x}{v_c} = \frac{150}{0.08} = 1875\ \mathrm{days} \approx \mathbf{5.1\ \mathrm{years}} \]
    \item \textbf{Dispersion Coefficient ($D_L$):}
    \[ D_L = \alpha_L v_c + D^* \approx 5 \times 0.08 = \mathbf{0.4\ \mathrm{m^2/day}} \]
    \item \textbf{Time for $C/C_0 = 0.5$:}
    Using the approximation for a continuous source: $C/C_0 = \frac{1}{2} \text{erfc}\left(\frac{x - v_c t}{2\sqrt{D_L t}}\right)$.
    Set $C/C_0 = 0.5$:
    \[ 0.5 = 0.5 \text{erfc}\left(\frac{x - v_c t}{2\sqrt{D_L t}}\right) \implies \text{erfc}(\dots) = 1 \implies \frac{x - v_c t}{2\sqrt{D_L t}} = 0 \]
    Thus $x = v_c t \implies t = \frac{x}{v_c} = \frac{150}{0.08} = 1875\ \mathrm{days} \approx \mathbf{5.1\ \mathrm{years}}$.
    \textbf{Note:} For this 1D step-input solution, the 50\% concentration contour arrives exactly at the advective travel time. Dispersion causes earlier arrival of low concentrations and later arrival of high concentrations, but the midpoint is unchanged.
\end{enumerate}
\end{exoresolu}

\newpage

\coursec{Exercises}

\courssub{I check my knowledge}

\begin{exercice}[Multiple choice questions]
For each question, choose the single correct answer.
\begin{enumerate}
\item The process by which water changes from liquid to vapour in the atmosphere is:
\begin{enumerate}
\item condensation \item evaporation \item infiltration \item precipitation
\end{enumerate}
\item A rock that is porous but does NOT allow water to flow through it easily is said to have:
\begin{enumerate}
\item high permeability \item low permeability \item high hydraulic conductivity \item no porosity
\end{enumerate}
\item The upper surface of the zone of saturation is called:
\begin{enumerate}
\item the aquiclude \item the piezometric surface \item the water table \item the capillary fringe
\end{enumerate}
\item Which of the following is the best aquifer?
\begin{enumerate}
\item shale \item unfractured granite \item well-sorted coarse sandstone \item clay
\end{enumerate}
\end{enumerate}
\end{exercice}

\begin{exercice}[True or false — justify]
State whether each statement is true or false, and justify your answer in one or two sentences.
\begin{enumerate}
\item All porous rocks are good aquifers.
\item The water table always lies at the same depth below the surface throughout the year.
\item An artesian well is one in which water rises above the top of the confined aquifer without pumping.
\item Groundwater contamination is easily reversed once the source of pollution is removed.
\end{enumerate}
\end{exercice}

\begin{exercice}[Definitions]
Define each of the following terms in one or two sentences, giving one example where possible:
\begin{enumerate}
\item infiltration;
\item aquiclude;
\item spring;
\item hydraulic gradient;
\item zone of aeration.
\end{enumerate}
\end{exercice}

\begin{exercice}[Direct calculation — water balance]
At a station near Bafoussam, the mean annual precipitation is $P = 1800\ \mathrm{mm}$, the mean annual evapotranspiration is $E = 1200\ \mathrm{mm}$ and the surface runoff is $R = 350\ \mathrm{mm}$.
\begin{enumerate}
\item Using the water balance equation $P = E + R + I$, calculate the annual infiltration $I$.
\item Express $I$ as a percentage of $P$.
\item State one factor that could increase $I$ at this station.
\end{enumerate}
\end{exercice}

\courssub{I practise}

\begin{exercice}[Data analysis — recharge in two stations]
The table below gives mean annual water-balance data (in mm) for two Cameroonian stations, A (Sahelian north) and B (humid south).

\begin{center}
\begin{tabularx}{\linewidth}{|l|X|X|}
\hline
\rowcolor{popblueL} Element & Station A & Station B \\
\hline
Precipitation $P$ & 650 & 2400 \\
\hline
Evapotranspiration $E$ & 520 & 1500 \\
\hline
Runoff $R$ & 100 & 500 \\
\hline
\end{tabularx}
\end{center}

\begin{enumerate}
\item Compute the infiltration $I$ for each station using $P = E + R + I$.
\item Compute the ratio $I/P$ for each station.
\item Compare the potential for groundwater recharge at the two stations and explain the difference using climatic and geological factors.
\item Explain why a high runoff value is generally unfavourable to groundwater recharge.
\end{enumerate}
\end{exercice}

\begin{exercice}[Darcy's law — discharge through an aquifer]
A horizontal sandy aquifer has a hydraulic conductivity $K = 20\ \mathrm{m/day}$, a cross-sectional area $A = 500\ \mathrm{m^2}$ perpendicular to flow, and a hydraulic gradient $i = 0.004$.
\begin{enumerate}
\item State Darcy's law, defining each symbol and its unit.
\item Calculate the groundwater discharge $Q$ through the aquifer in $\mathrm{m^3/day}$.
\item If the hydraulic gradient doubles while $K$ and $A$ remain constant, what happens to $Q$? Justify.
\item Give one reason why the real discharge in a fractured basalt aquifer (e.g. on the slopes of Mount Cameroon) may differ from the value predicted by Darcy's law.
\end{enumerate}
\end{exercice}

\begin{exercice}[Porosity and permeability of sediments]
Three sediment samples were tested in the laboratory:

\begin{center}
\begin{tabularx}{\linewidth}{|l|X|X|}
\hline
\rowcolor{popblueL} Sample & Porosity (\%) & Permeability \\
\hline
Well-sorted coarse sand & 35 & very high \\
\hline
Clay & 45 & very low \\
\hline
Poorly sorted gravel-sand-clay mixture & 20 & moderate \\
\hline
\end{tabularx}
\end{center}

\begin{enumerate}
\item Explain why the clay has the highest porosity but the lowest permeability.
\item Explain why the poorly sorted mixture has a lower porosity than the well-sorted sand.
\item Which sample would make the best aquifer? Justify your choice.
\item Which sample would best serve as an aquiclude? Justify.
\end{enumerate}
\end{exercice}

\begin{exercice}[Reading a cross-section]
A cross-section shows, from top to bottom: a permeable sandstone layer outcropping on a hill slope, underlain by an impermeable shale layer that dips towards a valley; a line of springs occurs where the sandstone/shale contact meets the valley side.
\begin{enumerate}
\item Identify the aquifer and the aquiclude, justifying each choice.
\item Name the type of spring formed.
\item With the aid of a simple labelled sketch, explain the origin of the spring.
\item Predict what would happen to the spring discharge during a long dry season, and explain why.
\end{enumerate}
\end{exercice}

\courssub{I prepare for the GCE}

\begin{exercice}[Structured question — spring and aquifer (GCE Paper 2 style)]
The figure below (to be drawn by the candidate as part of the answer) represents a cross-section through a synclinal basin: a permeable limestone layer is sandwiched between two impermeable marl layers; the limestone outcrops on both flanks at higher elevation; a borehole sunk in the centre of the basin yields water that rises to the surface without pumping.

\begin{enumerate}
\item Copy and complete the cross-section, labelling: aquifer, aquicludes, recharge zones, water table, piezometric surface. \textbf{(6 marks)}
\item Name the type of well obtained in the centre of the basin and explain, with reference to pressure, why water rises without pumping. \textbf{(4 marks)}
\item State two conditions necessary for such a system to develop. \textbf{(2 marks)}
\item Explain two ways in which human activity in the recharge zone could affect the well. \textbf{(4 marks)}
\item Name one region of Cameroon where confined groundwater is exploited and state its importance to the local population. \textbf{(2 marks)}
\end{enumerate}
\textbf{(Total: 18 marks)}
\end{exercice}

\begin{exercice}[Case study — the salty well of Douala]
A community well near the coast in Douala supplied fresh water for many years. After a decade of increasingly heavy pumping, the water has become progressively salty.

\begin{enumerate}
\item Name and explain the phenomenon responsible for the salinisation of the well. \textbf{(4 marks)}
\item With the aid of a labelled diagram, describe the relationship between fresh groundwater and salt water in a coastal aquifer before and after heavy pumping. \textbf{(6 marks)}
\item Explain why the excessive pumping disturbs this equilibrium. \textbf{(3 marks)}
\item Propose three practical measures that could be taken to remedy or prevent this situation. \textbf{(3 marks)}
\item Briefly state one other threat to groundwater quality that affects wells in the Douala urban area. \textbf{(2 marks)}
\end{enumerate}
\textbf{(Total: 18 marks)}
\end{exercice}

\begin{exercice}[Essay — storage and movement of groundwater]
\textbf{``Discuss the factors controlling the storage and movement of groundwater, illustrating your answer with Cameroonian examples.''}

Your essay should:
\begin{enumerate}
\item define the key terms (porosity, permeability, aquifer, aquiclude, water table);
\item explain how rock type, texture, fracturing and geological structure control groundwater storage and flow;
\item discuss the role of climate and relief in recharge;
\item use at least three Cameroonian examples (e.g. volcanic aquifers of the Cameroon Volcanic Line, sedimentary basins of the north, crystalline basement of the south);
\item be organised in an introduction, a developed body and a conclusion.
\end{enumerate}
\textbf{(Total: 20 marks)}
\end{exercice}

\newpage

\coursec{Solutions to the exercises}

\coursec{Corrected Exercises — Hydrogeology}

\begin{corrige}[Exercise 1]
\textbf{Multiple choice questions.}
\begin{enumerate}
\item \textbf{Answer: (b) evaporation.} Evaporation is the change of state from liquid water to water vapour, driven mainly by solar energy. Condensation (a) is the reverse process (vapour to liquid), infiltration (c) is the entry of water into the soil, and precipitation (d) is water falling from the atmosphere.
\item \textbf{Answer: (b) low permeability.} Porosity measures the \emph{volume} of voids; permeability measures the ability of the rock to \emph{transmit} water. A rock such as clay can be highly porous yet almost impermeable because its pores are too small and poorly connected for water to move.
\item \textbf{Answer: (c) the water table.} The water table is the upper boundary of the zone of saturation, where all pores are filled with water. The piezometric surface (b) applies to \emph{confined} aquifers, and the capillary fringe (d) lies just above the water table.
\item \textbf{Answer: (c) well-sorted coarse sandstone.} It combines high porosity with large, well-connected pores, giving high permeability. Shale and clay (a, d) are aquicludes, and unfractured granite (b) has almost no primary porosity.
\end{enumerate}
\end{corrige}

\begin{corrige}[Exercise 2]
\textbf{True or false.}
\begin{enumerate}
\item \textbf{False.} A good aquifer must be both porous \emph{and} permeable. Clay is highly porous (up to about 50\%) but its pores are minute and poorly connected, so it cannot transmit water and acts as an aquiclude.
\item \textbf{False.} The water table fluctuates seasonally: it rises during the rainy season when recharge is high and falls during the dry season when discharge and evaporation dominate. In the southern part of Cameroon it is typically highest around September–October.
\item \textbf{True.} In an artesian well, the confined aquifer is under pressure because its recharge zone lies at a higher elevation; when the well head is below the piezometric surface, water rises freely above the top of the aquifer, sometimes to the surface, without pumping.
\item \textbf{False.} Groundwater moves very slowly, so pollutants are flushed out only over years to centuries; some contaminants also sorb onto aquifer materials or degrade very slowly. Remediation is therefore slow, costly and often incomplete — prevention is far better than cure.
\end{enumerate}
\end{corrige}

\begin{corrige}[Exercise 3]
\textbf{Definitions.}
\begin{enumerate}
\item \textbf{Infiltration:} the process by which water on the ground surface soaks downward into the soil and underlying rocks under gravity. Example: rainwater entering the sandy soils of the Adamawa plateau after a storm.
\item \textbf{Aquiclude:} a rock formation that may store water (it can be porous) but cannot transmit it in significant amounts because of very low permeability. Example: a clay or shale layer confining an aquifer.
\item \textbf{Spring:} a natural outflow of groundwater at the Earth's surface, occurring where the water table intersects the topography or where water is forced out along a contact, fault or fracture. Example: the springs on the flanks of Mount Cameroon.
\item \textbf{Hydraulic gradient:} the change in hydraulic head per unit horizontal distance along the direction of groundwater flow, $i = \Delta h / L$; it is dimensionless and drives flow according to Darcy's law.
\item \textbf{Zone of aeration (unsaturated zone):} the layer of soil and rock above the water table in which the pores contain both air and water; the water is held mainly by capillary forces and is moving downward.
\end{enumerate}
\end{corrige}

\begin{corrige}[Exercise 4]
\textbf{Direct calculation: water balance.}
\begin{enumerate}
\item From the water balance equation $P = E + R + I$:
\[ I = P - E - R = 1800 - 1200 - 350 = 250\ \mathrm{mm}. \]
The annual infiltration is $\boxed{I = 250\ \mathrm{mm}}$.
\item As a percentage of precipitation:
\[ \frac{I}{P}\times 100 = \frac{250}{1800}\times 100 \approx 13.9\ \%. \]
\item One factor that could increase $I$: a \textbf{more permeable soil or fractured bedrock} (e.g. fractured volcanic rocks) allowing water to enter the ground more easily. (Also acceptable: gentler slopes, denser vegetation cover reducing runoff, or less intense, longer-duration rainfall.)
\end{enumerate}
\end{corrige}

\begin{corrige}[Exercise 5]
\textbf{Data analysis: recharge in two stations.}
\begin{enumerate}
\item Using $I = P - E - R$:
\begin{itemize}
\item Station A: $I_A = 650 - 520 - 100 = 30\ \mathrm{mm}$.
\item Station B: $I_B = 2400 - 1500 - 500 = 400\ \mathrm{mm}$.
\end{itemize}
\item Ratios:
\[ \left(\frac{I}{P}\right)_A = \frac{30}{650} \approx 0.046 \approx 4.6\ \%; \qquad \left(\frac{I}{P}\right)_B = \frac{400}{2400} \approx 0.167 \approx 16.7\ \%. \]
\item \textbf{Station B has a far greater recharge potential}, both in absolute terms ($400$ mm against $30$ mm) and as a fraction of rainfall. \emph{Climatic factors:} the humid south receives abundant, frequent rainfall, so the soil is often moist and water is available to percolate; in the Sahelian north, rainfall is scarce, brief and intense, and high temperatures drive strong evapotranspiration, leaving little surplus water. \emph{Geological factors:} recharge is enhanced where permeable formations (fractured basement, sands) outcrop; where the surface is clay-rich or indurated (e.g. lateritic crusts), infiltration is reduced even in humid areas.
\item High runoff means that a large proportion of the rainfall flows rapidly over the surface into rivers instead of soaking into the ground. The water is therefore \emph{lost to the subsurface} before it can infiltrate: the residence time on the surface is too short. Hence high runoff (steep slopes, intense storms, impermeable or crusted surfaces) reduces the water available for recharge and lowers $I$.
\end{enumerate}
\end{corrige}

\begin{corrige}[Exercise 6]
\textbf{Darcy's law.}
\begin{enumerate}
\item \textbf{Darcy's law:} $Q = K \cdot i \cdot A$, where:
\begin{itemize}
\item $Q$ = discharge (volume of water flowing per unit time), in $\mathrm{m^3\,day^{-1}}$;
\item $K$ = hydraulic conductivity (ability of the material to transmit water), in $\mathrm{m\,day^{-1}}$;
\item $i$ = hydraulic gradient ($\Delta h / L$), dimensionless;
\item $A$ = cross-sectional area perpendicular to flow, in $\mathrm{m^2}$.
\end{itemize}
\item Substituting the values:
\[ Q = K \cdot i \cdot A = 20 \times 0.004 \times 500 = 40\ \mathrm{m^3\,day^{-1}}. \]
\[ \boxed{Q = 40\ \mathrm{m^3\,day^{-1}}} \]
\item Since $Q$ is \textbf{directly proportional} to $i$ (with $K$ and $A$ constant), doubling the gradient doubles the discharge: $Q' = 2Q = 80\ \mathrm{m^3\,day^{-1}}$.
\item Darcy's law assumes \emph{laminar flow through a homogeneous porous medium}. In a fractured basalt, flow is concentrated in fractures and fissures rather than distributed through pores; it may be turbulent and highly anisotropic (much faster along fractures, almost nil in intact blocks). The effective $K$ is therefore very difficult to define, and the real discharge can be much larger along fracture zones than the value predicted from an average $K$.
\end{enumerate}
\end{corrige}

\begin{corrige}[Exercise 7]
\textbf{Porosity and permeability of sediments.}
\begin{enumerate}
\item Clay consists of extremely fine platy particles, so the total void space between them is large — hence the \textbf{high porosity} (45\%). However, the individual pores are microscopic, poorly connected, and water is held strongly by capillary and electrostatic forces on the large surface area of the clay particles. Water therefore cannot move through it: \textbf{permeability is very low}.
\item In a poorly sorted mixture, the fine particles (sand, silt, clay) fill the voids between the larger gravel clasts. The total void volume is thus reduced, giving a lower porosity (20\%) than a well-sorted sand (35\%), in which grains of similar size leave large open spaces between them.
\item \textbf{The well-sorted coarse sand} is the best aquifer: it combines high porosity (good storage) with very high permeability (easy transmission and extraction of water). The mixture is only moderate, and the clay cannot yield water.
\item \textbf{The clay} is the best aquiclude: its very low permeability prevents the movement of groundwater, so it can confine an aquifer or act as a barrier, despite its high porosity.
\end{enumerate}
\end{corrige}

\begin{corrige}[Exercise 8]
\textbf{Reading a cross-section.}
\begin{enumerate}
\item \textbf{Aquifer: the sandstone}, because it is permeable and can store and transmit groundwater. \textbf{Aquiclude: the shale}, because it is impermeable and blocks the downward percolation of water.
\item This is a \textbf{contact spring} (also called a scarp-foot or stratigraphic spring), formed where a permeable layer overlies an impermeable layer and the contact crops out on a slope.
\item \textbf{Origin:} rainwater infiltrates the sandstone where it outcrops on the hill slope and percolates downward until it reaches the impermeable shale. Unable to go deeper, the water flows \emph{laterally along the contact}, down the dip of the shale, towards the valley. Where the valley side cuts across the contact, the water table intersects the ground surface and the groundwater emerges as a line of springs.
\begin{popfigure}
\begin{tikzpicture}[scale=0.95]
% Topography and layers
\draw[thick,popdark] (0,3) -- (2,3.2) -- (5,2) -- (7,2) -- (7,0) -- (0,0) -- cycle;
% Sandstone layer (top)
\fill[popblueL] (0,3) -- (2,3.2) -- (5,2) -- (7,2) -- (7,1.2) -- (5,1.2) -- (2,2.4) -- (0,2.2) -- cycle;
% Shale layer (bottom)
\fill[gray!40] (0,2.2) -- (2,2.4) -- (5,1.2) -- (7,1.2) -- (7,0) -- (0,0) -- cycle;
% Contacts
\draw[thick] (0,2.2) -- (2,2.4) -- (5,1.2) -- (7,1.2);
% Rain
\draw[->,thick,popblue] (1,3.5) -- (1,3.0);
\draw[->,thick,popblue] (3,3.5) -- (3,2.8);
\node[popblue] at (1,3.65) {\small rain};
\node[popblue] at (3,3.65) {\small rain};
% Labels
\node[popdark] at (3.5,2.5) {\small Sandstone (aquifer)};
\node[popdark] at (3.5,0.6) {\small Shale (aquiclude)};
% Water table
\draw[dashed,popblue] (0.5,2.7) -- (2,2.55) -- (4.5,1.55) -- (6,1.3);
\node[popblue] at (1.5,2.3) {\tiny water table};
% Spring at outcrop of contact on valley side (x=7)
\draw[fill=popblue] (7,1.2) circle (0.08);
\node[popblue,anchor=west] at (7.1,1.2) {\small spring};
% Flow arrows
\draw[->,thick,popblue] (2.5,2.5) -- (4,1.7);
\draw[->,thick,popblue] (4,1.7) -- (6.5,1.35);
\end{tikzpicture}
\legende{Contact spring at the sandstone/shale boundary on the valley side.}
\end{popfigure}
\item During a long dry season the spring discharge \textbf{decreases and may cease entirely}. There is no recharge by rainfall, so the water table within the sandstone falls as the stored water drains out; when the water table drops below the level of the contact at the valley side, the spring dries up. Contact springs fed by thin aquifers with small storage are typically seasonal.
\end{enumerate}
\end{corrige}

\begin{corrige}[Exercise 9]
\textbf{Structured question: spring and aquifer (GCE Paper 2 style).}
\textbf{Indicative mark scheme and model answer (18 marks).}

\begin{enumerate}
\item \textbf{Labelled cross-section (6 marks):} 1 mark for each correctly placed and labelled feature: aquifer (limestone), upper aquiclude, lower aquiclude (marls), recharge zones (both outcrops), water table (in the unconfined flanks), piezometric surface (above the aquifer in the basin centre).
\begin{popfigure}
\begin{tikzpicture}[scale=0.9]
% Topography
\draw[thick,popdark] (-4,3.5) -- (-1.5,1.5) -- (0,1) -- (1.5,1.5) -- (4,3.5);
% Upper marl (aquiclude)
\fill[gray!30] (-4,3.5) -- (-1.5,1.5) -- (0,1) -- (1.5,1.5) -- (4,3.5) -- (4,3) -- (1.5,1.2) -- (0,0.7) -- (-1.5,1.2) -- (-4,3) -- cycle;
% Limestone aquifer
\fill[popblueL] (-4,3) -- (-1.5,1.2) -- (0,0.7) -- (1.5,1.2) -- (4,3) -- (4,2.2) -- (1.5,0.4) -- (0,0) -- (-1.5,0.4) -- (-4,2.2) -- cycle;
% Lower marl (aquiclude)
\fill[gray!30] (-4,2.2) -- (-1.5,0.4) -- (0,0) -- (1.5,0.4) -- (4,2.2) -- (4,1.5) -- (1.5,-0.2) -- (0,-0.5) -- (-1.5,-0.2) -- (-4,1.5) -- cycle;
% Labels
\node[popdark] at (0,0.35) {\small limestone aquifer};
\node[popdark] at (-2.6,3.2) {\small marl (aquiclude)};
\node[popdark] at (2.6,-0.15) {\small marl (aquiclude)};
% Recharge arrows
\draw[->,thick,popblue] (-3.6,4) -- (-3.6,3.3);
\draw[->,thick,popblue] (3.6,4) -- (3.6,3.3);
\node[popblue] at (-3.6,4.2) {\small recharge};
\node[popblue] at (3.6,4.2) {\small recharge};
% Piezometric surface
\draw[dashed,popblue] (-4,3.0) -- (4,3.0);
\node[popblue] at (0,3.2) {\small piezometric surface};
% Artesian well
\draw[thick,popdark] (0,2.8) -- (0,-0.2);
\draw[->,thick,popblue] (0,-0.2) -- (0,2.8);
\node[popdark,anchor=west] at (0.2,1.3) {\small artesian well};
\node[popdark,anchor=east] at (-0.2,2.9) {\small well head};
\end{tikzpicture}
\legende{Confined synclinal aquifer with artesian well (schematic).}
\end{popfigure}
\item \textbf{Artesian well (4 marks).} The well is an \textbf{artesian (flowing) well} \textbf{(1)}. The confined aquifer is recharged at high elevation on the flanks; the water in the basin centre is therefore under a hydrostatic pressure corresponding to the height of the column of water up to the piezometric surface \textbf{(2)}. When the borehole pierces the upper aquiclude, this pressure is released and water rises \textbf{(1)}; because the well head lies below the piezometric surface, the water reaches the surface without pumping \textbf{(1)}.
\item \textbf{Conditions (2 marks, any two, 1 each):}
\begin{itemize}
\item a permeable aquifer sandwiched between impermeable layers (confined aquifer);
\item the aquifer must outcrop at a higher elevation (recharge zone above the well head);
\item a suitable structure (syncline or dipping strata) and sufficient rainfall for recharge.
\end{itemize}
\item \textbf{Human activity in the recharge zone (4 marks, 2 per point):}
\begin{itemize}
\item \emph{Deforestation / cultivation / construction} reduces infiltration and increases runoff, so recharge decreases and the piezometric level (and well discharge) falls.
\item \emph{Pollution} (fertilisers, pesticides, waste dumps, fuel storage) infiltrates with recharge water and contaminates the whole confined system, including the well, and such pollution is very difficult to remove.
\end{itemize}
\item \textbf{Cameroonian example (2 marks):} the confined aquifers of the \textbf{Chad Basin / Logone plain in the Far North} (also acceptable: the Douala sedimentary basin). Importance \textbf{(1)}: in this arid region with very scarce surface water, deep confined groundwater provides a reliable year-round supply for drinking, livestock and small-scale irrigation \textbf{(1)}.
\end{enumerate}

\textbf{Examiner expectations:} a clear, correctly orientated synclinal diagram; precise use of the terms \emph{confined aquifer}, \emph{piezometric surface} and \emph{artesian}; explicit link between recharge elevation and pressure; human impacts linked logically to recharge and quality.
\end{corrige}

\begin{corrige}[Exercise 10]
\textbf{Case study: the salty well of Douala.}
\textbf{Indicative mark scheme and model answer (18 marks).}

\begin{enumerate}
\item \textbf{Phenomenon (4 marks):} \textbf{salt-water intrusion} (marine intrusion into the coastal aquifer) \textbf{(1)}. Coastal aquifers naturally contain fresh groundwater floating on denser sea water; the interface between them lies inland and at depth \textbf{(1)}. Excessive pumping lowers the water table/piezometric level, so the pressure of fresh water pushing seaward decreases \textbf{(1)}; the denser salt water then migrates inland and upward (upconing beneath the well), and the well begins to pump brackish then salty water \textbf{(1)}.
\item \textbf{Diagram and description (6 marks):} diagram \textbf{(3)}: correct layering, interface sloping inland, change after pumping; description \textbf{(3)}.
\begin{popfigure}
\begin{tikzpicture}[scale=0.85]
% Before pumping
\fill[popblueL] (0,0) rectangle (5,1.6);
\fill[popgold!60] (0,0) -- (5,0) -- (5,0.7) -- (1.5,0) -- cycle;
\draw[thick] (0,1.6) -- (5,1.6);
\draw[thick] (1.5,0) -- (5,0.7);
\draw[thick,popblue] (5,0) -- (5,1.9);
\node[popdark] at (2.5,1.2) {\small fresh water};
\node[popgold!70!black] at (3.6,0.25) {\small salt water};
\node[popblue] at (5.35,1.0) {\small sea};
\node at (2.5,-0.5) {\small Before pumping};
% After pumping
\begin{scope}[xshift=7cm]
\fill[popblueL] (0,0) rectangle (5,1.6);
\fill[popgold!60] (0,0) -- (5,0) -- (5,1.1) -- (0.8,0) -- cycle;
\draw[thick] (0,1.6) -- (5,1.6);
\draw[thick] (0.8,0) -- (5,1.1);
\draw[thick,popblue] (5,0) -- (5,1.9);
\draw[thick,popdark] (2.5,1.6) -- (2.5,0.1);
\draw[thick,popgold!70!black] (2.5,0.1) -- (2.2,0.55) -- (2.8,0.55) -- cycle;
\node[popdark,anchor=west] at (2.6,1.0) {\small well};
\node[popgold!70!black] at (1.2,0.75) {\small upconing};
\node at (2.5,-0.5) {\small After heavy pumping};
\end{scope}
\end{tikzpicture}
\legende{Fresh/salt water interface in a coastal aquifer before and after over-pumping.}
\end{popfigure}
Before pumping, fresh groundwater (lower density) floats above salt water, with a sloping interface: the fresh-water lens thickens inland and the interface toe lies seaward of the well. After heavy pumping, the fresh-water head falls, the interface moves inland, and a cone of salt water rises beneath the pumping well (upconing), contaminating it.
\item \textbf{Why pumping disturbs the equilibrium (3 marks):} the equilibrium is maintained by the seaward flow and hydrostatic pressure of fresh water balanced against the denser salt water \textbf{(1)}. Pumping removes fresh water faster than recharge replaces it, lowering the water table \textbf{(1)}; the reduced fresh-water pressure can no longer hold back the sea water, which advances inland and rises toward the well screen \textbf{(1)}.
\item \textbf{Measures (3 marks, any three, 1 each):}
\begin{itemize}
\item reduce or ration pumping rates and spread abstraction over several wells;
\item artificially recharge the aquifer (e.g. infiltration basins with surplus surface water) to rebuild the fresh-water head;
\item relocate wells further inland, away from the interface;
\item monitor water levels and salinity to regulate exploitation;
\item develop alternative supplies (treated surface water) to relieve demand on the aquifer.
\end{itemize}
\item \textbf{Other threat (2 marks):} contamination by \textbf{untreated domestic sewage and pit latrines} (nitrate and bacterial pollution) in the densely populated, unsewered parts of Douala, where the shallow water table is easily reached by pollutants. (Also acceptable: industrial effluents, hydrocarbon leaks from fuel depots and the port.)
\end{enumerate}

\textbf{Examiner expectations:} correct naming of salt-water intrusion; a density-based explanation (fresh water floats on salt water); the causal chain pumping $\rightarrow$ falling head $\rightarrow$ landward/upward movement of the interface; realistic, feasible remedies.
\end{corrige}

\begin{corrige}[Exercise 11]
\textbf{Essay: storage and movement of groundwater.}
\textbf{Model essay plan with indicative content (20 marks).}

\textbf{Introduction (2 marks).} Groundwater is the subsurface component of the hydrologic cycle and the main source of well and spring water in Cameroon. Its \emph{storage} and \emph{movement} are controlled by the physical properties of rocks, the geological structure, and the climate and relief that govern recharge.

\textbf{1. Key definitions (4 marks, 1 each).}
\begin{itemize}
\item \textbf{Porosity:} percentage of void space in a rock; controls storage capacity.
\item \textbf{Permeability:} ability of a rock to transmit water through interconnected pores or fractures; controls movement.
\item \textbf{Aquifer:} a permeable, saturated formation that stores and yields useful quantities of water (e.g. sandstone, fractured basalt).
\item \textbf{Aquiclude:} an impermeable formation (clay, shale, unfractured granite) that blocks flow and confines aquifers.
\item \textbf{Water table:} the upper surface of the zone of saturation, fluctuating with recharge.
\end{itemize}

\textbf{2. Rock type, texture, fracturing and structure (6 marks).}
\begin{itemize}
\item \emph{Texture:} well-sorted coarse sediments (sands, gravels) have high porosity and permeability; poorly sorted mixtures and clays store but do not transmit water. Storage without permeability (clay) gives an aquiclude, not an aquifer.
\item \emph{Fracturing:} crystalline rocks (granite, gneiss) and basalts have virtually no primary porosity; they become aquifers only through joints, faults and weathered zones (secondary porosity). Flow is then channelled along fractures — anisotropic and hard to predict by Darcy's law.
\item \emph{Structure:} alternating permeable/impermeable beds create confined aquifers; synclines and dipping strata produce artesian systems; faults may act as conduits (springs) or barriers.
\end{itemize}

\textbf{3. Climate and relief in recharge (4 marks).} Recharge $I = P - E - R$: abundant, moderate-intensity rainfall and low evapotranspiration (humid south) favour infiltration; scarce, intense Sahelian rains with high evaporation (north) give minimal recharge. Gentle slopes and vegetated surfaces increase infiltration; steep relief and crusted surfaces increase runoff and reduce recharge.

\textbf{4. Cameroonian examples (3 marks, 1 each, any three).}
\begin{itemize}
\item \emph{Cameroon Volcanic Line / Mount Cameroon:} fractured basalt and scoria aquifers feed abundant flank springs; flow is fracture-controlled.
\item \emph{Sedimentary basins of the north (Chad Basin, Benue trough):} layered sandstones and sands confined by clays give important artesian groundwater in an arid climate.
\item \emph{Crystalline basement of the south and centre:} groundwater occurs only in the weathered mantle and fracture zones; wells have modest, variable yields.
\item \emph{Douala coastal basin:} sedimentary aquifers threatened by salt-water intrusion under heavy pumping.
\end{itemize}

\textbf{Conclusion (1 mark).} Groundwater storage depends on porosity, movement on permeability; both are governed by lithology, texture, fracturing and structure, while climate and relief determine how much water enters the system. Sustainable exploitation in Cameroon requires matching abstraction to the recharge and storage characteristics of each geological setting.

\textbf{Examiner expectations:} accurate definitions; clear distinction porosity/permeability; logical cause–effect reasoning; at least three correctly used Cameroonian examples; coherent essay structure with introduction, development and conclusion; correct geological vocabulary throughout.
\end{corrige}

\newpage

\coursec{Summary sheet}

\begin{popfigure}
\begin{tikzpicture}[mindmap, grow cyclic, every node/.style={concept}, text=white,
root concept/.append style={concept color=popdark, font=\large\bfseries},
level 1/.append style={level distance=3.6cm, sibling angle=60, font=\small\bfseries},
level 2/.append style={level distance=2.2cm, sibling angle=45, font=\scriptsize}]
\node[root concept]{Hydrogeology}
  child[concept color=popblue]{ node {Hydrologic cycle}
    child[concept color=popblue!70]{ node {Evaporation, condensation}}
    child[concept color=popblue!70]{ node {Infiltration vs runoff}}
  }
  child[concept color=popgreen]{ node {Sources of groundwater}
    child[concept color=popgreen!70]{ node {Meteoric (main)}}
    child[concept color=popgreen!70]{ node {Juvenile, connate, magmatic}}
  }
  child[concept color=poporange]{ node {Circulation \& accumulation}
    child[concept color=poporange!70]{ node {Zone of aeration}}
    child[concept color=poporange!70]{ node {Zone of saturation; water table}}
  }
  child[concept color=poppurple]{ node {Springs, wells, rivers}
    child[concept color=poppurple!70]{ node {Spring types; artesian well}}
    child[concept color=poppurple!70]{ node {Effluent / influent rivers}}
  }
  child[concept color=poppink]{ node {Rock properties}
    child[concept color=poppink!70]{ node {Porosity $n=V_v/V_t$}}
    child[concept color=poppink!70]{ node {Permeability; Darcy $Q=KiA$}}
  }
  child[concept color=popgold]{ node {Aquifers \& quality}
    child[concept color=popgold!70]{ node {Aquifer, aquiclude, aquitard}}
    child[concept color=popgold!70]{ node {Contamination \& protection}}
  };
\end{tikzpicture}
\legende{Mind map of Chapter GEO07 -- Hydrogeology.}
\end{popfigure}

\begin{bilan}
\textbf{1. Key definitions}
\begin{itemize}
\item \cle{Hydrogeology}: study of the occurrence, circulation, distribution and quality of underground water.
\item \cle{Hydrologic cycle}: continuous movement of water between oceans, atmosphere and land (evaporation $\rightarrow$ condensation $\rightarrow$ precipitation $\rightarrow$ infiltration/runoff $\rightarrow$ back to oceans).
\item \cle{Infiltration}: entry of rainwater into the soil; \cle{percolation}: its downward movement under gravity towards the water table.
\item \cle{Water table}: upper surface of the zone of saturation, where water pressure equals atmospheric pressure.
\item \cle{Zone of aeration} (vadose): pores contain air + water; \cle{zone of saturation} (phreatic): all pores filled with water.
\item \cle{Spring}: natural outflow of groundwater at the surface where the water table intersects the topography.
\item \cle{Aquifer}: permeable rock that stores \emph{and} transmits exploitable water; \cle{aquiclude}: impermeable rock that stores but does not transmit (clay); \cle{aquitard}: transmits slowly; \cle{aquifuge}: neither stores nor transmits (massive granite).
\item \cle{Porosity} $n$: percentage of void space; \cle{permeability}: ability to transmit water through \emph{connected} pores.
\item \cle{Artesian well}: well in a confined aquifer where water rises (sometimes to the surface) under pressure.
\end{itemize}

\textbf{2. Sources of groundwater}
\begin{itemize}
\item \cle{Meteoric water}: from precipitation -- by far the main source.
\item \cle{Connate water}: trapped in sediments during deposition (fossil water).
\item \cle{Juvenile/magmatic water}: released from magma, reaching the surface for the first time (volcanic regions, e.g.\ Cameroon Volcanic Line).
\end{itemize}

\textbf{3. Laws and formulas to know}
\begin{itemize}
\item Porosity: $n = \dfrac{V_v}{V_t}\times 100\ \%$ ($V_v$ = volume of voids, $V_t$ = total volume).
\item \cle{Darcy's law}: $Q = K\,i\,A$ with hydraulic gradient $i = \dfrac{\Delta h}{L}$; $Q$ = discharge, $K$ = coefficient of permeability, $A$ = cross-sectional area.
\item Specific yield $+$ specific retention $=$ porosity.
\item Water balance: $P = E + R + I$ (precipitation = evapotranspiration + runoff + infiltration).
\end{itemize}

\textbf{4. Classification of aquifers and springs}
\begin{itemize}
\item \cle{Unconfined aquifer}: open to the atmosphere, bounded above by the water table.
\item \cle{Confined aquifer}: sandwiched between two aquicludes, water under pressure $\rightarrow$ artesian conditions.
\item \cle{Perched aquifer}: local lens of saturation above the main water table, on an impermeable layer.
\item Spring types: depression, contact, fracture/fault, karst, artesian springs.
\item \cle{Effluent river}: fed by groundwater (water table above river bed); \cle{influent river}: loses water to the aquifer.
\end{itemize}

\textbf{5. Rock properties controlling flow}
\begin{itemize}
\item High porosity + high permeability: gravel, sand, sandstone, fractured basalt/limestone $\rightarrow$ good aquifers.
\item High porosity but low permeability: clay (pores too small, water held by molecular attraction).
\item Low porosity: massive crystalline rocks -- aquifers only when fractured (e.g.\ basement rocks of the Adamawa Plateau).
\item Karst: dissolution of limestone by \ce{CO2}-rich water creates conduits, caves and large springs.
\end{itemize}

\textbf{6. Water quality and contamination}
\begin{itemize}
\item Quality judged by: physical characters (colour, turbidity, taste, temperature), chemical characters (pH, hardness from \ce{Ca^{2+}}/\ce{Mg^{2+}}, dissolved salts), biological characters (pathogens, e.g.\ \emph{E. coli}).
\item Pollution sources: latrines and septic tanks, fertilisers/pesticides, industrial waste, oil spills, saline intrusion in coastal aquifers (e.g.\ Douala, Kribi).
\item Protection: safe distances between wells and latrines, casing of wells, protection zones around catchments, regular water analysis.
\end{itemize}

\textbf{7. Common mistakes to avoid}
\begin{itemize}
\item Confusing \textbf{porosity} (storage) with \textbf{permeability} (transmission) -- clay is the classic trap.
\item Saying the water table is horizontal and fixed: it mimics the relief and fluctuates seasonally.
\item Believing an artesian well always flows at the surface: water only needs to rise \emph{above the top of the confined aquifer}.
\item Forgetting that groundwater flows from high to low \emph{hydraulic head}, not necessarily downwards.
\item Confusing aquiclude (stores but does not transmit) with aquifuge (neither).
\end{itemize}

\textbf{8. GCE tips}
\begin{itemize}
\item In structured questions, always \textbf{define the term first}, then describe, then give a named example (Cameroon examples score well: Mount Cameroon springs, Lake Nyos gas-rich groundwater, Douala coastal aquifer).
\item Practise drawing and labelling: the hydrologic cycle, a confined/unconfined aquifer cross-section with an artesian well, and the zones of aeration/saturation.
\item Numerical questions on Darcy's law: convert all units to metres and seconds before substituting.
\item For essays on contamination, structure your answer: sources $\rightarrow$ pathways $\rightarrow$ effects $\rightarrow$ prevention.
\end{itemize}
\end{bilan}

\findecours

\end{document}
